\providecommand{\recsinc}{\operatorname{sinc}}
\providecommand{\recgE}{\gamma_{\mathrm{E}}}
\providecommand{\recGen}{\operatorname{Gen}}
\providecommand{\recav}{\operatorname{av}}
\providecommand{\recA}{\mathcal{A}}
\providecommand{\recartanh}{\operatorname{artanh}}
\providecommand{\zerJ}{\mathcal{J}}
\providecommand{\zerK}{\mathcal{K}}
\providecommand{\zersinc}{\operatorname{sinc}}
\providecommand{\zergE}{\gamma_{\mathrm{E}}}
\providecommand{\zerLL}{\mathcal{L}}
\providecommand{\zerconv}{\operatorname{conv}}
\providecommand{\zerdiag}{\operatorname{diag}}
\providecommand{\zerSym}{\mathfrak{D}}
\providecommand{\zerLog}{\operatorname{Log}}
\section{Characters and arithmetic parameters from zeros}\label{sec:chz}

\subsection{The modulus and the twist from the signed zero spectrum}\label{rec:sec:spectrum}
\sources{dirichlet-grh-research; exponent_seven_and_grh_transfers; grh-torsion-and-zero-data}

Here the data are on the zero side and the unknowns are arithmetic. The function $F$ below is the
signed counterpart of the function $K_q$ of \cref{rec:thm:primepower}: its coefficients are
differences of multiplicities of an untwisted and a twisted family, and its derivative jumps
again sit at logarithms of primes.

The Mellin residues of the smoothed prime remainder in the source proof record only the
difference between two sums of zero multiplicities; the relevant declaration of the source is
\code{New.isPolar_primeRemainder} (see \cref{src:sec:mechanism}). Contributions at a common zero can cancel exactly,
and the names of the characters have disappeared. Nevertheless the signed data determine both
arithmetic parameters, and for the parameters of the lattice argument a finite part of the
data does so with a certificate.

\paragraph{Setting.}
Two classes of parameters occur.
\begin{itemize}
\item The \emph{general class}: $q>0$ with $8\mid q$, and a prime $d\equiv7\pmod8$ coprime
to $q$.
\item The \emph{lattice class}: $18480\mid q$ and $d$ a prime with $d\equiv q/2-1\pmod q$.
Then $(d,q)=1$, $8\mid q$, $d\equiv7\pmod8$ and $q\mid2(d+1)$, so that $q\le2(d+1)$; this is
a subclass of the general class, and it is the one used in the source.
\end{itemize}
Write $h=\varphi(q)$ and $\eta(n)=(n/d)$. For each $\chi\bmod q$ let $\chi^*$ be its inducing
primitive character. For a real ordinate $\gamma$ put $\rho_\gamma=\tfrac12+i\gamma$ and
\begin{equation}\label{rec:eq:mgamma}
 m(\gamma)=\frac1h\sum_{\chi\bmod q}\Bigl(\ord_{\rho_\gamma}L(s,\chi^*)-
   \ord_{\rho_\gamma}L(s,\chi^*\eta)\Bigr),\qquad
 b_\gamma=\frac{m(\gamma)}{1/4+\gamma^2},\qquad
 F(t)=\sum_\gamma b_\gamma e^{i\gamma t}.
\end{equation}
Removing imprimitive Euler factors preserves orders in the open critical strip, so $b_\gamma$
is precisely the residue in the source, and by \cref{src:thm:dgrh} every nontrivial zero in
these sums is of the form $\rho_\gamma$. The input to the inverse problem is the signed
coefficient at each frequency, with coincident frequencies combined; both signs of the
ordinate are used.

\begin{proposition}[The spectral function; \stG]\label{rec:prop:ode}
In the general class the series $F$ is absolutely and uniformly convergent, real and even.
For a prime power $n=p^k$ let $Q_p=p^{v_p(q)}$, $q^{(p)}=q/Q_p$ and
\[
 a(p^k)=\frac{1-\eta(p^k)}{\varphi(Q_p)}\,\mathbf 1_{p^k\equiv1\bmod q^{(p)}} .
\]
Then, as distributions on $\R$, with $\delta_u$ the unit point mass at $u$,
\begin{equation}\label{rec:eq:ode}
 F''-\tfrac14F=(\log d)\,\delta_0-\frac2h\cosh(t/2)\,\dd t
 +\sum_{p^k\ge2}\frac{\log p}{p^{k/2}}\,a(p^k)\bigl(\delta_{\log p^k}+\delta_{-\log p^k}\bigr).
\end{equation}
Equivalently, for every real test function $g\in C_c^\infty(\R)$, with
$\widehat g(\gamma)=\int g(t)e^{i\gamma t}\,\dd t$,
\begin{equation}\label{rec:eq:weil}
 -\sum_\gamma m(\gamma)\widehat g(\gamma)
 =(\log d)\,g(0)-\frac2h\int g(t)\cosh(t/2)\,\dd t
 +\sum_{n=p^k}\frac{\log p}{\sqrt n}\,a(n)\bigl[g(\log n)+g(-\log n)\bigr].
\end{equation}
\end{proposition}
\begin{proof}
The classical zero count $N(T)=O_{q,d}(T\log(2+T))$ ensures $\sum_\gamma\abs{b_\gamma}<\infty$
by partial summation. Both families are closed under complex conjugation, whence
$b_{-\gamma}=b_\gamma\in\R$. Since $F''-F/4=-\sum_\gamma m(\gamma)e^{i\gamma t}$, the two
displays are equivalent, and it suffices to check \eqref{rec:eq:weil}. The Guinand--Weil
explicit formula for a primitive character of conductor $f$ contributes $\log(f/\pi)\,g(0)$,
the gamma term for its parity, and a negative sum over prime powers with coefficients
$\Lambda(n)\chi^*(n)/\sqrt n$ at one sign of $\log n$ and the conjugate coefficients at the
other. The primitive principal character also contributes
$\widehat g(i/2)+\widehat g(-i/2)$ from the poles at $0,1$ of the completed zeta function.
This is the normalization of \cite[\S2.1]{ChirrePereiraDeLaat2021Primes}; see also
\cite{Booker2015Lfunctions}. Equivalently the formula follows by shifting the contour of the
logarithmic derivative of the completed primitive $L$-function.

As $q>2$, half the characters have each parity, and the odd character $\eta$ interchanges
the parities. All gamma terms therefore cancel between the two families. The two conductors
for $\chi^*$ and $\chi^*\eta$ are $f$ and $fd$, since $(f,d)=1$, which gives $\log d$ at the
origin after the signed subtraction. Only the untwisted family contains the primitive
principal character, which gives the density after division by $h$. At $p^k$ the average of
the untwisted primitive character values is
$\mathbf 1_{p^k\equiv1\bmod q^{(p)}}/\varphi(Q_p)$: a nonprincipal $p$-component makes the
primitive value zero, and orthogonality over the characters modulo $q^{(p)}$ evaluates the
remaining sum. Multiplication by $1-\eta(p^k)$ gives the displayed $a$, also at $p=d$.
\end{proof}

For a point $u$ write $J(u)=F'(u+)-F'(u-)$, with $J(u)=0$ where $F$ is differentiable.

\begin{theorem}[Jumps and uniqueness; \stG]\label{rec:thm:jumps}
In the general class, $F$ is analytic away from the atoms of \eqref{rec:eq:ode}, and
\begin{equation}\label{rec:eq:mass}
 J(0)=\log d,\qquad
 J(\log p)=\frac{\log p}{\sqrt p}\bigl(1-\eta(p)\bigr)\mathbf 1_{p\equiv1\bmod q}
 \quad\text{for every prime }p\nmid qd .
\end{equation}
Thus $F'$ has a jump at $\log p$ precisely if $p\equiv1\pmod q$ and $\eta(p)=-1$, in which
case the jump is $2\log p/\sqrt p$. Let $S$ be the set of primes $p$ at which $F'$ has a
positive jump at $\log p$. For all sufficiently large $X$,
\[
 d=\exp J(0),\qquad q=\gcd\set{p-1:\ p\in S,\ p>X}.
\]
Consequently the data $\set{(\gamma,b_\gamma)}$ determine $(q,d)$, even after they have been
changed at finitely many frequencies; more generally the same reconstruction works after
adding any continuously differentiable function to $F$.
\end{theorem}
\begin{proof}
The right side of \eqref{rec:eq:ode} is smooth apart from a locally finite set of point
masses. Integrating across one such point shows that the jump of $F'$ equals its mass. Away
from the atoms $F$ satisfies an ordinary differential equation with analytic coefficients,
hence is analytic there. This gives \eqref{rec:eq:mass}; evenness then gives
$F(t)=F(0)+\tfrac12\log d\,\abs t+O_{q,d}(t^2)$ near zero.

Every entry $p-1$ with $p\in S$ is divisible by $q$ after the finitely many primes dividing
$qd$ have been discarded. Suppose a prime $\ell$ divided the gcd after division by $q$. We
produce infinitely many primes in $S$ with $\ell\nmid(p-1)/q$. If $\ell\mid q$, require
$p\equiv1+q\pmod{q\ell}$ and independently choose a nonzero quadratic nonresidue modulo $d$.
If $\ell\nmid qd$, then $\ell$ is odd; choose a unit other than $1$ modulo $\ell$, the
residue $1$ modulo $q$, and a nonresidue modulo $d$. If $\ell=d$, the conditions
$p\equiv1\pmod q$ and $\eta(p)=-1$ already imply $p\not\equiv1\pmod\ell$. Each construction
gives a reduced class by the Chinese remainder theorem and infinitely many primes by
Dirichlet's theorem. Thus no $\ell$ divides the quotient. All these jumps survive a $C^1$
perturbation, including any finite trigonometric polynomial.
\end{proof}

\paragraph{The spectral function on actual prime sums.}
Let $R$ be the remainder of the source proof,
\[
 R(x)=\sum_{\substack{n\le x\\n\equiv1\bmod q}}\Lambda(n)\bigl(1-\eta(n)\bigr)
   \Bigl(\frac1n-\frac1x\Bigr)-\frac{\log x}h-c,
\]
where $c$ gives $R(x)\to0$ by the prime number theorem for fixed progressions; in the
source this is the declaration \code{New.primeRemainder}.

\begin{proposition}[\stG]\label{rec:prop:remainder}
In the general class,
\begin{equation}\label{rec:eq:remainder}
 \sqrt x\,R(x)=F(\log x)+O_{q,d}\Bigl(\frac{1+\log x}{\sqrt x}\Bigr)\qquad(x\ge1).
\end{equation}
In the lattice class the comparison is exact for every $x\ge1$:
\[
 F(0)=-2c-\log d,\qquad
 \sqrt x\,R(x)=F(\log x)+\frac{c+\log d-h^{-1}\log x}{\sqrt x} .
\]
\end{proposition}
\begin{proof}
Use $F'(0+)=(\log d)/2$ and solve \eqref{rec:eq:ode} on $t\ge0$. Write
$S^*(x)=\sum_{p^k\le x}\log(p)\,a(p^k)(p^{-k}-x^{-1})$. The homogeneous solutions are
$\cosh(t/2)$ and $\sinh(t/2)$, the particular solution for the continuous forcing is
$-2t\sinh(t/2)/h$, and an atom of mass $B$ at $u>0$ adds $2B\sinh((t-u)/2)$ for $t\ge u$.
Consequently
\[
 \frac{F(\log x)}{\sqrt x}=S^*(x)-\frac{1-x^{-1}}h\log x+\frac{F(0)}2(1+x^{-1})
 +\frac{\log d}2(1-x^{-1}).
\]
The difference between $S^*(x)$ and the sum in $R(x)$ comes exclusively from powers of
$p\mid q$. It is $C_{\mathrm{bad}}+O_q((1+\log x)/x)$, where
$C_{\mathrm{bad}}=\sum_{p\mid q}\sum_{k\ge1}\log(p)\,a(p^k)p^{-k}$; this follows by bounding
a geometric tail for each such $p$ and at most $\log x/\log p$ terms of size $2\log(p)/x$.
Taking $x\to\infty$ gives $c+C_{\mathrm{bad}}+(F(0)+\log d)/2=0$, proving
\eqref{rec:eq:remainder}. In the lattice class the source proves $\eta(p)=1$ for every
$p\mid q$ (\code{New.Parameters.chiK_eq_one_of_dvd}). Then $C_{\mathrm{bad}}=0$ and every
correction at $p\mid q$ vanishes.
\end{proof}

\begin{example}\label{rec:ex:9239}
The pair $(q,d)=(18480,9239)$ lies in the lattice class. The two prime jumps at $258721$ and
$794641$ have $p-1=14q$ and $p-1=43q$, respectively, so their gcd is already $q$. Trial
division and Euler's criterion verify both primes and both nonresidues; the output is in
\artifact{dirichlet-grh-research/critical-line/prime-jumps-output.txt}.
\end{example}

\begin{corollary}[Oscillation at the critical scale; \stG]\label{rec:cor:oscillation}
In the general class the logarithmic mean of $\sqrt x\,R(x)$ is $b_0$, and for every
$\gamma\ne0$ with $b_\gamma\ne0$,
\[
 \liminf_{x\to\infty}\sqrt x\,R(x)\le b_0-\abs{b_\gamma},\qquad
 \limsup_{x\to\infty}\sqrt x\,R(x)\ge b_0+\abs{b_\gamma}.
\]
Infinitely many nonzero frequencies have $b_\gamma\ne0$.
\end{corollary}
\begin{proof}
The mean is obtained by termwise averaging. The averages of $F(t)$ times
$1\pm\cos(\gamma t)$ over $[0,T]$ tend to $b_0\pm b_\gamma$, and each weight is nonnegative
with mean one. Equation \eqref{rec:eq:remainder} transfers the bounds to the prime sum. The
cusp of $F$ at zero is not removed by any finite trigonometric polynomial, so infinitely
many nonzero frequencies survive. The mean includes any genuine central zero; the argument
uses neither simplicity nor linear independence of the ordinates.
\end{proof}

\begin{corollary}[Four central values at once; \stG]\label{rec:cor:chowla}
Take $q=8$ and let $d\equiv7\pmod8$ be prime. Write $\chi_D(n)=(D/n)$ for the primitive
quadratic character of fundamental discriminant $D$, and $R_{8,d}$ for the remainder above.
Then
\[
 B_d:=\lim_{X\to\infty}\frac1{\log X}\int_1^XR_{8,d}(x)\,\frac{\dd x}{\sqrt x}
 =-\sum_{D\in\set{-d,\,4d,\,-8d,\,8d}}\ord_{s=1/2}L(s,\chi_D)\ \in\ -2\Z_{\ge0}.
\]
Thus $B_d>-2$ is equivalent to the simultaneous nonvanishing at $1/2$ of these four
quadratic $L$-functions.
\end{corollary}
\begin{proof}
Equation \eqref{rec:eq:remainder} identifies the mean with $b_0$. The untwisted functions
are $\zeta$ and the $L$-functions of the characters of discriminants $-4,8,-8$. The
alternating eta series gives $\zeta(1/2)<0$, and the alternating beta series gives
$L(1/2,\chi_{-4})>0$. For $8$ and $-8$ the signs on each block of four successive odd
integers are respectively $(+,-,-,+)$ and $(+,+,-,-)$; the sums of these blocks are positive
because $y^{-1/2}$ is strictly decreasing and strictly convex. The twist is $\chi_{-d}$ and
the four resulting discriminants are the ones displayed. Finally each completed quadratic
$L$-function has root number $+1$, so each central order is even.
\end{proof}

Proving the strict lower bound $B_d>-2$ for every such $d$ would establish this subfamily of
Chowla's conjecture on central nonvanishing \cite{Soundararajan2000Nonvanishing}. The
one-sided estimate $R(x)\ll x^{-1/2+\eps}$ of the source allows every fixed value of $B_d$
and supplies no such bound.

\paragraph{Reading a jump from one high band.}
The reconstruction can discard every zero of height below $T$, and its errors can be made
explicit.

\begin{lemma}[\stG]\label{rec:lem:Fbound}
In the general class $\norm F_\infty\le6\log q+3\log d+48$. In the lattice class
$\norm F_\infty\le9\log d+60$.
\end{lemma}
\begin{proof}
For a primitive character of conductor $f$, let $\Xi$ be its completed function, with an
additional factor $s(s-1)$ in the zeta case. Its Hadamard logarithmic derivative and
functional equation give
\[
 \sum_\rho\Re\frac1{2-\rho}=\Re\frac{\Xi'}{\Xi}(2)\le\tfrac12\log f+4 .
\]
Indeed the constant in the real part of the Hadamard logarithmic derivative vanishes by
evaluating at a regular point on the critical line, where both real parts vanish by GRH and
reflection. At $2$ the gamma contribution is at most $1/2$, since the digamma function is
increasing and $\psi(2)<1$. The logarithmic Euler derivative has absolute value at most
$\sum_{n\ge2}(\log n)/n^2<2$ by the integral test, and the zeta correction is $1+1/2$.
Discard $-\tfrac12\log\pi$. On GRH, $(1/4+\gamma^2)^{-1}\le6\Re(2-\rho)^{-1}$. Adding the
two families therefore gives $\norm F_\infty\le6\log q+3\log d+48$. In the lattice class
$q\le2(d+1)\le4d$.
\end{proof}

Use the convention $\widehat g(s)=\int g(t)e^{-ist}\,\dd t$ in the next statement and set
\[
 g(t)=\cos(3t/2)\Bigl(\frac{\sin(t/16)}{t/16}\Bigr)^8,\qquad g(0)=1 .
\]
Each sinc factor has frequency support $[-1/16,1/16]$, so $\widehat g$ is supported on
$[-2,-1]\cup[1,2]$; it is an explicit piecewise polynomial, obtained by eight convolutions
of interval indicators followed by two translations. Define the band reading
\begin{equation}\label{rec:eq:AT}
 A_T(u)=-\frac1T\sum_{T\le\abs\gamma\le2T}m(\gamma)\,\widehat g(\gamma/T)\,e^{i\gamma u}.
\end{equation}

\begin{proposition}[Explicit band reading; \stG]\label{rec:prop:band}
Let $(q,d)$ be in the general class and $u\in\R$. Choose $0<r\le1$ such that $(u-2r,u+2r)$
contains no potential atom of \eqref{rec:eq:ode} except possibly $u$; at $u=0$ one may take
$r=\log(2)/3$, and at $u=\log p$ one third of
$\min\set{\log(p/(p-1)),\log((p+1)/p)}$. For $Tr\ge1$ put
\[
 B_u(T)=\frac{512}{7T}\cosh\bigl((\abs u+2r)/2\bigr),\qquad
 K_u(T)=\frac{2\cdot16^8}{7r^7T^8}\Bigl(108T^2+\frac{40T}r+\frac6{r^2}+\frac14\Bigr).
\]
Then
\begin{equation}\label{rec:eq:error}
 \abs{A_T(u)-J(u)}\le B_u(T)+\norm F_\infty K_u(T),
\end{equation}
and in the lattice class $\norm F_\infty$ may be replaced by $9\log d+60$.
\end{proposition}
\begin{proof}
Write $\mu=F''-\tfrac14F$. The reading $A_T(u)$ equals the pairing
$\langle\mu,g(T(\cdot-u))\rangle$, defined as
$\int F(t)(D^2-\tfrac14)[g(T(t-u))]\,\dd t$; termwise integration follows from the absolute
uniform convergence of $F$. Multiply the test by a cutoff equal to one on $\abs{t-u}\le r$,
zero on $\abs{t-u}\ge2r$, and between these equal to $1-10z^3+15z^4-6z^5$ at
$z=\abs{t-u}/r-1$. Its first two derivatives are bounded by $2/r$ and $6/r^2$. On the local
part only the selected mass and the density of \eqref{rec:eq:ode} contribute. Since
$\int\abs g\le2(16+16/7)=256/7$, the density costs at most $B_u(T)$. For the complementary
part integrate by parts onto the test. Writing $g(t)=16^8t^{-8}\cos(3t/2)\sin^8(t/16)$ shows
that for $\abs t\ge1$ its derivatives of orders $0,1,2$ are bounded by $16^8\abs t^{-8}$
times $1,10,108$, respectively: the trigonometric factor has frequency at most $2$ and total
absolute exponential coefficients at most $1$. The product rule and
$\int_{\abs{t-u}\ge r}\abs{t-u}^{-8}\,\dd t=2/(7r^7)$ give $K_u$. The cutoff is $C^2$, which
suffices by approximation, and the integrable differentiated tests justify passage to the
full real line. This last step permits localization even though the positive prime masses
and the smooth density on the right of \eqref{rec:eq:ode} individually grow exponentially at
infinity.
\end{proof}

\begin{remark}\label{rec:rem:fastband}
A test with vanishing moments removes the term $B_u(T)$ at the cost of explicit constants.
Take a real even $w\in C_c^\infty(\set{s:1<\abs s<2})$ with $\int_\R w=2\pi$ and
$g(v)=(2\pi)^{-1}\int w(s)e^{isv}\,\dd s$. Then $g(0)=1$ and $\int v^jg(v)\,\dd v=0$ for
every $j\ge0$, since $w$ vanishes near zero, and the same localization gives, for every
fixed $u\in\R$ and $A>0$ in the general class,
\[
 J(u)=-\frac1T\sum_{T<\abs\gamma<2T}m(\gamma)w(\gamma/T)e^{i\gamma u}+O_{q,d,u,w,A}(T^{-A}):
\]
the atom contributes $J(u)$, the smooth compactly supported density contributes
$O_A(T^{-A})$ by Taylor expansion and the vanishing moments, and away from $u$ the $L^1$
norms of the test and of its first two derivatives are $O_A(T^{-A})$ by Schwartz decay.
\end{remark}

\paragraph{Errors in the spectral data.}
The formulas quantify errors in the data. Let $E_T$ bound the total absolute missing,
inserted or altered weight and let $\delta_T\le1$ bound the displacement of matched
ordinates; use records in a padded band, so that every entry that can cross an endpoint is
included. If $M_T$ bounds the absolute weight of the matched entries, a valid absolute error
for the reading \eqref{rec:eq:AT} is
\begin{equation}\label{rec:eq:noise}
 e_u(T)=\frac{\norm{\widehat g}_\infty E_T}T+\frac{M_T\delta_T}T
 \Bigl(\abs u\,\norm{\widehat g}_\infty+\frac{\norm{\widehat g'}_\infty}T\Bigr)
 +e_{\mathrm{arithmetic}},
\end{equation}
because $\widehat g(v/T)e^{ivu}$ is Lipschitz in $v$ with constant
$\norm{\widehat g'}_\infty/T+\abs u\norm{\widehat g}_\infty$. Both norms can be bounded from
the explicit spline. For each fixed packet $M_T=O_{q,d}(T\log(2+T))$ by zero counting, so
$E_T=o(T)$ and $\delta_T=o(1/\log T)$ preserve every fixed jump in the limit; for actual
certificates $M_T$ can be bounded by the observed absolute weight plus the certified weight
budget. The readings use signed entries and do not require pairing two nearby ordinates
before summation.

\begin{corollary}[\stG]\label{rec:cor:discrepancy}
Two distinct parameter pairs of the general class have absolute signed multiplicity
discrepancy $\gg T$ in each sufficiently high dyadic band $T<\abs\gamma<2T$.
\end{corollary}
\begin{proof}
Their jump functions differ at some point by \cref{rec:thm:jumps}, and \cref{rec:rem:fastband}
at that point bounds the discrepancy from below.
\end{proof}

\begin{theorem}[A certified stopping rule from one band; \stG]\label{rec:thm:band}
Let $(q,d)$ be in the lattice class. From a sufficiently high band $T\le\abs\gamma\le2T$ of
the aggregate signed data one can recover $(q,d)$ with a computable certificate. No upper
bound for $q$ or $d$ is part of the input. On an insufficient band the procedure may return
\emph{undetermined}. If the band readings \eqref{rec:eq:AT} have certified absolute errors
$e_u(T)$ tending to zero for each fixed observation site $u$, the procedure succeeds on
every sufficiently high band.
\end{theorem}
\begin{proof}
Let $a$ be the real part of the observed $A_T(0)$ and write $K=K_0(T)$,
$D=B_0(T)+60K+e_0(T)$. When $9K<1$, \eqref{rec:eq:error} confines $\log d$ to
\[
 \Bigl[\frac{a-D}{1+9K},\ \frac{a+D}{1-9K}\Bigr].
\]
Use outward-rounded arithmetic and test the primes in the exponential of this interval. If
there is a unique eligible prime, it is $d$; otherwise decline this band. For fixed true $d$
the interval shrinks to $\log d$. List the divisors $q'$ of $2(d+1)$ with $18480\mid q'$ and
$d\equiv q'/2-1\pmod{q'}$.

For every pair of distinct candidate moduli, find a prime $p$ coprime to all candidates and
to $d$, with $\eta(p)=-1$, for which their indicators in \eqref{rec:eq:mass} differ.
Enumerating primes terminates: choose a prime power $\ell^a$ occurring to a higher exponent
in one modulus than in the other. Modulo the least common multiple of the candidates choose
$1$ at all other components and a unit congruent to $1$ modulo the lower $\ell$-power but
distinct from $1$ modulo $\ell^a$. If the lower power is $1$, then $\ell$ is odd, since $8$
divides each candidate. Set a nonresidue at $d$ independently. Dirichlet's theorem supplies
infinitely many primes in this class. For each witness intersect the two possible masses,
$0$ and $2\log p/\sqrt p$, with the certified interval around the observed
$\Re A_T(\log p)$ from \eqref{rec:eq:error} plus the measurement error. Retain compatible
moduli. Each witness eventually separates its pair because $p$ is fixed and $B_u,K_u,e_u$
tend to zero. A surviving singleton is certified.
\end{proof}

A single accepted band certifies both integers. This assumes a genuine packet in the stated
parameter class; authenticating a putative spectrum is a separate problem. The affine bound
\eqref{rec:eq:error} is what supplies a test for accepting a band before either parameter is
known.

\begin{example}\label{rec:ex:band}
Take $d=46199$ and either $q=18480$ or $q=92400$. A genuine band at $T=10^{12}$ suffices,
even with absolute reading error $10^{-7}$ at each of the two sites. Use $r=1/8$ at zero and
$r=1/(4p)$ at $p=110881$. The origin interval contains only $d$, which leaves exactly these
two candidate moduli. The witness $p$ has $\eta(p)=-1$ and $p-1=6\cdot18480$. The two
possible masses differ by more than $2/334$. This is an a priori guarantee for either
packet; no zero ordinate at height $10^{12}$ was computed. The error margins are verified
by rational arithmetic in the following files.
\begin{center}
\artifact{exponent_seven_and_grh_transfers/other-results/check_decoder_example.py}\\
\artifact{exponent_seven_and_grh_transfers/other-results/decoder_example.txt}
\end{center}
\end{example}

\paragraph{Recovery from an initial segment, with a height bound.}
When the signed list is known through a height $T$, probes that annihilate the smooth
density of \eqref{rec:eq:ode} exactly give a bound for the height that suffices.

\begin{theorem}[Finite recovery with a height bound; \stG]\label{rec:thm:recovery}
Let $(q,d)$ be in the lattice class. Knowing the signed list $\set{(\gamma,m(\gamma))}$
through a sufficiently high finite height determines $q$ and $d$. The algorithm needs no
prior numerical upper bound on either. For approximate input it accepts enclosures for the
finite sums it uses; with complete lists and arbitrarily refinable certified errors it
terminates. For exact input a sufficient height is $O(d^2\log^3(2d))$ with an absolute
implied constant.
\end{theorem}

The height estimate is independent of the cost of computing or certifying the input lists.

\begin{proof}
\emph{Probes.} In this proof $\widehat g(\gamma)=\int g(t)e^{i\gamma t}\,\dd t$, as in
\eqref{rec:eq:weil}. If a real compactly supported $g$ kills the $\cosh(t/2)$ integral, has
value one at $u$, and its support isolates $u$ among the potential atoms, then by
\eqref{rec:eq:weil} the quantity $-\sum_\gamma m(\gamma)\widehat g(\gamma)$ equals $J(u)$,
with $J$ as in \eqref{rec:eq:mass}. The origin and prime identities alone would only give
an asymptotic reconstruction; to certify a finite stopping time one must estimate the unseen
zero contribution before knowing the conductors.

\emph{A bound containing the unknown on both sides.} For a primitive character of conductor
$f$ and parity $a\in\set{0,1}$, let $B_f$ denote the sum of $(1+\gamma^2)^{-1}$ over its
nontrivial zeros, with multiplicity. Hadamard factorization and the functional equation give
at $s=2$
\[
 \sum_\gamma\frac{3/2}{9/4+\gamma^2}=\Re\frac{L'}L(2)+\tfrac12\log(f/\pi)
 +\tfrac12\psi(1+a/2)+\tfrac32\mathbf 1_{f=1}.
\]
This is the real logarithmic derivative of the completed entire function, with $s(s-1)$
included for zeta \cite{LamzouriLiSoundararajan2015Conditional}; GRH gives the indicated
denominators. Multiplying the summand by $3/2$ majorizes $(1+\gamma^2)^{-1}$. For $f>1$,
$\abs{L'/L(2)}\le\sum_{n\ge2}(\log n)/n^2<1$ (isolate $2$ and $3$, integrate from $3$). For
$f=1$ the logarithmic derivative is negative. Since $\psi(1+a/2)\le\psi(2)<\log\pi$, in
both cases $B_f\le\tfrac34\log f+\tfrac94$. Averaging absolute values before taking the
signed difference yields
\[
 W:=\sum_\gamma\frac{\abs{m(\gamma)}}{1+\gamma^2}
 \le\tfrac32\log q+\tfrac34\log d+\tfrac92\le\tfrac94\log d+7,
\]
using $q\le3d$. For any chosen test $g$, suppose $A$ bounds
$\sup_{\abs\gamma>T}(1+\gamma^2)\abs{\widehat g(\gamma)}$, and let $[\ell_-,\ell_+]$ enclose the real
part of the negative sum through height $T$. The imaginary part of the exact full sum
vanishes. For the origin probe put $y=\log d$. Whenever $9A/4<1$, the arithmetic certificate
is
\begin{equation}\label{rec:eq:certificate}
 \frac{\ell_--7A}{1+9A/4}\ \le\ y\ \le\ \frac{\ell_++7A}{1-9A/4}.
\end{equation}
Exponentiate with outward rounding and stop the first stage when there is just one
admissible prime in the resulting interval. After that stage the error for each new probe
is at most $A(\tfrac94\log d+7)$, independent of $q$. A finite interval can likewise be
certified by summing enclosures of $m(\gamma)\widehat g(\gamma)$; omitted or uncertain
entries require an enclosure of their possible total contribution. Without such an input
error guarantee a list may omit a decisive term.

\emph{Two explicit probes.} A simple choice at location $v$ is the translate by $v$ of
\[
 g_a(t)=\frac{f_a''(t)-f_a(t)/4}{f_a''(0)-1/4},\qquad f_a(t)=\bigl(1-(t/a)^2\bigr)_+^8 .
\]
It equals one at its center and kills both $e^{t/2}$ and $e^{-t/2}$ by integration by parts.
Use $a=1/2$ at zero, or $a=1/[3(p+1)]$ at $\log p$; the distance to the nearest other
logarithm of a positive integer is at least $\log(1+1/p)\ge1/(p+1)$. Four integrations by
parts permit $A=\norm{g_a^{(4)}}_1(1+T^2)/T^4$.

For the sharper height estimate, take $b=1/6$ at zero and $b=1/[10(p+1)]$ at $\log p$. Set
$k=\floor{bT/2}\ge4$. Convolve the uniform probability density on $[-b/k,b/k]$ with itself
$k$ times and divide by its value at zero to obtain $w$. It is symmetric and unimodal
(averaging any symmetric unimodal density on a centered interval preserves this property).
Before normalization its maximum is at least $1/(2b)$ by total mass and support. Define
\[
 g(t)=w(t-v)-\frac{w(t-v-2b)+w(t-v+2b)}{2\cosh b}.
\]
The support lies within $3b$ of the center, $g(v)=1$, and translation of $e^{\pm t/2}$ by
$\pm2b$ cancels its integral. The Fourier transform of the unnormalized $w$ is
$\recsinc(b\gamma/k)^k$, where $\recsinc(x)=\sin(x)/x$. For $\abs\gamma>T$ its absolute
value is at most $2^{-k}(T/\abs\gamma)^k$. The translation multiplier is at most two, so one
can use
\[
 A=4b(1+T^2)\,2^{-\floor{bT/2}} .
\]
The probes have sufficient differentiability for the explicit formula: convolution with
smooth approximate identities gives the same identity in the limit, with spectral domination
by a multiple of $(1+\gamma^2)^{-1}$. Increasing $T$ and refining enclosures therefore
certifies $d$ at the origin; one can, for example, choose $T$ dyadically and refine the
finite-sum enclosure to width $1/T^2$ or smaller. For exact data $T=O(\log(2d))$ suffices at
the origin, since adjacent integer logarithms near $d$ are separated by $\gg1/d$.

\emph{Separating the finitely many moduli.} Form $M=2(d+1)$. The remaining candidates are
the divisors $q$ of $M$ with $18480\mid q$ and $M/q$ odd. For two distinct candidates
$q_1<q_2$ there is a unit $y\bmod M$ congruent to one modulo $q_1$ and different from one
modulo $q_2$: choose a prime $\ell$ with $v_\ell(q_1)<v_\ell(q_2)$, prescribe the residue
$1+\ell^{v_\ell(q_1)}$ modulo the full $\ell$-power in $M$, and one at all other prime
powers. If $v_\ell(q_1)=0$, this residue is two and $\ell$ is odd. Now require
$p\equiv y\bmod M$ and $\eta(p)=-1$. These conditions specify one coset of index
$2\varphi(M)$ modulo $Md$, since $(M,d)=1$. Dirichlet's theorem gives a prime in this
coset, and a search with integer tests finds one. The two possible values of $J(\log p)$
are separated by $2\log p/\sqrt p$. A certified interval excluding one value eliminates a
candidate. Continuing until only one remains proves finite recovery.

\emph{The height.} Theorem~1.4 of \cite{LamzouriLiSoundararajan2015Conditional} bounds the
least prime in a coset of index $I>1$ modulo $Q\ge20000$ by $O(I^2\log^2Q)$ under Dirichlet
GRH; the explicit alternatives are $p\le10^9$ or
$p\le((I-1)\log Q+3(I+1)+\tfrac52(\log\log Q)^2)^2$. Here $I\le2M$ and $Q=Md$, so each
witness can be taken with $p\ll d^2\log^2(2d)$. For a sufficiently large absolute constant,
$T$ of order $p\log(2pd)$ makes the spline tail times $W$ less than half the two-value
separation. All witnesses use the same spectral list, so the maximum of these heights
proves the claimed bound.
\end{proof}

\begin{example}\label{rec:ex:recovery}
Take $d=109909799$. Sixteen candidate moduli survive the divisor tests. Even knowing the
totient would leave the pair $1127280,1201200$ at totient $230400$. The prime $34945681$ is
a quadratic nonresidue modulo $d$, and $\gcd(p-1,2(d+1))=1127280$; its jump distinguishes
the two. At $T=768$ the spline at the origin has a tail below $1.045\cdot10^{-12}$ for this
$d$. At $T=6\cdot10^{10}$ the witness probe has a tail below $5.220\cdot10^{-11}$, against a
jump greater than $0.005875$. These are calculations on known parameters; an observed
spectral list of these heights was not computed. Primes, characters, candidate enumeration
and tail arithmetic are checked in the following files.
\begin{center}
\artifact{grh-torsion-and-zero-data/check_recovery.py}\\
\artifact{grh-torsion-and-zero-data/verification/check_recovery.txt}
\end{center}
\end{example}

The band rule of \cref{rec:thm:band} and the initial-segment rule of
\cref{rec:thm:recovery} use different data and neither contains the other: the first
discards all zeros below $T$ and comes with the noise model \eqref{rec:eq:noise}, the second
uses the whole list through $T$ and comes with an explicit sufficient height. Both are
confined to the lattice class, where $q\mid2(d+1)$ leaves finitely many candidate moduli
once $d$ is known; \cref{rec:thm:jumps} and \cref{rec:rem:fastband} cover the general class
in the limit.




\subsection{Omitted primes from normalized zero spacings}\label{rec:sec:gaps}
\sources{dirichlet_zero_recovery}

The zeros of an imprimitive $L$-function on $\Re s=0$ come from the omitted Euler factors. Let
$\chi$ modulo $q$ be induced by $\chi^*$ of conductor $f$, and $P=\set{p\mid q:p\nmid f}$, nonempty,
$m=\abs P$. Then $L(s,\chi)=L(s,\chi^*)\prod_{p\in P}(1-\chi^*(p)p^{-s})$, and by the functional
equation and nonvanishing on $\Re s=1$ every zero on the positive imaginary axis comes from the
finite product. Writing $\chi^*(p)=e^{2\pi i\alpha_p}$ with rational $\alpha_p\in[0,1)$, the
ordinates contributed by $p$ are the positive points of $(\Z+\alpha_p)/\lambda_p$,
$\lambda_p=\log p/(2\pi)$; these lattices are disjoint at positive height by unique factorization,
and each factor has simple zeros. So the intensity is $\Lambda=\frac1{2\pi}\sum_{p\in P}\log p$.
List the positive ordinates as $t_1<t_2<\cdots$, put $g_j=\Lambda(t_{j+1}-t_j)$ and
$S_P(x)=\lim_N\frac1N\#\set{j\le N:g_j>x}$; common rescaling of the heights leaves $S_P$ unchanged.
Put $w_p=\log p/\sum_{r\in P}\log r$, $b=1/\max_pw_p$ and $A_P(x)=\prod_{p\in P}(1-w_px)$.

\begin{proposition}[Gap law for the omitted factors; \stU]\label{rec:prop:gaplaw}
For every $x\ge0$ the limit $S_P(x)$ exists, is independent of the phases $\alpha_p$, and equals
$-A_P'(x)$ for $0\le x<b$ and $0$ for $x\ge b$. If $p_*$ is the largest prime of $P$, the gap law has
an atom at $b$ of mass $w_{p_*}\prod_{r\ne p_*}(1-w_r/w_{p_*})$.
\end{proposition}

\begin{proof}
This is the Barra--Gaspard spacing law for superposed lattices, as stated by Marklof and
Str\"ombergsson \cite{MarklofStrombergsson2013Gaps}. Directly: from the $n$th point
$(n+\alpha_p)/w_p$ of the normalized $p$-lattice, its position in the $r$-lattice is the fractional
part of $nw_r/w_p+\alpha_pw_r/w_p-\alpha_r$. Since $1$ and the $w_r/w_p$, $r\ne p$, are rationally
independent (a relation would give an integer relation between logarithms of distinct primes),
these positions equidistribute in the unit cube by Weyl's criterion. For $0\le x<b$ the next
$r$-point lies farther than $x$ with limiting frequency $1-w_rx$, jointly for all $r$, and the next
$p$-point is $1/w_p>x$ away; starting points of type $p$ have frequency $w_p$. Hence
$S_P(x)=\sum_pw_p\prod_{r\ne p}(1-w_rx)=-A_P'(x)$. Every gap is at most $b$, and the left limit at $b$
gives the atom.
\end{proof}

\begin{theorem}[Recovery from short gaps; \stU]\label{rec:thm:gaps}
Let $\chi,\widetilde\chi$ have nonempty omitted sets $P,Q$, and suppose $S_P=S_Q$ on $(0,\eps)$ for
some $\eps>0$. Then $\abs P=\abs Q$, and if $\abs P\ge3$, then $P=Q$. No absolute density, modulus,
conductor or character value is used.
\end{theorem}

\begin{proof}
Near zero $A_P(x)=1-\int_0^xS_P$ is a polynomial, so the germ determines $A_P$, whose degree is $m$
and whose roots are the $1/w_p$. Hence the sorted weights agree, and with
$c=\sum\log q_i/\sum\log p_i$ the ordered primes satisfy $q_i=p_i^c$. For $m\ge3$, if $c$ were
irrational, the six exponentials theorem \cite{Waldschmidt2005Further} applied to
$\log p_1,\log p_2,\log p_3$ and $1,c$ would make one of the primes $p_i,q_i$ transcendental. So
$c=a/d$ in lowest terms, $q_i^d=p_i^a$, and unique factorization gives $q_i=p_i$.
\end{proof}

For a squarefree principal modulus with at least three prime divisors $P$ is its prime divisor set,
so the short-gap law recovers the modulus. For one prime the normalized law is a point mass at $1$;
for two primes $p_1<p_2$ and $q_1<q_2$ the laws agree exactly when
$\log q_1/\log p_1=\log q_2/\log p_2$, which the four exponentials conjecture would rule out for
distinct pairs. The inverse map is not continuous.

\begin{proposition}[Approximation by unrelated prime sets; \stU]\label{rec:prop:gapinstab}
For every $P$ with $m\ge2$ there are pairwise disjoint sets $Q_n$ of $m$ primes, disjoint from $P$,
with $\min Q_n\to\infty$, whose gap laws converge weakly to that of $P$ and whose short-gap
polynomials converge coefficientwise.
\end{proposition}

\begin{proof}
Order the weights $w_1<\dots<w_m$ of $P$ and, for large $T$, choose by Bertrand's postulate
primes $q_i\in[\floor{e^{Tw_i}},2\floor{e^{Tw_i}}]$; then $\log q_i=Tw_i+O(1)$ and the
normalized weights are $w_i+O_m(1/T)$. Let $T$ grow fast enough along a sequence; the cutoffs
converge to $b$ and the survival functions converge at every $x\ne b$.
\end{proof}

\begin{example}\label{rec:ex:gaps30}
The principal character modulo $30$ and the character modulo $210$ induced by the quadratic
character modulo $7$ both have $P=\set{2,3,5}$, with phases $(0,0,0)$ and $(0,\tfrac12,\tfrac12)$.
With $e_j$ the elementary symmetric functions of $\log2/\log30,\log3/\log30,\log5/\log30$,
\[
 S_P(x)=1-2e_2x+3e_3x^2\approx1-0.63021791x+0.09344797x^2\qquad(0\le x<b),
\]
with $b=\log30/\log5\approx2.11328275$, and about $8.55\%$ of the gaps equal $b$. The empirical
frequencies of $541{,}314$ and $541{,}316$ gaps below height $10^6$ agree with the law within
$1.2\cdot10^{-5}$ at the six tested events; see
\begin{center}
\artifact{dirichlet_zero_recovery/gaps/verify_gaps.py}\\
\artifact{dirichlet_zero_recovery/gaps/example_output.txt}
\end{center}
\end{example}

\subsection{Characters and multiplicity vectors in short windows at large height}\label{rec:sec:localwin}
\sources{local_euler_recovery; translated_dirichlet_zero_divisors; dirichlet_prime_frequency_inverse; prime_pole_windows; grh-moving-zero-windows}

The zeros are now observed only in one window $\abs{\Im\rho-T}\le H$ at a large height $T$. A test
function holomorphic across the closed critical strip and decaying fast enough there makes the
Guinand--Weil formula local, and the radius needed to suppress the unseen zeros is
$\pi^{-1}\log\log\log T+O(1)$; windows centred at $T$ of half-width $\pi/(2\log U)+O((\log U)^{-2})$, $U=\log(q(2+T)/\pi)$, already
contain $\gg U/\log^2U$ zeros (\cref{rec:prop:movingwin}), so the larger radius serves to suppress the unseen zeros. One family of tests reads a single Euler coefficient uniformly in the
conductor; a second, binomial, family gives full rank and both signs of every balanced
combination at the same radius, and survives translation of the functions. Several statements
use complex zeros and no hypothesis on their location. In this subsection
$\widehat f(v)=\int_\R f(t)e^{-ivt}\dd t$, $z_\rho=(\rho-\tfrac12)/i$, and for a primitive $\chi$ of
conductor $q$ and parity $a$,
\[
 G_\chi(t)=\log(q/\pi)+\Re\frac{\Gamma'}\Gamma\Bigl(\frac{1/2+a+it}2\Bigr),\qquad
 G_a^\circ(t)=\Re\frac{\Gamma'}\Gamma\Bigl(\frac{1/2+a+it}2\Bigr).
\]
For tests $h$ holomorphic and rapidly decreasing in a strip wider than $\abs{\Im z}\le\tfrac12$ the
explicit formula reads, with each zero counted with multiplicity
\cite{ChirrePereiraDeLaat2021Primes,Davenport2000Multiplicative},
\begin{equation}\label{rec:eq:gwstrip}
 2\pi\sum_\rho h(z_\rho)=2\pi\mathbf 1_{q=1}\bigl\{h(\tfrac i2)+h(-\tfrac i2)\bigr\}+\int_\R h(t)G_\chi(t)\dd t
 -\sum_{n\ge2}\frac{\Lambda(n)}{\sqrt n}\bigl\{\chi(n)\widehat h(\log n)+\overline{\chi(n)}\widehat h(-\log n)\bigr\}.
\end{equation}
It needs no hypothesis on the zeros; complex tests are allowed by linearity, and the tests below
are reached from compactly supported smooth Fourier transforms by truncation, using the
unit-interval zero count.

\paragraph{One Euler coefficient from one window.}
For $Q\ge2$ and $T\ge2$ put
\[
 \ell=\log\bigl(e^eQ(T+3)\bigr),\qquad R=32\log\ell,\qquad D=D(Q,T)=2+\pi^{-1}\log(4R),
\]
so that $D=\pi^{-1}\log\log\log(QT)+O(1)$. With $w(z)=\cosh(\pi z)$ let
\[
 k_R(z)=c_R\exp\Bigl[-\cosh(\pi z/2)-\frac{Rw(z)}{R+w(z)}\Bigr],
\]
with $c_R$ chosen so that $\int_\R k_R=1$. For a prime $p$ and an integer $m\ge4$ put $u=\log p$,
$b=8p$, $J_b=\mathbf 1_{[-b,b]}/(2b)$, $g=J_b^{*m}$, $K=g*k_R$, and
\[
 h(z)=\frac1{iu}\frac{\dd}{\dd z}\bigl(e^{iuz}K(z-T)\bigr),\qquad
 \mathcal E_p=\frac{\sqrt p}{\log p}\Bigl[\int_\R h(t)G_0^\circ(t)\dd t-2\pi\sum_{\abs{\Im\rho-T}\le H}h(z_\rho)\Bigr].
\]
The integral is known before any zero is observed; for a product of $k$ factors it is multiplied
by $k$.

\begin{theorem}[Local Euler coefficient; \stU]\label{rec:thm:localcoeff}
Let $q_\chi\le Q$, let $p$ be prime and $m\ge4$ an integer, and put $B=8pm$, $H=2B+D(Q,T)$, with
$T\ge H+2$. Then
\[
 \abs{\mathcal E_p-\chi(p)}\le C\,2^{-m}
\]
for an absolute constant $C$. The only $\chi$-dependent input is the multiset of complex zeros
$\rho$ with $\abs{\Im\rho-T}\le H$; under \cref{src:thm:dgrh} their ordinates with multiplicities
suffice. Ramified primes, where $\chi(p)=0$, are included.
\end{theorem}

\begin{lemma}[Decay on the closed strip; \stU]\label{rec:lem:lekernel}
Uniformly for $R\ge2$ and real $v$, $\abs{\widehat k_R(v)}\ll_Ne^{-\abs v/2}(1+\abs v)^{-N}$ for each
$N\ge0$. For fixed $M>0$, $j=0,1$ and $z=x+iy$ with $\abs y\le\tfrac12$,
\[
 \abs{k_R^{(j)}(z)}\ll_Me^{-M\abs x},\qquad \abs{k_R^{(j)}(z)}\ll_Me^{-R/4}e^{-M\abs x}\quad(\abs x\ge D).
\]
\end{lemma}

\begin{proof}
In this strip $\Re w\ge0$, $\abs{R+w}\ge\max(R,\abs w)$ and
$\Re\frac{Rw}{R+w}=R\frac{R\Re w+\abs w^2}{\abs{R+w}^2}\ge0$; when $\abs w\ge R$ the last expression is at
least $R/4$, and $\abs{w(x+iy)}\ge\abs{\sinh(\pi x)}$ covers $\abs x\ge D$. Also
$\Re\cosh(\pi z/2)\ge2^{-1/2}\cosh(\pi x/2)$. The normalizer and its reciprocal are bounded
uniformly (on $\abs x\le1$ use $Rw/(R+w)\le w$). Each derivative of fixed positive order of
$Rw/(R+w)=R-R^2/(R+w)$ is $O_j(\exp(C_j\abs x))$ on the strip uniformly in $R$, and the
double-exponential first factor absorbs it. This gives the pointwise bounds and uniform $L^1$
bounds for derivatives of each fixed order on both edges of the strip. Shift the Fourier
contour to $\Im z=-\operatorname{sgn}(v)/2$ and integrate by parts. For fixed $R$ the kernel is
analytic and rapidly decreasing on $\abs{\Im z}\le\tfrac34$, since the nearest zeros of $R+w$ have
$\abs{\Im z}=1$; this justifies the shift and the use of \eqref{rec:eq:gwstrip}.
\end{proof}

\begin{lemma}[Frequency response; \stU]\label{rec:lem:lefreq}
One has $\widehat h(0)=0$, $\widehat h(u)=1$,
$\widehat h(v)=\frac vue^{-i(v-u)T}\widehat k_R(v-u)\recsinc(b(v-u))^m$, and
\[
 \frac{\sqrt p}u\Bigl(\sum_{n\ge2,\,n\ne p}\frac{\Lambda(n)}{\sqrt n}\abs{\widehat h(\log n)}
 +\sum_{n\ge2}\frac{\Lambda(n)}{\sqrt n}\abs{\widehat h(-\log n)}\Bigr)\ll2^{-m}.
\]
\end{lemma}

\begin{proof}
Differentiation, translation and convolution give the formula. For $p/2\le n\le2p$, $n\ne p$,
$b\abs{\log(n/p)}\ge4\abs{n-p}$, the factors $(\sqrt p/\sqrt n)(\Lambda(n)/u)(\log n/u)$ are bounded,
and $\abs{\widehat k_R}\le1$, so these terms contribute $O(\sum_{j\ge1}(4j)^{-m})$. In the remaining
sums $\abs{v-u}\ge\log2$, and before the sinc factors
$\sum_n\frac{\Lambda(n)}{\sqrt n}\frac{\log n}ue^{-\abs{\log n-u}/2}(1+\abs{\log n-u})^{-4}\ll\sqrt p$: for
$n\le p$ bound it by $p^{-1/2}\sum_{n\le p}\Lambda(n)$, and on $2^jp<n\le2^{j+1}p$ by
$O(\sqrt p(1+j)^{-3})$ by Chebyshev; negative frequencies have the same majorant since
$\abs{-\log n-u}\ge\abs{\log n-u}$. Two sinc factors give $O(p^{-2})$ and each further one is at most
$1/2$.
\end{proof}

\begin{proof}[Proof of \cref{rec:thm:localcoeff}]
In \eqref{rec:eq:gwstrip}, $\widehat h(0)=0$ kills the conductor term, and \cref{rec:lem:lefreq} leaves
the positive frequency $\log p$, with value $\chi(p)\log p/\sqrt p$, up to $O(2^{-m})$ after
normalization. \emph{Unseen zeros.} The density $g$ is supported in $[-B,B]$ with
$\norm g_\infty\le(2b)^{-1}$, so $\abs K+\abs{K'}\ll b^{-1}$ on the closed strip, and if
$d_x=\abs{x-T}-B\ge D$, \cref{rec:lem:lekernel} gives
$\abs{K(x+iy-T)}+\abs{K'(x+iy-T)}\ll_Mb^{-1}e^{-R/4}e^{-Md_x}$ for $\abs y\le\tfrac12$. The modulation
costs at most $\sqrt p$ on the strip. A unit interval at height $y$ contains
$O(\log(q_\chi(\abs y+3)))$ zeros \cite{BennettMartinOBryantRechnitzer2021Counting}; on the $j$th unit
interval outside the window $d_x\ge B+D+j$ and the count is $O(\ell+j+1)$, since $H<T$. The omitted
contribution is therefore $\ll\frac p{bu}e^{-R/4}\sum_{j\ge0}(\ell+j+1)e^{-(B+D+j)}\ll e^{-B}$, and the
argument keeps $\Re\rho$ throughout. \emph{Parity and poles.} Reflection for $\psi$ at
$z=\tfrac14+it/2$ gives $G_1^\circ-G_0^\circ=\pi\operatorname{sech}(\pi t)$. As $\int\abs h\ll1$ and
$\operatorname{sech}(\pi t)\ll e^{-\pi T/4}$ off $\abs t\le T/4$, where the box support is at distance at least
$T/4$, the parity error is $O(e^{-cT})$. At $z=\pm i/2$ the distance from the support is at least
$T/2$, so the pole readings are $O(e^{-T/2})$; with $M=4$ these errors are $O(2^{-m})$ since
$T\ge16pm$.
\end{proof}

\begin{theorem}[Identification and factorization from one window; \stG]\label{rec:thm:localid}
There are absolute constants $C_0,C_1$. Fix $d\ge1$, put $U=d^2\log^2(2Q)\log(2d)$ and
$H=C_0U+D(Q,T)$. If $T\ge C_1U$, the multiset of ordinates in $[T-H,T+H]$ of the zeros of
$\prod_{j=1}^kL(s,\chi_j)$, $0\le k\le d$, $q_{\chi_j}\le Q$, determines $k$ and the multiset
$\set{\chi_1,\dots,\chi_k}$ of primitive characters, repetitions allowed, from the total
multiplicities alone.
\end{theorem}

\begin{proof}
Subtract two putative products as multisets of factors, leaving distinct primitive
$\chi_1,\dots,\chi_s$, $s\le2d$, with nonzero integer differences $a_j$, $\norm a_1\le2d$. Take $X=Ad^2\log^2(2Q)$
with $A$ large and a smooth $W\ge0$ supported in $(1,2)$ with $\int W=1$. By \cref{rec:lem:gram},
\[
 \frac1X\sum_p(\log p)W(p/X)\chi_i(p)\overline{\chi_j(p)}=\mathbf 1_{i=j}+O\bigl(\log(2Q)/\sqrt X\bigr),
\]
so the Gram matrix
differs from the identity entrywise by at most $1/(2s)$, and the weighted mean square of
$\sum_ja_j\chi_j(p)$ is at least $\tfrac12\norm a_2^2$. The total weight divided by $X$ is bounded, so some prime $X<p<2X$ has
\begin{equation}\label{rec:eq:localsep}
 \Bigl\lvert\sum_ja_j\chi_j(p)\Bigr\rvert\ge c_*\norm a_2\ge c_*>0,
\end{equation}
with $c_*$ absolute, also when $p$ ramifies in some factors. Take $m=\ceil{A'\log(2d)}\ge4$ with
$2dC2^{-m}<c_*/2$; for $C_0,C_1$ large the window of \cref{rec:thm:localcoeff} for this $p$ lies in
$[T-H,T+H]$ and $T\ge16pm+D+2$. To allow different degrees use only the zero sum in $\mathcal E_p$:
integrating by parts, $\bigl\lvert\frac{\sqrt p}u\int hG_0^\circ\bigr\rvert\le\frac{\sqrt p}{u^2}\int K(t-T)\abs{G_0^{\circ\prime}(t)}\dd t
\ll\sqrt p/(Tu^2)$, since $G_0^{\circ\prime}(t)\ll(1+\abs t)^{-1}$. Each zero-only reading differs from the
corresponding sum of character values by at most $d(C2^{-m}+C'\sqrt p/(Tu^2))<c_*/2$ after enlarging
$A',C_1$, and equal spectra in the window would contradict \eqref{rec:eq:localsep}, including against
the empty product.
\end{proof}

For $d=1$ the conductor, parity, order and all Euler factors are determined, and a finite decoder
enumerates candidates. A partial matching of two lists pairs occurrences bijectively; let $N$ be
the number of unpaired occurrences and $S$ the sum of the displacements of the paired ones.

\begin{corollary}[Local matching cost; \stG]\label{rec:cor:localmatch}
For distinct primitive characters of conductor at most $Q$ and the window of \cref{rec:thm:localid}
with $d=1$, every partial matching of their lists satisfies
$S+N/\log(2+\log(2Q))\gg\log(2Q)$. With only exact coincidences paired, the symmetric difference
counted with multiplicity is $\gg\log(2Q)\log(2+\log(2Q))$.
\end{corollary}

\begin{proof}
Use the prime of \eqref{rec:eq:localsep}, $p\asymp\log^2(2Q)$, and the statistic with its sum extended
to the whole window; the tail estimate covers the complement of the larger window as well. The
readings differ by at least $c_*/2$, while $\norm h_\infty\ll p^{-1}$ and $\norm{h'}_\infty\ll(\log p)/p$ on
the real axis. Hence $1\ll(N/u+S)/\sqrt p$.
\end{proof}

At $T=1000$ the statistic, computed in floating point from self-generated zero lists of $\zeta$,
$L(s,\chi_{-4})$ and the quartic character modulo $5$ with $\chi(2)=i$, reproduces $\chi(p)$ for
$p=2,3,5,7$ with errors below $10^{-4}$, and four zero-only readings select
$\zeta(s)L(s,\chi)^2$ among $20$ candidate products, with residual $0.000203$ against $1.000077$ for
the next candidate;
\begin{center}
\artifact{local_euler_recovery/numerical/local_readout.py}\\
\artifact{local_euler_recovery/numerical/output.txt}
\end{center}

\paragraph{A binomial cutoff and full rank at the endpoint.}
Fix distinct primitive characters $\chi_1,\dots,\chi_d$ of conductors $q_j$, $Q=\lcm(q_j)$, and put
$m(\rho)=(\ord_\rho L(s,\chi_j))_j$ for a nontrivial zero of at least one of them, and
$m(\gamma)=m(\tfrac12+i\gamma)$ under \cref{src:thm:dgrh}. Write $\log_3T=\log\log\log T$. Constants
may depend on the family.

\begin{lemma}[A cutoff across the critical strip; \stU]\label{rec:lem:binomkernel}
Fix $W\ge1$ and an integer $J\ge4$. Let $k_M(z)=c_M\operatorname{sech}^2(z)\bigl(1+M^{-1}\cosh(\pi z)\bigr)^{-M}$ with
$\int_\R k_M=1$, and let $K=K_{M,W}$ be the convolution of $k_M$ with $J$ uniform probability
densities on $[-W/J,W/J]$. Uniformly in integers $M\ge2$,
\[
 \abs{K(x+iy)}\ll_{W,J}e^{-2\abs x},\qquad \abs{K(x+iy)}\ll_{W,J}2^{-M/2}e^{-2\abs x}\quad
 \bigl(\abs x\ge W+\pi^{-1}\operatorname{arsinh}M\bigr)
\]
for $\abs y\le\tfrac12$, and $\abs{\widehat K(\xi)}\ll e^{-\abs\xi/2}\min\set{1,(J/(W\abs\xi))^J}$. The function $K$
is even, nonnegative on $\R$, admissible in \eqref{rec:eq:gwstrip}, and $\widehat K(0)=1$. The same
bounds hold for $K^{(j)}$, $0\le j\le J-4$, with constants depending on $j,W,J$.
\end{lemma}

\begin{proof}
The constants $c_M$ are bounded above and below, by comparison with $\operatorname{sech}^2x$ and by integration
over $[-1,1]$ using $(1+y/M)^M\le e^y$. In the closed strip $w=\cosh(\pi z)$ has $\Re w\ge0$ and
$\abs w^2=\sinh^2(\pi x)+\cos^2(\pi y)$, so $\abs{1+w/M}\ge1$, and $\abs{1+w/M}\ge\sqrt2$ when $\abs w\ge M$,
which holds for $\abs x\ge\pi^{-1}\operatorname{arsinh}M$; and $\abs{\operatorname{sech}^2(x+iy)}\ll e^{-2\abs x}$. The nearest
poles of the binomial factor have $\abs{\Im z}=1$ (if $\cosh\pi z=-M$ then $\sin\pi y=0$), so the
kernel is analytic with exponential decay on $\abs{\Im z}\le\tfrac34$, and shifting the Fourier
integral to $\Im z=-\operatorname{sgn}(\xi)/2$ gives $\abs{\widehat k_M(\xi)}\ll e^{-\abs\xi/2}$. Averaging translates
over the compact support $[-W,W]$ gives the pointwise bounds, and the multiplier
$(\sin(W\xi/J)/(W\xi/J))^J$ the last one. For $K^{(j)}$ put one derivative on each of $j$ uniform
factors; each is a signed measure on two endpoints.
\end{proof}

At centre $U$, $T/2\le U\le2T$, take $M=\ceil{10\log\log T}$ and $H=W+\pi^{-1}\operatorname{arsinh}M=\pi^{-1}\log_3T+O_W(1)$.
Multiplying $K(z-U)$ by a fixed finite sum of exponentials $e^{i\lambda z}$ costs a fixed factor on
the strip, and by the lemma and the unit-interval count the zeros with $\abs{\Im\rho-U}>H$
contribute $O_W(2^{-M/2}\log(2T))=o(1)$; the pole terms are $O_W(e^{-2U})$; and Stirling's formula on
$\abs{t-U}\le U/2$ gives
\begin{equation}\label{rec:eq:binomgamma}
 \int h(t)G_\chi(t)\dd t=\log\bigl(qU/(2\pi)\bigr)\int h(t)\dd t+O_W(T^{-1})
\end{equation}
for these tests. For $t\mapsto K(t-U)e^{i\omega t}$ the transform at $\xi$ is $e^{iU(\omega-\xi)}\widehat K(\xi-\omega)$,
so, for any fixed finite list of frequencies, the prime sum in \eqref{rec:eq:gwstrip} tends, as
$W\to\infty$, to its terms with $\pm\log n$ in the list, uniformly in $U$ and $M$; the large-$n$
majorant $C_\omega\log n/(n(1+\log n)^J)$ is summable. The primes are fixed first, then $W$, then $T$.

\begin{theorem}[Local independence of multiplicity vectors; \stU]\label{rec:thm:localrank}
There are $C,T_0$ such that for $T\ge T_0$ the vectors $m(\rho)$, $\abs{\Im\rho-T}\le\pi^{-1}\log_3T+C$,
indexed by distinct complex zeros, span $\C^d$. In particular every nonzero complex linear
combination of the $L'/L(s,\chi_j)$ has a pole in this height range. For a family closed under
conjugation the same holds for $T\le-T_0$.
\end{theorem}

\begin{proof}
The induced characters on $(\Z/Q\Z)^\times$ are distinct, so the value table has rank $d$; choose $d$
independent rows and a prime $p_i$ in each corresponding class, so that
$B_{ij}=\chi_j(p_i)\log(p_i)/\sqrt{p_i}$ is invertible. For $\lambda=\log p_i$ use
$h_i(t)=K(t-T)\bigl(e^{i\lambda t}-e^{i\lambda T}\widehat K(\lambda)\bigr)$, whose integral vanishes by evenness of $K$,
so that \eqref{rec:eq:binomgamma} is $O_W(1/T)$ despite the growing common gamma term. Frequency
selection gives $2\pi\sum_{\abs{\Im\rho-T}\le H}h_i(z_\rho)m_j(\rho)=-B_{ij}+o_W(1)+o_{T;W}(1)$, the unmodulated
correction having a bounded prime sum and coefficient $o_W(1)$, and the tail including zeros off
the line. Choose $W$, then $T_0$, to make the error smaller than half the least singular value of
$B$; the left matrix has rank $d$ and its rows lie in the span of the $m(\rho)$. At a zero the residue
of $\sum_jv_jL'/L(s,\chi_j)$ is $v\cdot m(\rho)$. Conjugation gives negative heights.
\end{proof}

This sharpens \cref{zer:thm:rank}, whose windows have length $K_q\log T$, and the radii
$b\log_3T+O_{q,b}(1)$, $b>1/\pi$, obtained earlier for the characters of one modulus. The
coefficient $1/\pi$ is reached with an additive constant; its optimality is open.

\begin{lemma}[Distinguishing primes; \stU]\label{rec:lem:frame}
For every $D_0>0$ there is a finite set $P$ of primes coprime to $Q$ with
$\sum_{p\in P}\frac{(\log p)^2}p\bigl\lvert\sum_ju_j\chi_j(p)\bigr\rvert^2>D_0^2\norm u_2^2$ for all nonzero $u\in\C^d$.
\end{lemma}

\begin{proof}
In each reduced class modulo $Q$ the sum of $1/p$ diverges; choose primes so that
\[
 \sum(\log p)^2/p>D_0^2/\varphi(Q)
\]
in every class, and use orthogonality.
\end{proof}

\begin{theorem}[Both signs at the endpoint; \stG]\label{rec:thm:bothsigns}
Let $d\ge2$. There are $C,c>0$ and $T_0$ such that for $T\ge T_0$, every real unit vector
$u\perp\mathbf 1$ and $s\in\set{\pm1}$,
\[
 \sum_{\abs{\gamma-T}\le\pi^{-1}\log_3T+C}\bigl(s\,u\cdot m(\gamma)\bigr)_+\ge c,
\]
and for each fixed $b>1/\pi$ and large $T$,
$\sum_{\abs{\gamma-T}\le b\log_3T}(s\,u\cdot m(\gamma))_+\gg_b\log_3T$. In particular, for every nonzero $n\in\Z^d$
with $\sum_jn_j=0$, $\prod_jL(s,\chi_j)^{n_j}$ has a zero and a pole in every window of the first
kind, after all cancellations, and $\gg\norm n_2\log_3T$ of each in the windows of the second kind.
\end{theorem}

\begin{proof}
Put $c_u=\sum_ju_j\log q_j$ and $b_p(u)=\frac{\log p}{\sqrt p}\sum_ju_j\chi_j(p)$, and fix $P$ by
\cref{rec:lem:frame} with $\norm{b(u)}_2>\norm{(\log q_j)_j}_2+3$ for every such $u$. For $a=(a_0,(a_p)_{p\in P})$
of norm one let $A_a(t)=a_0+\sum_pa_pe^{it\log p}$ and $h_{a,U}(t)=K(t-U)\abs{A_a(t)}^2$, extended as
$A_a(z)\overline{A_a(\bar z)}$. The frequencies of $\abs{A_a}^2$ are $0$ (coefficient one), $\pm\log p$ (from the
constant term) and $\log(p/r)$ for distinct $p,r\in P$, which are not signed prime-power logarithms;
so $\int h_{a,U}=1+o_W(1)$. Combine \eqref{rec:eq:gwstrip} over $j$ with coefficients $u_j$: by
\eqref{rec:eq:binomgamma} the gamma contribution is $c_u\int h_{a,U}+O_W(T^{-1})$, the common term
$\log(U/2\pi)$ cancelling because $\sum u_j=0$, and
\[
 2\pi\sum_{\abs{\gamma-U}\le H}h_{a,U}(\gamma)\,u\cdot m(\gamma)=c_u-2\Re\Bigl(\overline{a_0}\sum_pa_pb_p(u)\Bigr)+o_W(1)+o_{T;W}(1),
\]
uniformly in $T/2\le U\le2T$, $a$ and $u$. By \cref{src:thm:dgrh} each $h_{a,U}(\gamma)$ is nonnegative.
Take $a_0=1/\sqrt2$ and $a_p=\pm\overline{b_p(u)}/(\sqrt2\norm{b(u)}_2)$: the main terms are $c_u\mp\norm{b(u)}_2$,
of either sign with margin greater than three. Since $\abs{A_a}^2\le1+\abs P$ and $K\ge0$,
\begin{equation}\label{rec:eq:weightedsign}
 \sum_{\abs{\gamma-U}\le H}K(\gamma-U)\bigl(s\,u\cdot m(\gamma)\bigr)_+\ge c_0=\frac1{2\pi(1+\abs P)}\qquad(s=\pm1),
\end{equation}
and $U=T$ with $\norm K_\infty<\infty$ gives the first claim. For the second, with $R=b\log_3T$ and
$M,W,H$ fixed at their values for $T$, integrate \eqref{rec:eq:weightedsign} over $\abs{U-T}\le R-H$:
each ordinate receives total weight at most $\int K=1$, so the sum is at least $2c_0(R-H)$.
\end{proof}

\Cref{zer:thm:short} gives, for the same unshifted families, mass $\gg H\log\log H$ on intervals of
length $H\ge(2/\pi+\eps)\log_3T$; the theorem above removes the $\eps$ at the price of a bounded
mass, and its tests survive translation of the functions (\cref{rec:thm:moving}).

\begin{corollary}[Exact balancing; \stG]\label{rec:cor:balance}
In each window of the first kind in \cref{rec:thm:bothsigns} the points
$v_\gamma=m(\gamma)/\sum_jm_j(\gamma)$ have a convex hull in the hyperplane $\sum_jx_j=1$ containing the
ball of radius $c_1/\log T$ about $\mathbf 1/d$, and some at most $d$ of these ordinates with
positive integer weights have a weighted sum of multiplicity vectors with all coordinates equal.
\end{corollary}

\begin{proof}
The weight $\sum_{\abs{\gamma-T}\le H}K(\gamma-T)\sum_jm_j(\gamma)$ is $O(\log T)$, so
\eqref{rec:eq:weightedsign} gives $\max_\gamma u\cdot v_\gamma\gg1/\log T$ for every unit $u\perp\mathbf 1$;
separation gives the ball, Carath\'eodory's theorem a rational convex combination equal to
$\mathbf 1/d$ on at most $d$ points, and clearing denominators the integer weights.
\end{proof}

\begin{corollary}[A cost for moving or deleting zeros; \stG]\label{rec:cor:binomtransport}
For two fixed distinct primitive characters and $b>1/\pi$, charge $1$ for each unmatched
occurrence and $\min(1,\abs{\gamma-\gamma'})$ for each matched pair in a partial matching of their
zero lists in $[T-b\log_3T,T+b\log_3T]$. Every matching costs $\gg_b\log_3T$.
\end{corollary}

\begin{proof}
Take $u=(1,-1)/\sqrt2$ and the positive test above; by the tail bound the indicator of $[-H,H]$
may be replaced by a taper $\eta_H$, equal to one on $[-H,H]$ and linear to zero on
$H\le\abs x\le H+1$. For $J'=[T-R+H+1,T+R-H-1]$, $R=b\log_3T$, the function
$f(t)=\abs{A_a(t)}^2\int_{J'}K(t-U)\eta_H(t-U)\dd U$ is nonnegative, supported in the window, and has
supremum and Lipschitz constant bounded independently of $T$. Integrating the signed estimate over
$J'$ gives $\sum f(\gamma_1)-\sum f(\gamma_2)\gg\abs{J'}\gg\log_3T$, while a matched pair contributes at most
$\max(2\norm f_\infty,\operatorname{Lip}f)\min(1,\abs{\gamma_1-\gamma_2})$ and an unmatched zero at most $\norm f_\infty$.
\end{proof}

\paragraph{Translated functions and coalescing shifts.}
Consider $L(s+i\tau_j,\chi_j)$ with real shifts $\tau=(\tau_j)$ in a fixed compact set $E$, allowing
repeated primitive characters but with $(\chi_j,\tau_j)\ne(\chi_k,\tau_k)$ on $E$ for $j\ne k$; the
multiplicity at $\tfrac12+it$ is $m_{\chi_j}(t+\tau_j)$.

\begin{theorem}[Uniform translated families; \stG]\label{rec:thm:moving}
\Cref{rec:thm:localrank}, \cref{rec:thm:bothsigns} and \cref{rec:cor:balance} hold for these shifted
functions uniformly in $\tau\in E$, the rank statement without \cref{src:thm:dgrh}; for a pair, so does \cref{rec:cor:binomtransport}, and its
lower bound persists after the infimum over all allowed translations.
\end{theorem}

\begin{proof}
Testing the original list with $h(t-\tau_j)$ replaces the prime coefficient by $\chi_j(n)n^{-i\tau_j}$
and the gamma density by $G_{\chi_j}(t+\tau_j)$; poles and errors are uniform for bounded shifts. It
suffices to choose the primes of \cref{rec:lem:frame} uniformly. For fixed $\tau$ the matrix
$G_X(j,k)=\sum_{p\le X,p\nmid Q}\frac{(\log p)^2}p\chi_j(p)\overline{\chi_k(p)}p^{-i(\tau_j-\tau_k)}$ satisfies
$G_X/(\tfrac12\log^2X)\to I_d$, by the prime number theorem in progressions and partial summation
($\int_2^X(\log x)x^{-1-iv}\dd x$ is $\tfrac12\log^2X+O(1)$ for $v=0$ and $O_v(\log X)$ otherwise); where the
characters coincide the shifts differ. The condition $\lambda_{\min}(G_X)>D_0^2$ is open in $\tau$ and
persists as $X$ grows, since $G_X$ increases as a quadratic form, and compactness gives one $X$ for
all of $E$. All later arguments use finitely many frequencies with uniformly bounded coefficients.
\end{proof}

For a primitive $\chi$, $r\ge1$ and real $h\ne0$ put
\[
 F_{r,h}(s)=\prod_{k=0}^rL(s+ikh,\chi)^{(-1)^{r-k}\binom rk},
\]
so that $\zeta(s+ih)/\zeta(s)$ and $\zeta(s+2ih)\zeta(s)/\zeta(s+ih)^2$ are the cases $r=1,2$; at $h=0$ all
factors cancel. Zeros and poles are counted after cancellation.

\begin{theorem}[Uniformity as the shifts coalesce; \stG]\label{rec:thm:coalesce}
Fix $\chi$, $r\ge1$ and $B_0>0$. There are $C,T_0$ independent of $h$ such that for $T\ge T_0$ and
every $0<\abs h\le B_0$, $F_{r,h}$ has a zero and a pole on the critical line in
$\abs{\Im s-T}\le\pi^{-1}\log_3T+C$; if $N_\pm$ count them there, $N_\pm\ge c\abs h^r$, and in
$\abs{\Im s-T}\le b\log_3T$, $b>1/\pi$ fixed, $N_\pm\gg_{\chi,r,B_0,b}\max\set{1,\abs h^r\log_3T}$.
\end{theorem}

\begin{proof}
For a test $f$ put $\mathcal D_{r,h}(f)=h^{-r}\sum_kd_k\sum_\gamma f(\gamma-kh)$, $d_k=(-1)^{r-k}\binom rk$, zeros with
multiplicity; its signed measure has atom $\ord_{1/2+it}F_{r,h}$ at $t$, and
$h^{-r}\sum_kd_kg(x+kh)=\int_{[0,1]^r}g^{(r)}(x+h(v_1+\dots+v_r))\dd v$ controls it although $\abs{d_k/h^r}$ is large.
Choose distinct primes $p_1,p_2\nmid q_\chi$ and put $b_p(h)=\chi(p)\frac{\log p}{\sqrt p}\bigl(\frac{e^{-ih\log p}-1}h\bigr)^r$,
extended continuously to $h=0$. They have no common zero on $[-B_0,B_0]$ (that would make
$\log p_1/\log p_2$ rational), so $\max_{p=p_1,p_2}\abs{b_p(h)}\ge\delta>0$. Use $J=r+4$ in
\cref{rec:lem:binomkernel}, and for the maximizing $p$ the two positive tests
$f=K(t-U)\bigl(1\pm\Re(\alpha e^{it\log p})\bigr)$, $\alpha=\overline{b_p(h)}/\abs{b_p(h)}$. In
$\mathcal D_{r,h}(f)$ the coefficient $\chi(n)$ is multiplied by $((n^{-ih}-1)/h)^r$, of size at most
$(\log n)^r$, the off-frequency majorant $O_p((\log n)^{r+1}/(n(1+\log n)^J))$ is summable, and the
peak gives $\mp\abs{b_p(h)}+o_W(1)$ uniformly in $h$. The gamma density
$h^{-r}\sum_kd_kG_\chi(t+kh)$ is $O(T^{-r})$ near $U$ and bounded globally, since
$G_\chi^{(r)}(t)=O((1+\abs t)^{-r})$; the pole terms are $O(e^{-2U})$. For the truncation take a smooth
taper $\eta_H$ with derivatives through order $r$ bounded independently of $H$ and
$g=(1-\eta_H(t-U))f$: every derivative through order $r$ is $O(2^{-M/2}e^{-2\abs{t-U}})$, and applying the
integral identity before summing over zeros gives $\abs{\mathcal D_{r,h}(g)}=o(1)$ with no factor
$\abs h^{-r}$. Hence the tapered nonnegative tests have readings of opposite signs bounded below,
and since $0\le1\pm\Re(\cdot)\le2$, $\sum_{\abs{t-U}\le H+1}K(t-U)(s\ord F_{r,h})_+\gg\abs h^r$ for $s=\pm1$. The
supremum of $K$ gives both signs in the endpoint window, integration over $U$ the larger window,
and each positive count is an integer.
\end{proof}

\begin{theorem}[Microscopic steps; \stG]\label{rec:thm:micro}
Fix $r\ge1$, $\chi$ and $0<\kappa<2\pi$. For any $C>\frac\pi2+\frac{\pi\kappa}{2\pi-\kappa}$ and all large $T$, $F_{r,h}$ has
a zero and a pole in $\abs{\Im s-T}\le C/\log\log T$ simultaneously for all $0<\abs h\le\kappa/\log T$;
in particular $[T,T+1]$ contains $\gg_{\kappa,\chi,r}\log\log T$ zeros and as many poles, uniformly in these
steps.
\end{theorem}

\begin{proof}
Put $\lambda=\log T$, $\ell=\log\log T$. Under \cref{src:thm:dgrh}, $\abs{S_\chi(t)}\le(\tfrac14+o(1))\lambda/\ell$ for
$\abs{t-T}=O(1)$ \cite{CarneiroChandeeMilinovich2015Note}, so for fixed $a>0$ and $\abs{x-T}\le2$,
$\bigl\lvert N_\chi((x,x+a/\ell))-\frac a{2\pi}\frac\lambda\ell\bigr\rvert\le(\tfrac12+o(1))\frac\lambda\ell$. A central interval of length
$a/\ell$, $a>\pi$, contains a zero $\gamma_0$. For $A>\pi\kappa/(2\pi-\kappa)$, each of the two adjacent intervals
of length $A/\ell$ from $\gamma_0$ has at least $A\lambda/(\kappa\ell)-O(1)$ sites of the grid $\gamma_0+nh$ and at most
$(A/(2\pi)+\tfrac12+o(1))\lambda/\ell$ zeros, and the first coefficient is larger; so each side has at least
$r$ empty sites. With $v(n)=m_\chi(\gamma_0+nh)$, $v(0)>0$, the order-$r$ divided difference on $r$
negative empty indices and $0$ is positive by the Lagrange formula, and on $r-1$ negative empty
indices, $0$ and one positive empty index it is negative. Every divided difference on integer
nodes is a convex combination of divided differences on $r+1$ consecutive sites in their span
(interpolate through a missing interior integer $y$ with weights $(y-x_0)/(x_r-x_0)$ and
$(x_r-y)/(x_r-x_0)$ and induct), and on consecutive sites it is $\Delta^rv(n)/r!$, the order of
$F_{r,h}$ at the corresponding ordinate. All these ordinates lie within $(a/2+A)/\ell$ of $T$, and
$a/2+A<C$ is possible. For the unit interval use disjoint windows with a slightly larger $\kappa'<2\pi$.
\end{proof}

At spacing exactly $2\pi/\log T$ a lattice is invariant under the step, so a density comparison
alone gives nothing at $\kappa=2\pi$.

\begin{corollary}[Nonintegral multiples of the mean spacing; \stG]\label{rec:cor:nonres}
Fix $K_0,\delta>0$. If $\nu=\abs h\log T/(2\pi)$ satisfies $0<\nu\le K_0$ and $\operatorname{dist}(\nu,\set{1,2,\dots})\ge\delta$,
then $F_{r,h}$ has a zero and a pole in $\abs{\Im s-T}\le C_{\chi,r,K_0,\delta}/\log\log T$, uniformly in $h$.
\end{corollary}

\begin{proof}
For $\nu\le\tfrac12$ use \cref{rec:thm:micro}. Otherwise take consecutive half-open bins of length
$s=\abs h$, $r+1$ blocks of $k=\floor{A/(s\ell)}$ bins from $T$, and bin counts $a_n$. The counting formula
gives block means $\nu+O(1/A)+o(1)$, the endpoint error $O(\log T/\ell)$ being divided by
$k\asymp A\log T/(\nu\ell)$; for $A$ large each block mean lies strictly between $\floor\nu$ and $\ceil\nu$, so
each block has counts below and above $\nu$. Choosing one bin per block with $a_n-\nu$ of alternating
signs makes every term of the order-$r$ Lagrange divided difference of one strict sign, and the
reverse choice the other. The consecutive-site identity gives $\Delta^ra_n$ of both signs, and for
$h=s>0$, $\Delta^ra_n$ is the signed divisor count of $F_{r,h}$ in bin $n$, translation by $ks$ preserving the
bins; for $h=-s$ reverse the grid.
\end{proof}

The archive checks the binomial identity on $456$ node choices, translated bin identities, and,
for $\chi=1$, $r=2$, $\abs h\le1$, the two-prime lower bound $\sum_{p=2,3}\abs{b_p(h)}^2>\tfrac12$
and, for $\zeta$ and $\chi_{-4},\chi_8,\chi_{-8}$, that the twelve primes $3,\dots,47$ of \cref{rec:lem:frame} already
exceed the conductor bias $\sqrt6\log2$ in every balanced direction:
\begin{center}
\artifact{translated_dirichlet_zero_divisors/checks/translation_check.txt}\\
\artifact{translated_dirichlet_zero_divisors/checks/lattice_check.txt}
\end{center}

\paragraph{The zeros of a differentiated completion.}
For primitive $\chi$ of conductor $q$ and parity $a$ let $A(s)=(q/\pi)^{(s+a)/2}\Gamma((s+a)/2)$, with
the factor $s(s-1)/2$ included for $q=1$, and $\Xi_\chi=AL(s,\chi)$. Let $P(X)=\prod_{j=1}^d(X-i\alpha_j)$,
$\alpha_j\in\R$, be monic ($P=1$ allowed), $F=P(\dd/\dd s)\Xi_\chi$ and, for a prime $p$, $h=\log p$ and
\[
 S_p(T,H)=\frac{\sqrt{2\pi p}}H\sum_\gamma\exp\Bigl(-\frac{(\gamma-T)^2}{2H^2}+ih\gamma\Bigr),\]\[
 z=\tfrac12\log\frac T{2\pi}+\frac{i\pi}4,\quad c=\tfrac12\log q,\quad Q_P(X)=P(X+c),
\]
the sum over the zeros $\tfrac12+i\gamma$ of $F$ with multiplicity.

\begin{lemma}[Zeros after differentiation; \stG]\label{rec:lem:interlace}
$F$ has order at most one, all its zeros lie on $\Re s=\tfrac12$, a unit interval contains
$O_{\chi,P}(\log(2+\abs u))$ of their ordinates, and $F(s)=(-1)^d\epsilon_\chi\overline{F(1-\bar s)}$.
\end{lemma}

\begin{proof}
Dividing $\Xi_\chi(\tfrac12+iw)$ by a square root of $\epsilon_\chi$ gives a real entire $f$ of order at most
one with only real zeros, by \cref{src:thm:dgrh}. Its genus-one product gives
$\Im(f'/f)(w)=-\Im(w)\sum_tm_t/\abs{w-t}^2$, so $(\dd/\dd w+\alpha)f$ has only real zeros for real $\alpha$: a zero of
order $m$ keeps order $m-1$, and $f'/f$ decreasing from $+\infty$ to $-\infty$ on each gap gives one more
simple zero. Interval counts change by at most one per step, and $D-i\alpha$ corresponds to
$-i(\dd/\dd w+\alpha)$; induct. The last identity follows by conjugation.
\end{proof}

\begin{theorem}[Prime-frequency formula; \stG]\label{rec:thm:primefreq}
For fixed $\chi,P,p$ there is $\eta_p>0$ such that uniformly for $1\le H\le T^{2/3}$,
\[
 S_p(T,H)=-h\chi(p)\frac{Q_P(z-h)}{Q_P(z)}+O_{\chi,P,p}\bigl(T^{-1/5}+(\log T)e^{-\eta_pH^2}\bigr),
\]
including $p\mid q$. With $H=\log\log T$ the error is $o((\log T)^{-k})$ for every $k$.
\end{theorem}

\begin{proof}
In the Banach algebra of Dirichlet coefficient sequences with norm $\sum\abs{u_n}n^{-2}$, put $g=z+c$,
$U_g(s)=\sum_n\chi(n)P(g-\log n)n^{-s}$ and $V_g(s)=-\sum_n\chi(n)(\log n)P(g-\log n)n^{-s}$. Since
$U_g/g^d\to L(s,\chi)$ in norm with error $O(1/\abs g)$ and $L^{-1}$ lies in the algebra, a Neumann series
gives $\norm{U_g^{-1}}=O(\abs g^{-d})$, so $B_g=V_g/U_g=\sum b_n(g)n^{-s}$ has bounded norm, $b_1=0$, and
comparison of coefficients at $p$ in $B_gU_g=V_g$ gives $b_p(g)=-(\log p)\chi(p)P(g-\log p)/P(g)$. For
$s=2+it$, $\abs{t-T}\le T^{3/4}$, Stirling gives $A'/A=g+O(T^{-1/4})$ and $(A'/A)^{(j)}=O(T^{-j})$, hence
$A^{(k)}/A=g^k+O(T^{-1/4}(1+\abs g)^k)$; replacing $\dd/\dd s$ acting on $A$ by $g$ in $F/A$ and $F'/A$ gives $U_g$ and
$gU_g+V_g$ up to $O(T^{-1/4}(1+\abs g)^{d+1})$, and division yields $(F'/F)(2+it)=g+B_g(2+it)+O(T^{-1/5})$.
Apply the residue theorem on $\sigma=2,-1$ to $(F'/F)(s)p^s\exp((s-\tfrac12-iT)^2/(2H^2))$, with horizontal
edges at heights $\gg1/\log(2+\abs t)$ from all ordinates (\cref{rec:lem:interlace}), where
$F'/F=O(\log^2(2+\abs t))$; reflection gives $(F'/F)(-1+it)=-\overline{(F'/F)(2+it)}$. A term of frequency $v$ on
the edge $\sigma$ integrates to $e^{ivT}e^{-H^2v^2/2-(\sigma-1/2)v}$; the only resonance is $n=p$ on the right
edge, and all other frequencies have $\abs v\ge\log(1+1/p)=\delta$, costing $O((\log T)e^{-\delta^2H^2/4})$.
\end{proof}

\begin{theorem}[One-prime reconstruction; \stG]\label{rec:thm:operatorinverse}
For a prime $p\nmid q$, the values $S_p(T,\log\log T)$ along any unbounded sequence of $T$ determine
$\chi(p)$ and the monic $P$, and, if $\deg P>0$, the conductor $q$. Writing the normalized expansion as
$1+\sum_jr_jz^{-j}$, $r_1=-dh$, $r_2=\binom d2h^2+ha_1$ with $Q_P=X^d+a_1X^{d-1}+\dots$, and
\[
 q=\exp\Bigl(\frac2{dh}\Re\Bigl[r_2-\binom d2h^2\Bigr]\Bigr).
\]
Terms through $z^{-d-1}$ determine $P$, and for each $d\ge1$ terms through $z^{-d}$ can agree for distinct
operators of degree $d$, even with $q$ and $\chi$ fixed.
\end{theorem}

\begin{proof}
Divide by the limit $-h\chi(p)\ne0$ and read the coefficients successively. The roots of $Q_P$ are
$-c+i\alpha_j$, so $\Re a_1=dc$. If monic $Q_1,Q_2$ of degree $d$ first differ by $\Delta a_j$ in position $j$, then
$\frac{Q_1(z-h)}{Q_1(z)}-\frac{Q_2(z-h)}{Q_2(z)}=jh\Delta a_jz^{-j-1}+O(z^{-j-2})$, by taking logarithms near infinity, which
recovers $a_j$ from $r_{j+1}$ inductively. For sharpness let $f(x)=\prod_{j=1}^d(x-j)$; for small real $t\ne0$,
$f+t$ has $d$ simple real roots, and $P_1(X)=i^df(X/i)$, $P_2(X)=i^d(f(X/i)+t)$ differ by a constant,
so they agree through $z^{-d}$ and differ at $z^{-d-1}$.
\end{proof}

For the quadratic character modulo $5$, $p=2$ and $P=X^2-iX+2$, the limit is $\log2$, the first
correction gives degree two, $r_2$ gives $q=5$, and the third correction gives $a_2$. For $P=X$ and
$q=1$ the direct sums at $p=3$ over $626$ and $1572$ zeros of $\xi'$ agree with $-\log3(1-\log3/z)$ to
about $10^{-4}$:
\begin{center}
\artifact{dirichlet_prime_frequency_inverse/check_algebra.py}\\
\artifact{dirichlet_prime_frequency_inverse/zero_example_output.txt}
\end{center}

\begin{corollary}[Local distinction of operators; \stG]\label{rec:cor:operatorlocal}
For each $k$ there are $A,B$ such that $S_p$ with $H=A\sqrt{\log\log T}$, restricted to
$\abs{\gamma-T}\le B\log\log T$, equals the main term up to $o((\log T)^{-k})$. Two distinct pairs
$(\chi_1,P_1)\ne(\chi_2,P_2)$ have different zero multisets in $[T-C\log\log T,T+C\log\log T]$ for all
large $T$, and $[X,2X]$ contains $\gg X/\log\log X$ unmatched occurrences; constants depend on the pairs.
\end{corollary}

\begin{proof}
The weights $e^{-(\gamma-T)^2/(4H^2)}$ have sum $O(H\log T)$, so the omitted tail is
\[
 O\bigl((\log T)^{1-B^2/(4A^2)}\bigr);
\]
take $\eta_pA^2$ and $B^2/(4A^2)$ larger than $k+2$. Distinct characters
differ at some prime, where the constant terms differ; for equal characters distinct monic
operators differ at a finite order by \cref{rec:thm:operatorinverse}. Pack disjoint windows in
$[X,2X]$.
\end{proof}

If $d$ is known, interpolating the truncated sums at the $d+2$ centres $T^j$, $1\le j\le d+2$, through
degree $d+1$ in $1/z$ (the nodes $\log T/z(T^j)$ tend to $2/j$, so the Vandermonde inverses are bounded)
estimates the coefficient of $z^{-k}$ with error $O((\log T)^{k-d-2})$, and for $d>0$ the conductor
estimate eventually rounds to $q$; no numerical threshold is claimed.

For $P=1$ the window here is much longer than that of \cref{rec:thm:localid}; the content of the
corollary is the operator. Without \cref{src:thm:dgrh} the same holds for the sum
$\frac{\sqrt{2\pi}}H\sum_{F(\rho)=0}p^\rho\exp((\rho-\tfrac12-iT)^2/(2H^2))$ over complex zeros (\stU): the zeros of
$\Xi_\chi$ lie in $0\le\Re s\le1$, for a real entire $f$ of order at most one with zeros in
$\abs{\Im w}\le\tfrac12$ the paired product gives $\Im(f'/f)(w)=-\sum_v(\Im w-\Im v)/\abs{w-v}^2$, of strict sign
outside the strip, a zero-free intermediate function is excluded by the right-edge lower bound
$\abs{P(D)\Xi_\chi(2+it)/A(2+it)}\gg(\log\abs t)^d$, and Jensen's formula on disks of radius $5$ about
$2+iu$ gives the unit-interval count; the rest of the proof is unchanged.

\subsection{Signed multiplicity differences in short intervals and in phase boxes}\label{zer:sec:short}
{\emergencystretch=6em\sources{grh-methods-polynomial-density; grh_cm_four_terms_and_prime_bias; grh_spectral_rigidity; grh_filtered_prime_races; grh_prime_error_balancing}}

The next two theorems locate both signs of differences of zero multiplicities on the
critical line: in every interval of triple-logarithmic length for an arbitrary fixed finite
family, and in every prescribed box of prime phases for the twist-family difference of the
source proof. A third statement records that the multiplicity vectors of all characters to a
fixed modulus have full rank in windows of logarithmic length, and a fourth that no nonzero
polynomially weighted combination of the multiplicity divisors of distinct primitive
characters has finite support.

\paragraph{Intervals of triple-logarithmic length.}
Let $\chi_1,\dots,\chi_r$ be distinct primitive Dirichlet characters, of conductors
$q_1,\dots,q_r$, and let $m_j(t)$ be the order of $L(s,\chi_j)$ at $s=\tfrac12+it$. The
conductor-one character gives the Riemann zeta function. Fix $r\ge2$. For $c\in\R^r$ put
$z_c(t)=\sum_jc_jm_j(t)$ and $z_{c,\pm}=\max(\pm z_c,0)$. Write
$\log_3T=\log\log\log T$. Sums below count each ordinate once.

\begin{theorem}[\stG]\label{zer:thm:short}
Fix $\eps>0$ and the characters. For all sufficiently large $T$, every
\[
 (2/\pi+\eps)\log_3T\le H\le T
\]
satisfies, simultaneously for both signs and every $c$ with $\sum_jc_j=0$ and
$\norm c_2=1$,
\[
 \sum_{T<t\le T+H}z_{c,\pm}(t)\ \gg_{\eps,\chi_1,\dots,\chi_r}\ H\log\log H .
\]
\end{theorem}

For two characters take $c=(1,-1)/\sqrt2$: each such interval contains both an ordinate
where the first multiplicity is greater and one where the second is greater. If the
multiplicities at a common zero are $2$ and $1$, the sum counts one unit there. Srinivas
proved that for distinct functions in the Selberg class of equal degree each direction of
excess occurs within distance $O(\log\log T)$ of every large $T$ \cite{Srinivas2002Distinct};
the use of real critical-line ordinates gives a smaller interval for fixed Dirichlet
families and a lower bound for the mass in it. Banks gives bounds uniform in the character
modulus \cite{Banks2023Comparing}; Fujii and Bombieri--Perelli obtain much larger total
multiplicity counts in global ranges \cite{Fujii1976Zeros,BombieriPerelli1998Distinct}, and
Banks and Loftus detect genuinely unshared zeros in longer windows
\cite{BanksLoftus2026Unshared}.

\begin{proof}
\emph{A positive weight and its prime frequencies.} Let $Q$ be the least common multiple of
the conductors and $C(n)=\sum_jc_j\chi_j(n)$. The induced characters of $(\Z/Q\Z)^\times$
are distinct by uniqueness of the primitive inducer, and character orthogonality gives the
mean square of $C(v)$ over the group as $1$. Choose $v$ with $\abs{C(v)}\ge1$ and $\theta$
with $e^{i\theta}C(v)=\abs{C(v)}$. Put $Y=\lambda\log H$ for a sufficiently small fixed
$\lambda>0$. Take all primes $p\le Y$ congruent to $v$ modulo $Q$ and write
\[
 w_p=\frac{\log p}{\sqrt p},\qquad V^2=\sum_pw_p^2\sim\frac{(\log Y)^2}{2\varphi(Q)},
\]
the asymptotic following by partial summation from the prime number theorem in the finitely
many fixed progressions. Form the real trigonometric sums
\[
 A(t)=\Im\prod_p\bigl(1+i(w_p/V)\cos(t\log p+\theta)\bigr),\qquad
 P_\pm(t)=1\pm e^{-1/2}A(t).
\]
On the real axis $\abs A\le\prod_p\exp(w_p^2/(2V^2))=\sqrt e$, so $0\le P_\pm\le2$. Extend
their finite exponential expansions to complex $t$, and write either of the two as
$P(t)=\sum_Rb_Re^{it\log R}$, where the positive rationals $R$ are in lowest terms. Unique
factorization gives $b_1=1$ and $b_p=\pm e^{-1/2}e^{i\theta}w_p/(2V)$ at the chosen primes.
Every other nonconstant frequency, apart from their reciprocals, has at least two different
primes with nonzero exponents and cannot be a prime power or its reciprocal. Also
$b_{1/R}=\overline{b_R}$. Both the numerator and the denominator of $R$ are at most
$D=\prod_{p\le Y}p\le e^{2Y}$ for large $Y$. Expanding the product and using
Cauchy--Schwarz gives $\sum_R\abs{b_R}\le1+\prod_p(1+w_p/V)\le1+\exp(\sqrt{\pi(Y)})$.

\emph{An explicit-formula test localized in height.} Choose $0<a<\pi$ and $0<\kappa<1/2$
such that $a(1/2-\kappa)(2/\pi+\eps)>1$. Set $B=\kappa H$ and $X=T+H/2$. Let $k(t)$ be
$e^{-\cosh(at)}$ divided by its real integral, and put
\[
 K_B(t)=\frac1{2B}\int_{-B}^Bk(t-y)\,\dd y,\qquad h(t)=P(t)K_B(t-X).
\]
Thus $0\le h(t)\le1/B$ on $\R$. Use the convention
$\widehat h(u)=\int_\R h(t)e^{-itu}\,\dd t$. For fixed $b$ with $1/2<b<\pi/(2a)$, shifting
the contour to $\Im t=\pm b$ gives $\abs{\widehat k(u)}\ll_{a,b}e^{-b\abs u}$, since
$\abs{e^{-\cosh(a(t+ib))}}=e^{-\cos(ab)\cosh(at)}$. Consequently
\[
 \widehat h(u)=\sum_Rb_Re^{-iX(u-\log R)}\,\widehat k(u-\log R)\,
 \zersinc\bigl(B(u-\log R)\bigr),
\]
with value $1$ of the last two factors at $u=\log R$: resonant frequencies $u=\log R$
contribute exactly $b_R$, independently of the center $X$.

For a primitive $\chi_j$ write $a_j=(1-\chi_j(-1))/2$ and $\psi=\Gamma'/\Gamma$. Under
Dirichlet GRH the Guinand--Weil formula in this normalization reads
\begin{align*}
 2\pi\sum_tm_j(t)h(t)={}&2\pi\,\mathbf 1_{q_j=1}\set{h(i/2)+h(-i/2)}\\
 &+\int_\R h(t)\Bigl[\log(q_j/\pi)+\Re\psi\bigl((1/2+a_j+it)/2\bigr)\Bigr]\dd t\\
 &-2\Re\sum_{n\ge2}\frac{\Lambda(n)\chi_j(n)}{\sqrt n}\,\widehat h(\log n).
\end{align*}
The real-valued test need not be even. The formula follows by integrating the completed
logarithmic derivative on the two sides of the critical strip, using its functional
equation on the left; the strip of decay just proved justifies both contour shifts and
absolute summation on the prime side. For the convention see
\cite[\S5]{FourierOptimizationQuadratic}.

After multiplication by $c_j$ and summation, the archimedean contribution is $O(1)$
uniformly: the common $\log(\abs t/2)$ from Stirling's formula cancels since
$\sum_jc_j=0$, the combined density is bounded near zero, and $\int\abs h\le2$. The pole
contribution is $o(1)$: on $\Im t=\pm1/2$ one has $\abs{P(t)}\le D^{1/2}\sum_R\abs{b_R}$,
whereas each translated kernel in $h(\pm i/2)$ has real displacement at least $X-B\ge T$.

On the prime side the total of the off-diagonal terms is $O(e^{CY}/B)$ for a fixed $C$. For
$n\le2D$ and $n\ne R$, the numerator of $n-R$ in the common denominator is a nonzero
integer, hence $\abs{\log n-\log R}\ge1/(2D^2)$. The sinc factor costs at most $2D^2/B$, and
$\sum_{n\le2D}(\log n)/\sqrt n\ll D^{1/2}\log(2D)$. For $n>2D$ the logarithmic difference
exceeds $\log2$ and the Fourier bound gives $O(B^{-1}D^bn^{-b})$; the resulting sum
converges since $b>1/2$. Summing these bounds with $\abs{b_R}$ proves the assertion,
uniformly in the phases and in the vector $c$. Choose $\lambda$ so that $C\lambda<1/2$. The
prime diagonal now gives
\[
 2\pi\sum_tz_c(t)P_\pm(t)K_B(t-X)=\mp e^{-1/2}\abs{C(v)}\,V+O(1).
\]
In particular the two positive test functions have signed sums of opposite signs and size
$\gg\log\log H$.

\emph{Removing the tails.} The number of zeros of any fixed primitive function in a unit
interval at height $y$, with multiplicity, is $O(\log(2+\abs y))$
\cite{BennettMartinOBryantRechnitzer2021Counting}. Put $L=H/2-B$. In the $j$th unit
interval outside either end of $(T,T+H]$, the maximal kernel value is
$O(\exp(-c_ae^{a(L+j)}))$, $j\ge0$; this bound also covers a zero exactly at the excluded
endpoint $T$. Summing the unit-interval bounds, using $H\le T$, gives a tail of size
\[
 O\bigl(\log(2+T)\exp(-c_a'e^{aL})\bigr)=o(1),
\]
because $e^{aL}\ge(\log\log T)^{1+\eta}$ for some $\eta>0$. The same estimates apply to
negative ordinates. Since $0\le h\le1/B$, each sign inside the interval has mass $\gg BV$.
\end{proof}

The constant $2/\pi$ is a sufficient limiting coefficient in this construction; its
optimality is open.

\begin{example}\label{zer:ex:seven}
For $\zeta$ and the Dirichlet beta function $L(s,\chi_{-4})$, use the unnormalized
difference of their multiplicities. The unamplified conductor term is $-\log4$. Using
$P=1\pm\cos(t\log7)$, the two limiting signed readings after multiplication by $2\pi$ are
$-\log4\mp2\log7/\sqrt7$. The absolute prime term exceeds $\log4$, so this one prime
already certifies both signs for sufficiently large windows of the stated length. The
growing family of primes in the proof supplies the factor $\log\log H$.
\end{example}

\paragraph{Prescribed phase boxes.}
Fix $q\ge1$ and an odd prime $d\nmid q$. Put $h=\varphi(q)$ and $\eta=(\cdot/d)$. For each
character $\chi\bmod q$ use its primitive inducing character $\chi^*$, of conductor
$f_\chi$; the character $\chi^*\eta$ is primitive of conductor $df_\chi$. At positive height
put
\[
 z(\gamma)=\sum_{\chi\bmod q}\Bigl(\ord_{1/2+i\gamma}L(s,\chi^*)
   -\ord_{1/2+i\gamma}L(s,\chi^*\eta)\Bigr),\qquad z_\pm=\max(\pm z,0).
\]
Omitted Euler factors have their zeros on $\Re s=0$, so the same integer is obtained with
the imprimitive functions; for $8\mid q$ and $d\equiv7\pmod8$ it is the residue numerator of
the source, studied as a whole in \cref{sec:rec}. Write $\mathbb T=\R/(2\pi\Z)$. For a
finite set $\mathcal B$ of primes put
$\theta_{\mathcal B}(t)=(t\log b\bmod2\pi)_{b\in\mathcal B}$.

\begin{theorem}[\stG]\label{zer:thm:phasebox}
Let $U$ be a nonempty open box of $\mathbb T^{\mathcal B}$, and fix $0<\alpha\le1$. There
are positive constants $c,T_0$ depending only on $q,d,\mathcal B,U,\alpha$ such that,
simultaneously for both signs, for $T\ge T_0$ and $T^\alpha\le H\le T$,
\begin{equation}\label{zer:eq:phase}
 \sum_{\substack{T<\gamma\le T+H\\\theta_{\mathcal B}(\gamma)\in U}}z_\pm(\gamma)\ \ge\
 cH\log\log T .
\end{equation}
Each distinct ordinate is summed once.
\end{theorem}

The left side counts excess multiplicity after all coincidences, including coincidences
across different characters, have been canceled. The case $\mathcal B=\varnothing$ gives
the unrestricted statement. For $q=1$ these are respectively the zeros and the poles, with
multiplicities, of $\zeta(s)/L(s,\eta)$.

\begin{lemma}[A uniform rational-frequency estimate; \stG]\label{zer:lem:rational}
For a primitive character $\xi$ of fixed conductor (including the zeta case), put
$D_\xi(x)=\sum_{T<\gamma_\xi\le T+H}x^{i\gamma_\xi}$. If $x=a/b>1$ is in lowest terms with
$a,b\le Q$, $Q\ge2$, then
\begin{equation}\label{zer:eq:rationallandau}
 D_\xi(x)=-\frac H{2\pi}\frac{\Lambda(x)\xi(x)}{\sqrt x}+O_\xi\bigl(Q^4\log^2(2T)\bigr),
\end{equation}
where the numerator in the main term is zero unless $x$ is an integral prime power. The
value at $x^{-1}$ is the conjugate. Moreover, in the notation above,
\begin{equation}\label{zer:eq:conductordifference}
 \sum_{T<\gamma\le T+H}z(\gamma)=-\frac{hH}{2\pi}\log d+O_{q,d}(\log(2T)).
\end{equation}
\end{lemma}
\begin{proof}
Banks proves a sharper uniform estimate for $x\ge2$ \cite[Theorem~1.2]{Banks2023Comparing};
the coarser estimate is given for the full range $x>1$ because some ratios of products of
primes approach $1$. Integrate $x^sL'(s,\xi)/L(s,\xi)$, divided by $2\pi i$, on the
rectangle with vertical sides $-1,2$. Move each horizontal side by at most $1$ to a height
at distance $\gg_\xi1/\log(2T)$ from all nearby zero ordinates; the standard local zero
count $N_\xi(t+1)-N_\xi(t-1)\ll_\xi\log(2t)$ permits this choice. The local
partial-fraction expansion for $L'/L$ then gives $O_\xi(\log^2(2T))$ on the two horizontal
sides \cite[Lemma~2.2]{Banks2023Comparing}. The horizontal integrals cost
$O_\xi(Q^2\log^2(2T))$. Restoring the endpoints costs $O_\xi(\sqrt x\log(2T))$ in the sum
of $x^\rho$ by GRH. On the right side expand $L'/L=-\sum_{n\ge2}\Lambda(n)\xi(n)n^{-s}$.
The diagonal $n=x$ gives the displayed main term before division by $\sqrt x$. For $n\ne x$
the oscillatory integral of $(x/n)^{it}$ is $O(1/\abs{\log(x/n)})$. For $x/2<n<2x$ we have
$\abs{x-n}\ge1/Q$ and $\abs{\log(x/n)}\gg Q^{-2}$; elsewhere the reciprocal is bounded.
Since $\sum_{n\ge2}\Lambda(n)/n^2<\infty$, the total off-diagonal cost is $O(Q^4)$. On the
left side the functional equation and Stirling's formula give a logarithmic term
$-\log(f_\xi t/(2\pi))$, a uniformly $O_\xi(1/t)$ remainder, and a logarithmic derivative
at $2-it$. Integration by parts bounds the logarithmic term by $O_\xi(\log(2T)/\log x)$,
with $1/\log x\ll Q$. The absolutely convergent Dirichlet series at $2-it$ has frequencies
$\log(nx)\ge\log2$, so costs $O(1)$. These estimates prove \eqref{zer:eq:rationallandau}
after division by $\sqrt x$. At frequency $1$ use the Riemann--von Mangoldt formula:
subtracting the equal numbers of degree-one functions cancels the common $t\log t$ term and
gives \eqref{zer:eq:conductordifference}.
\end{proof}

\begin{proof}[Proof of \cref{zer:thm:phasebox}]
\emph{A bounded test with a growing prime contribution.} Let $Y=c_0\log T$ for a
sufficiently small fixed $c_0>0$ to be specified by the error bound below. Let $S$ consist
of the primes $p\le Y$ outside $\mathcal B$ satisfying $p\equiv1\pmod q$ and $\eta(p)=-1$.
Define $w_p=(\log p)/\sqrt p$ and $\sigma^2=\sum_{p\in S}w_p^2$. This union of reduced
classes modulo $qd$ has density $1/(2h)$ among primes, so the prime number theorem in fixed
progressions and partial summation give $\sigma^2\sim(\log Y)^2/(4h)$. For large $T$ define
the real trigonometric polynomial
\[
 A(t)=\Im\prod_{p\in S}\Bigl(1+i\frac{w_p}\sigma\cos(t\log p)\Bigr).
\]
Pointwise $\abs A\le\exp(\sum_pw_p^2/(2\sigma^2))=\sqrt e$. Its constant term is zero. The
coefficient of $\cos(t\log p)$ arising from the singleton $p$ is $w_p/\sigma$. All
remaining nonzero terms involve at least two distinct primes, each with exponent $1$ or
$-1$ in the frequency ratio, so neither this ratio nor its inverse is an integral prime
power. At a singleton $p$ we have $\sum_\chi\chi^*(p)(1-\eta(p))=2h$. Taking the real part
in \eqref{zer:eq:rationallandau} consequently gives the weighted singleton contribution
$-hH\sigma/\pi$.

\emph{Restricting the phases.} Take a smooth function $0\le\Phi\le1$ supported inside $U$
with normalized Haar integral $\mu>0$ (for empty $\mathcal B$ take $\Phi=1$). Write
$\Phi_M$ for its rectangular Fourier truncation to exponents $\abs{k_b}\le M$, where
$M=\floor{\sqrt{\log T}}$. Smoothness gives
$\norm{\Phi-\Phi_M}_\infty\ll_{\Phi,N}M^{-N}$ for every fixed $N$ and a uniformly bounded
sum of absolute Fourier coefficients. The sum of $\abs z$ in the window is
$O_{q,d}(H\log T+\log T)$ by the zero count. Thus $\Phi_M$ can replace $\Phi$ in sums
against $z$ and $zA$ with an error $o(H)$; here $\Phi$ in a sum over ordinates stands for
$\Phi(\theta_{\mathcal B}(\gamma))$.

In $\Phi_MA$ every frequency with a nonzero exponent in $\mathcal B$ also has a nonzero
exponent in $S$. Unique factorization makes each such frequency, and its reciprocal,
ineligible for a prime-power main term. Only the constant Fourier coefficient $\mu$ of
$\Phi_M$ couples to a singleton in $A$. For $\Phi_M$ alone the constant frequency
contributes $O_{q,d}(H+\log T)$ by \eqref{zer:eq:conductordifference}. Its nonconstant
prime-power terms contribute $O_{\Phi,q,d,\mathcal B}(H)$, since the factor
$\log b/b^{\abs k/2}$ is bounded for the fixed primes $b\in\mathcal B$ and the Fourier
coefficients are absolutely summable.

All the rational frequencies occurring here, in lowest terms, have numerator and
denominator at most
\[
 Q=\exp\Bigl(\sum_{p\le Y}\log p+M\sum_{b\in\mathcal B}\log b\Bigr)
 \le\exp\bigl(2Y+O_{\mathcal B}(M)\bigr)
\]
for large $T$. The sum of the absolute Fourier coefficients of $A$ is at most
$\prod_{p\in S}(1+w_p/\sigma)\le\exp(\sqrt{\abs S})$ by Cauchy--Schwarz. Using
\eqref{zer:eq:rationallandau} termwise therefore costs
\[
 O_{\Phi,q,d}\bigl(Q^4\log^2(2T)\exp(\sqrt{\abs S})\bigr).
\]
Choose, for example, $c_0=\alpha/32$; this error is at most $T^{\alpha/4+o(1)}=o(H)$. We have proved the two estimates
\[
 \sum_{T<\gamma\le T+H}z(\gamma)\Phi=O(H),\qquad
 \sum_{T<\gamma\le T+H}z(\gamma)\Phi A(\gamma)=-\frac{\mu h}\pi H\sigma+o(H).
\]
The functions $\Phi(\sqrt e-A)$ and $\Phi(\sqrt e+A)$ take their values in
$[0,2\sqrt e]$ and vanish outside the box. Their signed sums are, respectively,
$+(\mu h/\pi)H\sigma+O(H)$ and $-(\mu h/\pi)H\sigma+O(H)$. Dividing the first by $2\sqrt e$
bounds the positive mass from below, and negating the second does the same for the
negative mass. Since $\sigma\asymp\log\log T$, \eqref{zer:eq:phase} follows.
\end{proof}

Liu and Zaharescu study smooth tests on multidimensional tori in races between the zeros of
two Dirichlet $L$-functions \cite{LiuZaharescu2021Races}. In \cref{zer:thm:phasebox} the
cancellation is between two whole products, the angular box is fixed arbitrarily before the
amplifier is chosen, and the same height window works in both directions. The count allows
different multiplicities at a common zero; simplicity, disjointness of zero sets, and the
conjectural $H\log T$ scale for almost all excess multiplicity are stronger assertions.

\paragraph{Full rank of the multiplicity vectors.}
The next statement is used in \cref{zer:sec:balancing} and \cref{zer:sec:cov}.

\begin{corollary}[Multiplicity rank; \stU]\label{zer:thm:rank}
Fix $q\ge1$ and put $n_q=\varphi(q)$. There is $K_q$ such that if $T$ is sufficiently large
and $K_q\log T\le H\le T$, the vectors
\[
 M(\rho)=\bigl(\ord_\rho L(s,\chi^*)\bigr)_{\chi\bmod q}\qquad
 (0<\Re\rho<1,\ T<\Im\rho<T+H)
\]
span $\R^{n_q}$. Distinct zeros index the vectors; a zero absent from one $L$-function
contributes a zero entry. Under \cref{src:thm:dgrh} the same vectors are indexed by
distinct positive ordinates.
\end{corollary}
\begin{proof}
For large $T$ the window $(T,T+H)$ contains a window of the much shorter length allowed in
\cref{rec:thm:localrank}, around its midpoint, so the vectors $M(\rho)$ span over $\C$, and
being real they span $\R^{n_q}$.
\end{proof}

This yields a nonsingular integer multiplicity matrix of size $\varphi(q)$ in every such
height window, including when different characters share zeros. In particular one can
choose $\gg_qT/\log T$ pairwise disjoint bases with ordinates in $[T,2T]$, by dividing it
into intervals of length $2K_q\log(2T)$. The rank argument keeps $q$ fixed; nonsingularity
places no condition on rational linear relations among the heights themselves.

\subsection{Blurred zero measures and a critical width}\label{rec:sec:blur}
\sources{dirichlet-zero-recovery}

For a primitive $\chi$ of conductor $q_\chi$ and parity $\epsilon_\chi\in\set{0,1}$ let $Z_\chi=\sum_\gamma m_\chi(\gamma)\delta_\gamma$ over its distinct
ordinates, both signs kept; under \cref{src:thm:dgrh} this is a measure on $\R$. For a finite family $\mathcal X$ of
distinct primitive characters and finite complex measures $\nu_\chi$ one observes
\begin{equation}\label{rec:eq:blurobs}
 M=\sum_{\chi\in\mathcal X}Z_\chi*\nu_\chi .
\end{equation}
An atom of $\nu_\chi$ at $u$ translates the whole zero list of $\chi$ by $u$; a probability law gives the intensity of the
list after each zero is displaced independently with that law. Unlike \cref{rec:thm:moving}, a single
character may carry a continuous law with unbounded support, and only the sum is observed. For $a,b>0$ let
$\mathcal M_{a,b}$ be the finite complex measures with
$\norm\nu_{a,b}=\int_{u<0}e^{-au}\dd\abs\nu(u)+\int_{u\ge0}e^{bu}\dd\abs\nu(u)<\infty$, and $\mathcal M_d=\mathcal M_{d,d}$. The bound
$Z_\chi([v,v+1])\ll_\chi\log(3+\abs v)$ \cite{BennettMartinOBryantRechnitzer2021Counting} makes $M$ locally finite and tempered.

\begin{theorem}[Sharp recovery threshold; \stG]\label{rec:thm:blur}
If $a+b\ge\pi$, the map \eqref{rec:eq:blurobs} is injective on $\mathcal M_{a,b}^{\mathcal X}$ for every finite family $\mathcal X$. If $0<a+b<\pi$, it
is injective exactly when $\mathcal X$ is a singleton, a pair of nonprincipal characters of opposite parity, or $\zeta$
together with one nonprincipal character; for every other family, distinct tuples of smooth strictly
positive probability densities with finite $\norm\cdot_{a,b}$ give the same $M$.
\end{theorem}

At $a+b\ge\pi$ the family need not be known: apply the theorem to the union of two candidate families,
completing with zero measures. For symmetric tails the critical exponent is $\pi/2$.

\begin{theorem}[One law erases all prime frequencies; \stG]\label{rec:thm:blurlaw}
There is a smooth, strictly positive, even probability density $K$ with $\int e^{d\abs u}K(u)\dd u<\infty$ for every
$d<\pi/2$ such that $Z_\chi*K=Z_{\chi'}*K$ if and only if $q_\chi=q_{\chi'}$ and $\chi(-1)=\chi'(-1)$. For a finite multiset of primitive
characters blurred by $K$, the combined observation determines exactly the number of factors, the number
of odd factors, the number of $\zeta$ factors and the product of the conductors.
\end{theorem}

For example, the quartic character modulo $5$ with $\chi(2)=i$ and its conjugate have the same conductor and
parity, so their zero measures become identical after convolution with $K$, while \cref{rec:thm:blur} forbids this
for every law with a finite moment at the critical exponent; the difference between the endpoint moment
and all smaller ones is therefore necessary.

Write $\widehat f(x)=\int f(t)e^{-itx}\dd t$, extended by duality, $N_\chi(z)=\int e^{-izu}\dd\nu_\chi(u)$, bounded and holomorphic on
$-a<\Im z<b$, so that $\widehat M=\sum_\chi N_\chi\widehat Z_\chi$ by Fubini and the unit-interval bound.

\begin{lemma}[The transform of a zero measure; \stG]\label{rec:lem:blurspec}
At $x=\log n$ and $x=-\log n$, $n\ge2$, the Dirac coefficients of $\widehat Z_\chi$ are $-\Lambda(n)\overline{\chi(n)}/\sqrt n$ and
$-\Lambda(n)\chi(n)/\sqrt n$. At $0$ there is the conductor mass $\log(q_\chi/\pi)\delta_0$, and elsewhere $\widehat Z_\chi$ has the even density
\[
 D_\chi(x)=-\frac{e^{-(1/2+\epsilon_\chi)\abs x}}{1-e^{-2\abs x}}+\mathbf 1_{q_\chi=1}\,2\cosh(x/2).
\]
\end{lemma}

\begin{proof}
Insert $h=\widehat\varphi$, $\varphi\in C_c^\infty(\R)$, in \eqref{rec:eq:gwstrip}, where \cref{src:thm:dgrh} makes every $z_\rho$ real; then
$\widehat h(v)=2\pi\varphi(-v)$, which fixes the conjugations. Up to constants, the digamma series for
$\Re\psi((\tfrac12+\epsilon+it)/2)$ has the terms $-2d_j/(d_j^2+t^2)$, $d_j=2j+\epsilon+\tfrac12$, with transforms $-2\pi e^{-d_j\abs x}$; the partial sums
are $O(\log(3+\abs t))$ uniformly, so they converge as tempered distributions, and off $0$ the exponential
series sums to the first term of $D_\chi$. The pole terms give $2\cosh(x/2)$.
\end{proof}

Let $Q$ be a common multiple of the conductors, $G=(\Z/Q\Z)^\times$, and inflate the characters to $G$ (they stay
distinct). For $r\in G$ put
\begin{equation}\label{rec:eq:blurrow}
 A_r(z)=\sum_{\chi\in\mathcal X}\overline{\chi(r)}N_\chi(z),\qquad N_\chi(z)=\frac1{\abs G}\sum_{r\in G}\chi(r)A_r(z).
\end{equation}
By \cref{rec:lem:blurspec}, the atom of $\widehat M$ at $\log p$, divided by $-\log p/\sqrt p$, is $A_r(\log p)$ for $p\equiv r\pmod Q$, and the
atom at $-\log p$ gives $A_r(-\log p)$ for $p\equiv r^{-1}\pmod Q$.

\begin{lemma}[Prime logarithms determine a strip function; \stU]\label{rec:lem:blurstrip}
Let $a+b\ge\pi$. If $F$ is bounded and holomorphic on $-a<\Im z<b$ and $F(\log p)=0$ for all primes $p$ in one reduced
class modulo $Q$, then $F=0$.
\end{lemma}

\begin{proof}
With $W_0=a+b$ and $c_0=(b-a)/2$, $w=\exp(\pi(z-ic_0)/W_0)$ maps the strip onto the right half-plane and $(w-1)/(w+1)$
onto the disk. The zero at $\log p$ goes to a point at distance $\asymp_{a,b}p^{-\pi/W_0}$ from the circle, since
$1-\abs{(w_p-1)/(w_p+1)}^2=4\Re w_p/\abs{w_p+1}^2$ and $\abs{c_0}<W_0/2$. As $\pi/W_0\le1$ and $\sum1/p$ diverges in each class
\cite{Williams1974Mertens}, the Blaschke condition fails.
\end{proof}

\begin{proof}[Proof of \cref{rec:thm:blur}, the case $a+b\ge\pi$]
Apply \cref{rec:lem:blurstrip} to the difference of the rows of two candidate tuples; orthogonality in
\eqref{rec:eq:blurrow} and uniqueness of Fourier transforms of finite measures finish the proof.
\end{proof}

In the symmetric critical strip the disk coordinate is $\tanh(z/2)$ and the deficit at a prime is $2/(p+1)$.
Any set of primes whose reciprocal sum diverges in each class suffices, and laws supported on a fixed
half-line are recovered from one positive exponential moment on the unbounded side.

\begin{proof}[Proof of \cref{rec:thm:blur}, the case $a+b<\pi$]
For a difference of two tuples with the same $M$, let $N_\chi$ be the difference transforms,
$E=\sum_{\epsilon_\chi=0}N_\chi$ (including $\zeta$), $O=\sum_{\epsilon_\chi=1}N_\chi$, $Z=N_\zeta$ if $\zeta\in\mathcal X$ and $0$ otherwise, and $H(z)=2+2\cosh z$.
The continuous part of \cref{rec:lem:blurspec} gives $E+e^{-x}O=P(x)Z$ for $x>0$ and $E+e^xO=P(-x)Z$ for $x<0$ between
prime frequencies, with $P(z)=(e^z+1)(1-e^{-2z})=H(z)(1-e^{-z})$; by continuation both hold on the strip, and
subtracting gives
\begin{equation}\label{rec:eq:blurpar}
 E=HZ,\qquad O=-HZ.
\end{equation}
So without $\zeta$ both parity sums are determined, hence all measures if each parity occurs at most once; with
$\zeta$ and one other character, $O=0$ or $E=Z$, and \eqref{rec:eq:blurpar} forces $Z=0$. This gives the injective cases.
For the converse put $c_0=(b-a)/2$ and
\[
 F_{c_0}(z)=e^{-(z-ic_0)\sinh(z-ic_0)}\sin(\pi e^z)\sin(\pi e^{-z}),
\]
entire, nonzero (evaluate at $\log\tfrac32$), vanishing to order two at $0$ and at every $\pm\log n$, $n\ge2$. On each substrip
$\abs{\Im z-c_0}\le d<\pi/2$ it decays faster than every $e^{-L\abs{\Re z}}$, since
$\Re((x+iv)\sinh(x+iv))=x\sinh x\cos v-v\cosh x\sin v$, and $F_{c_0}(-\bar z)=\overline{F_{c_0}(z)}$; contour shifts make its inverse
transform a smooth real density of finite $(a,b)$ norm, also after multiplication by $H$, a polynomial in $\cosh z$
or $z^{2k}$. If two nonprincipal characters have the same parity give them $F_{c_0},-F_{c_0}$; if $\zeta$, an even
nonprincipal and an odd character occur, give them $F_{c_0}$, $(H-1)F_{c_0}$, $-HF_{c_0}$; all others get $0$. These cover every
other family. Then \eqref{rec:eq:blurpar} holds identically, every prime-power atom and the conductor term vanish,
and the observation vanishes also at $0$, since $x^2\widehat g_\epsilon=-\widehat{g_\epsilon''}$ is continuous by Stirling and the
multipliers are smoothly divisible by $x^2$. If $f_j$ are the real inverse transforms ($\int f_j=0$),
$b_j=g_j/\int g_j$ with $g_j=\sqrt{f_j^2+e^{-2t^2}}$, then for small $\lambda>0$ the tuples $(b_j+\lambda f_j)$ and $(b_j-\lambda f_j)$ are strictly
positive smooth probability densities with the same observation.
\end{proof}

For $c_0=0$ the densities are even and have every subcritical exponential moment, and the factor $z^{2k}$
makes any finite number of ordinary moments agree.

\begin{lemma}[Positivity by convolution; \stU]\label{rec:lem:blurpos}
Let $A_0>0$, let $h$ be real with $\int\abs{h(s)}e^{d\abs s}\dd s<\infty$ for $0\le d<A_0$, and suppose its bilateral Laplace
transform $L(y)=\int h(s)e^{ys}\dd s$ satisfies $\inf_{\abs y<A_0}L(y)=\mu_0>0$. There is an even integrable $k>0$ with the same
subcritical moments and $(h*k)(t)\ge(\mu_0/2)k(t)$ for all $t$.
\end{lemma}

\begin{proof}
Choose $0=d_0<d_1<\dots\uparrow A_0$, increasing $L_j\ge j$ with $\int_{\abs s>L_j}\abs{h(s)}e^{d_j\abs s}\dd s<\mu_0/4$, and decreasing $\eta_j>0$
with $(\eta_jL_j^2/2)\int\abs{h(s)}e^{d_j\abs s}\dd s<\mu_0/4$. Let $V$ be even, convex and $C^1$ with $V(0)=V'(0)=0$, whose derivative on
the positive half-line rises in stage $j$ from $d_{j-1}$ to $d_j$ at rate $\eta_j$ and then stays at $d_j$ for length $L_{j+1}$; put
$k=e^{-V}$. In stage $j$, $\abs{V'}\le d_j$ and $V'$ is $\eta_j$-Lipschitz on $[t-L_j,t+L_j]$, so by convexity and Taylor
$\abs{k(t-s)/k(t)-e^{sV'(t)}}\le e^{d_j\abs s}\eta_jL_j^2/2$ for $\abs s\le L_j$, and the difference is at most $e^{d_j\abs s}$ beyond. Hence
$\abs{(h*k)(t)/k(t)-L(V'(t))}<\mu_0/2$; and $V(t)/\abs t\to A_0$.
\end{proof}

\begin{proof}[Proof of \cref{rec:thm:blurlaw}]
Let $S(w)=\sin(\pi w)/(1-w)$, $S(1)=\pi$, and $U_0(z)=e^{-z\sinh z}S(e^z)S(e^{-z})$. Its inverse transform $h$ is real, even, smooth,
with all exponential moments of rate below $\pi/2$, and $L(y)=U_0(iy)=e^{y\sin y}\abs{S(e^{iy})}^2$ extends to a positive
continuous function on $[-\pi/2,\pi/2]$, since the only integer on that half of the unit circle is $1$. By
\cref{rec:lem:blurpos}, $h*k>0$; normalized, it is $K$, with
\begin{equation}\label{rec:eq:blurkill}
 \widehat K(0)=1,\qquad\widehat K(\pm\log n)=0\quad(n\ge2).
\end{equation}
For two packets let $e,o,z_0$ be the differences of their numbers of even, odd and $\zeta$ factors. Equal
observations give $e=H(x)z_0$ and $o=-H(x)z_0$ near $0$ by \eqref{rec:eq:blurpar}, so $e=o=z_0=0$; then the gamma and pole
densities cancel and what remains is the mass $\log$ of the ratio of the conductor products at $0$.
Conversely, equal invariants cancel these terms and \eqref{rec:eq:blurkill} kills all arithmetic atoms.
\end{proof}

For $\abs{c_0}<\pi/2$ the same argument applied to $e^{c_0t}h_{c_0}(t)$, with $h_{c_0}$ the inverse transform of
$U_{c_0}(z)=e^{-(z-ic_0)\sinh(z-ic_0)}S(e^z)S(e^{-z})$, gives a positive law with \eqref{rec:eq:blurkill} and finite bilateral Laplace moments
at every parameter in $(c_0-\pi/2,c_0+\pi/2)$, so every subcritical asymmetric moment class contains such a law.
Positive convolution factors for signed functions of positive integral are studied in
\cite{RuzsaSzekely1983Convolution}.

\paragraph{Finite prime data and the endpoint.}
At the symmetric endpoint $A_0=\pi/2$ let $M_0=\sum_\chi\norm{\nu_\chi}_{A_0,A_0}$ bound a row $A_r$ on the strip. For $t\in\R$ put $w=\tanh(t/2)$,
take the distinct disk nodes $s_j=(p-1)/(p+1)$ for $p\le X$, $p\equiv r$, and $s_j=-(p-1)/(p+1)$ for $p\le X$, $p\equiv r^{-1}\pmod Q$, the
Blaschke product $B(w)=\prod_j(w-s_j)/(1-s_jw)$ and the measured values $D_j=A_r(2\operatorname{artanh}s_j)$.

\begin{proposition}[A convergent reconstruction; \stG]\label{rec:prop:blurdecode}
The interpolant $I_X(w)=\sum_jD_jB(w)/(B'(s_j)(w-s_j))$ satisfies $\abs{A_r(t)-I_X(w)}\le M_0\cosh(t/2)\abs{B(w)}$, and
\[
 \abs{B(w)}=\prod_{\substack{p\le X\\p\equiv r}}\Bigl\lvert\frac{p-e^t}{p+e^t}\Bigr\rvert\prod_{\substack{p\le X\\p\equiv r^{-1}}}\Bigl\lvert\frac{p-e^{-t}}{p+e^{-t}}\Bigr\rvert\ll_{Q,t}(\log X)^{-4\cosh(t)/\varphi(Q)},
\]
uniformly in $r$. Errors $\delta_j$ in the normalized atoms add $\sum_j\delta_j\abs{B(w)/(B'(s_j)(w-s_j))}$.
\end{proposition}

\begin{proof}
For $f(w)=A_r(2\operatorname{artanh}w)$, bounded by $M_0$ on the disk, integrate $f(v)/(B(v)(v-w))$ over circles of radius tending
to one; the residues give $(f(w)-I_X(w))/B(w)$, $\min\abs B\to1$ on these circles, and the mean of $\abs{e^{i\theta}-w}^{-1}$ is at
most $(1-\abs w^2)^{-1/2}=\cosh(t/2)$. For the rate expand $\log\abs{(p-u)/(p+u)}=-2u/p+O_u(p^{-3})$ with $u=e^{\pm t}$ and use
Mertens' theorem in progressions \cite{Williams1974Mertens}.
\end{proof}

For the quartic example ($Q=5$, $r=2$, $r^{-1}=3$) the product is $0.06579295$ at $t=0$, $X=100$, and $0.02623577$ at
$X=10^5$. Exact erasure has a finite cost: if a probability law $K_*$ with $M_*=\int e^{A_0\abs u}\dd K_*(u)$ satisfies
$\widehat K_*(\pm\log n)=0$ for $2\le n\le L$, dividing its transform by the disk factors and evaluating at $w=0$ gives
$M_*\ge\prod_{n=2}^L((n+1)/(n-1))^2=(L(L+1)/2)^2$.

\begin{corollary}[Polynomial losses at the endpoint; \stG]\label{rec:cor:blurend}
Injectivity holds if for each measure and some $k_0\ge0$, $\int e^{\pi\abs u/2}(1+\abs u)^{-k_0}\dd\abs{\nu_\chi}(u)<\infty$.
\end{corollary}

\begin{proof}
For $\beta<A_0$ a difference row is $O((A_0-\beta)^{-k_0})$ on $\abs{\Im z}<\beta$. With $s=\pi/(2\beta)>1$ and $t>0$ fixed, the zeros at
$\log p$, $p>2e^t$, $p\equiv r$, and $\log((1-v)/(1+v))\le-2v$ give
$\abs{A_r(t)}\ll(A_0-\beta)^{-k_0}\exp[-2e^{st}\sum_{p>2e^t,\,p\equiv r}p^{-s}]$, and
$\sum_{p\equiv r}p^{-s}=\varphi(Q)^{-1}\log(1/(s-1))+O_Q(1)$ as $s\downarrow1$, with $s-1\asymp A_0-\beta$; so the bound tends to zero when
$2e^t/\varphi(Q)>k_0$, the row vanishes on a ray and hence everywhere.
\end{proof}

So the density $K$ of \cref{rec:thm:blurlaw} fails every polynomially discounted critical moment.

\begin{proposition}[Sparse high windows; \stG]\label{rec:prop:blurwin}
Let $T_j\to\infty$ arbitrarily fast and $H_j-\pi^{-1}\log\log\log T_j\to\infty$, and assume
$\int\log(3+\abs y)\dd\abs{\nu_\chi}(y)<\infty$. The restrictions of $M$ to $[T_j-H_j,T_j+H_j]$, even with error measures whose total
variation on the windows tends to zero, determine $\sum_\chi\chi(p)\widehat\nu_\chi(-\log p)$ and $\sum_\chi\overline{\chi(p)}\widehat\nu_\chi(\log p)$ for
every prime $p$. Hence the injectivity of \cref{rec:thm:blur} for $a+b\ge\pi$ holds from these windows alone.
\end{proposition}

\begin{proof}
Fix $p$, $u=\log p$, $m\ge4$, and use the test $h$ of \cref{rec:thm:localcoeff} with $T=T_j$, so $\widehat h(0)=0$, $\widehat h(u)=1$, and
translation by $y$ multiplies $\widehat h(v)$ by $e^{ivy}$. By \cref{rec:lem:lefreq} the other prime-power terms are $O(2^{-m})$
uniformly in $T$ and $y$. The conductor constant integrates to zero; integration by parts bounds the gamma term
by $u^{-1}\int K(s-T)\abs{G'_\chi(s-y)}\dd s$, which tends to zero for each $y$, and so do the pole terms, by
\cref{rec:lem:lekernel}; bounded convergence handles integration against $\nu_\chi$. The unit-interval count and the
logarithmic moment give $\abs M([v,v+1])\ll\log(3+\abs v)$, and outside $\abs{t-T}\le D+2B$ the kernel is
$\ll_{p,m}e^{-R/4}e^{-(\abs{t-T}-B)}$, so the discarded part is $O(e^{-R/4}\log T)=o(1)$. Thus
$E_{p,m,j}=-(2\pi\sqrt p/u)\int_{T_j-H_j}^{T_j+H_j}h\dd M$ has
$\limsup_j\abs{E_{p,m,j}-\sum_\chi\chi(p)\widehat\nu_\chi(-u)}\le C2^{-m}\sum_\chi\norm{\nu_\chi}_{\mathrm{TV}}$; let $m\to\infty$, and use $\bar h$ for the other
reading.
\end{proof}

\subsection{Weights and multiplicity vectors from poles of prime series}\label{rec:sec:poles}
\sources{prime_pole_windows; dirichlet_zero_recovery; large_kernels_and_dirichlet_zeros; dirichlet_inverse_problems; local_euler_recovery; translated_dirichlet_zero_divisors}

Fix $q\ge2$, $G=(\Z/q\Z)^\times$, a weight $f:G\to\C$, and put
\[
 F_f(s)=\sum_{p\nmid q}\frac{f(p\bmod q)\log p}{p^s},\qquad G^m=\set{a^m:a\in G},\qquad U_mf(a)=f(a^m).
\]
With $c_\chi(g)=\abs G^{-1}\sum_ag(a)\overline{\chi(a)}$, let $n_\psi(\rho)=\ord_\rho L(s,\psi^*)$ for
$\rho$ in the open strip and $Z_g(\rho)=\sum_\psi c_\psi(g)n_\psi(\rho)$. The question is what the
poles and residues of $F_f$, or of the progression series, determine about $f$ and about the
multiplicity vectors $(n_\psi(\rho))_\psi$, which may have several nonzero entries at a common
zero.

\begin{proposition}[Residues on power lines; \stG]\label{rec:prop:powerline}
$F_f$ continues meromorphically to $\Re s>0$ with at most simple poles. For $\rho=\tfrac12+i\gamma$,
$\gamma\ne0$, and $m\ge1$,
\[
 \Res_{s=\rho/m}F_f(s)=-\frac{\mu(m)}m\,Z_{U_mf}(\rho),
\]
and every nonreal pole in $\Re s>0$ has this form. In particular only $f|_{G^m}$ enters the line
$\Re s=1/(2m)$, and for nonsquarefree $m$ that line carries no nonreal pole.
\end{proposition}

\begin{proof}
For $\Re s>1$ the Euler products and M\"obius inversion give
\[
 F_f(s)=-\sum_{k\ge1}\mu(k)\sum_\chi c_\chi(f)\frac{L'}L(ks,\chi^k)
 =-\sum_{k\ge1}\mu(k)\sum_\psi c_\psi(U_kf)\frac{L'}L(ks,\psi),
\]
since the coefficient of $p^{-js}$, $p\nmid q$, is $f(p^j)\log p\sum_{k\mid j}\mu(k)$. On
$\Re s\ge\sigma_0>0$ the terms with $k\sigma_0\ge2$ converge normally, each logarithmic derivative
being at most $2^{-(k\sigma_0-2)}\sum_n\Lambda(n)n^{-2}$ and $\sum_\psi\abs{c_\psi(U_kf)}\le\sqrt{\abs G}\norm f_\infty$.
Omitted Euler factors have zeros on $\Re s=0$, and by \cref{src:thm:dgrh} a nonreal pole of the
$k$-term lies on $\Re s=1/(2k)$; a simple pole of $L'/L(ms,\psi)$ at $\rho/m$ has residue
$n_\psi(\rho)/m$.
\end{proof}

\paragraph{Triple-logarithmic windows.}
The multiplicity vectors have full rank in very short windows, without any hypothesis on the
zeros. Write $\mathcal Z_q$ for the set of distinct nontrivial zeros of the primitive inducers of the
characters modulo $q$. These inducers are distinct and form a family
closed under conjugation, so \cref{rec:thm:localrank} shows that the vectors $(n_\psi(\rho))_\psi$,
$\rho\in\mathcal Z_q$, $\abs{\Im\rho-T}\le\pi^{-1}\log\log\log\abs T+C_q$, span $\C^{\widehat G}$ for all large $\abs T$. The
local coefficient statistic of \cref{rec:thm:localcoeff} makes this quantitative.

\begin{corollary}[An approximate inverse in each window; \stU]\label{rec:cor:localmap}
Fix $q$ and $\delta>0$. There is $C_{q,\delta}$ such that for all sufficiently large $\abs T$ there is a linear
map $D_T$ from the data $\set{Z_g(\rho)}$, $\rho\in\mathcal Z_q$, $\abs{\Im\rho-T}\le\pi^{-1}\log\log\log\abs T+C_{q,\delta}$, to
$\C^G$ with $\norm{D_TZ_g-g}_\infty\le\delta\norm g_\infty$ for all $g$.
\end{corollary}

\begin{proof}
Choose a prime $p_a$ in each class $a$, fix $m$, and for $T>0$ let $h_a$ be the test of
\cref{rec:thm:localcoeff} for $p_a$, $m$ and $Q=q$, with window radius $H_a=16p_am+D(q,T)$. Put
$(D_TZ)(a)=-\frac{2\pi\sqrt{p_a}}{\log p_a}\sum_{\abs{\Im\rho-T}\le H_a}h_a(z_\rho)Z(\rho)$. For $g=\psi$ one has
$Z_\psi(\rho)=n_\psi(\rho)$, and $(D_TZ_\psi)(a)$ is the zero sum of $\mathcal E_{p_a}$ for the inducer of $\psi$; by the
theorem and the bound $O(\sqrt p/(T\log^2p))$ for the known integral (proof of \cref{rec:thm:localid}) it
differs from $\psi(a)$ by at most $C2^{-m}+o(1)$. By linearity and
$\sum_\psi\abs{c_\psi(g)}\le\sqrt{\abs G}\norm g_\infty$, choosing $m$ and then the height threshold gives the claim, with
$C_{q,\delta}=\max_a16p_am+O(1)$. For $T<0$ use $\overline{Z_g(\rho)}=Z_{\bar g}(\bar\rho)$.
\end{proof}

\begin{theorem}[Local recovery of the weights; \stG]\label{rec:thm:localrecovery}
Fix $b>1/\pi$. There are $C,T_0$ depending only on $q,b$ such that for every $m\ge1$ and
$\abs T\ge T_0$, the poles and residues of $F_f$ on $s=(\tfrac12+i\gamma)/m$,
$\abs{\gamma-T}\le b\log\log\log\abs T+C$, determine exactly $f|_{G^m}$ when $m$ is squarefree: two
weights give the same observations if and only if they agree on $G^m$, and the observation map has
rank $\abs{G^m}$. For nonsquarefree $m$ it has rank zero.
\end{theorem}

\begin{proof}
By \cref{rec:prop:powerline} the normalized residues are $Z_{U_mf}(\rho)$ on the window, which
misses the real axis for large $T_0$, and \cref{rec:thm:localrank} recovers $U_mf$, hence $f$ on
$G^m$. Conversely all residues factor through $f|_{G^m}$, and every function on $G^m$ extends to $G$.
\end{proof}

Exact recovery of complex weights is a finite calibration: by \cref{rec:cor:localmap} with
$\delta=\tfrac12$ the matrix with columns $D_TZ_{\mathbf 1_{\set{a'}}}$, computed from the known family, is within
$1/2$ of the identity, so its inverse has norm at most $2$.

\begin{corollary}[Lines with poles; natural boundary; \stG]\label{rec:cor:boundary}
The line $\Re s=1/(2m)$ contains a nonreal pole of $F_f$ exactly when $m$ is squarefree and
$\supp f\cap G^m\ne\varnothing$, and then it has poles in every sufficiently high window of the
theorem, in both directions. If $f\ne0$, the imaginary axis is a natural boundary of $F_f$.
\end{corollary}

\begin{proof}
Compare $f$ with $0$. For every prime $m$ not dividing the exponent of $G$ the power map is an
automorphism, so such lines carry poles; for real $t\ne0$ take $T=mt$ to get poles $s_m$ with
$\Re s_m=1/(2m)$ and $\abs{\Im s_m-t}\le(b\log\log\log\abs{mt}+C)/m$, accumulating at $it$ (for $t=0$ use one
large $T$).
\end{proof}

The windows in \cref{rec:thm:localrecovery,rec:cor:boundary} are those of the source. The spanning
statement \cref{rec:thm:localrank} used in their proofs now holds at the endpoint radius
$\pi^{-1}\log\log\log\abs T+C$.

Let $P$ be the set of prime divisors of the exponent of $G$, $S(m)=\set{\ell\in P:\ell\mid m}$ and
$B(y)=\set{\ell\in P:y\in G^\ell}$.

\begin{theorem}[Prescribing the lines; \stU]\label{rec:thm:design}
On squarefree indices, the patterns of lines with nonreal poles are exactly the downward closed
families of sets $S(m)$, and the least $\abs{\supp f}$ realizing a family is its number of maximal
members.
\end{theorem}

\begin{proof}
On the Sylow $\ell$-part an exponent prime to $\ell$ acts invertibly, so $G^m=\bigcap_{\ell\in S(m)}G^\ell$
for squarefree $m$. A supported point $y$ makes exactly the subsets of $B(y)$ visible. Every
$B\subseteq P$ is some $B(y)$: take the identity on the $\ell$-part for $\ell\in B$ and a non-$\ell$th
power for $\ell\notin B$. Each maximal member needs its own point.
\end{proof}

Among squarefree indices a family $\mathcal D$ has relative density
$\sum_{S\in\mathcal D}\prod_{\ell\in S}\frac1{\ell+1}\prod_{\ell\in P\setminus S}\frac\ell{\ell+1}$. Modulo
$31$, for $f=\delta_9-\delta_{27}$ (that is, $3^2$ and $3^3$) the lines carry poles exactly when
$5\nmid m$ and $6\nmid m$, a family with maximal sets $\set2,\set3$ and density $55/72$, so two
supported classes are needed; for $f=\delta_3$ the density is $5/12$. Modulo $4$, a simple zero of
the Dirichlet beta function not shared by $\zeta$ and a simple zeta zero not shared by it give, after
the sign change, $\tfrac12(f(1)-f(3))$ and $\tfrac12(f(1)+f(3))$; such zeros occur at heights
$12.98809801\ldots$ and $14.13472514\ldots$ (\artifact{prime_pole_windows/code/window_demo.txt}).

For a pair of characters, every window of radius $\pi^{-1}\log\log\log T+O(1)$ contains zeros with
$n_{\chi_1}>n_{\chi_2}$ and with $n_{\chi_2}>n_{\chi_1}$, for arbitrary conductors and every balanced real
direction of a larger family (\cref{rec:thm:bothsigns}).

\paragraph{Dyadic blocks with corrupted observations.}
Now let $\mathcal Z_T$ be the set of distinct $\gamma\in(T,2T)$ at which some $\nu_\psi(\gamma)=n_\psi(\tfrac12+i\gamma)$
is nonzero, $\nu(\gamma)=(\nu_\psi(\gamma))_\psi$, $d=\abs G$ and $L_T=T\log T$. Observations are
the normalized residues $X_{m,f}(\gamma)=\nu(\gamma)\cdot\widehat g$, $g=U_mf$, at positions matched to
the known zeros with their multiplicities.

\begin{theorem}[A positive proportion of distinct pole ordinates; \stG]\label{rec:thm:distinct}
There are $c_q>0$ and $T_q$ such that for $T\ge T_q$ and every nonzero $b\in\C^{\widehat G}$,
\[
 c_qT\log T\le\#\set{\gamma\in\mathcal Z_T:\textstyle\sum_\psi b_\psi\nu_\psi(\gamma)\ne0}\ll_qT\log T .
\]
\end{theorem}

\begin{proof}
The family of primitive inducers satisfies the hypotheses of the joint short-interval limit theorem
of Bombieri and Perelli \cite[Theorem~2]{BombieriPerelli1998Distinct}: Euler products of degree one,
functional equations, $\sum_{p\le x}\psi^*(p)\overline{\eta^*(p)}/p=\mathbf 1_{\psi=\eta}\log\log x+C+O_q(1/\log x)$,
and, by \cref{src:thm:dgrh}, no zeros to the right of $\sigma=1/2$. (For a single character adjoin a
second one and take a marginal.) With $t$ uniform on $[T,2T]$, $M\to\infty$,
$M\le\log T/\log\log T$, $h=M/\log T$ and $S_\psi(t)=\pi^{-1}\arg L(\tfrac12+it,\psi^*)$, the vector
$(\Delta_hS_\psi(t)/\sqrt{\log M})_\psi$ converges in law to a centred Gaussian vector $Y$ with
independent coordinates of positive variance, and $\Delta_hN_\psi(t)=M/(2\pi)+\Delta_hS_\psi(t)+O_q(1+h)$.
If the lower bound failed along $T\to\infty$, choose unit real vectors $d_T\to d$ (real or imaginary
parts of $b$) with supports of size $K=o(T\log T)$ and
$M=\min\set{\log T/\log\log T,\sqrt{T\log T/(1+K)}}$. Outside probability $(K+1)h/T=o(1)$ the window
$(t,t+h]$ avoids the support, so $\sum_\psi d_{T,\psi}\Delta_hN_\psi(t)=0$ and
$d_T\cdot Y_T=-M\sum_\psi d_{T,\psi}/(2\pi\sqrt{\log M})+o(1)$, a deterministic quantity up to $o(1)$,
contradicting the nondegenerate limit $d\cdot Y$. The upper bound is the zero count.
\end{proof}

\begin{theorem}[Frequent independent multiplicity vectors; \stG]\label{rec:thm:frequent}
There are $\delta_q>0$ and $B_q$ such that every sufficiently high block contains $d$ linearly
independent integer vectors $v_j$ with $n_T(v_j)=\#\set{\gamma\in\mathcal Z_T:\nu(\gamma)=v_j}\ge\delta_qL_T$
and $\norm{v_j}_1\le B_q$. They can be selected from the observed frequencies alone, by sorting the
types by decreasing frequency and keeping those that increase the rank.
\end{theorem}

\begin{proof}
By the zero count, $\sum_{\gamma\in\mathcal Z_T}\norm{\nu(\gamma)}_1\le C_qL_T$. With $c=c_q$ choose $B$ with
$C_q/B<c/4$; at most $cL_T/4$ ordinates have row sum above $B$, and discarding the at most
$(B+1)^d$ remaining types with fewer than $\delta_qL_T=cL_T/(4(B+1)^d)$ occurrences removes at most
$cL_T/4$ more. A vector annihilating the remaining rows would be supported on at most $cL_T/2$
ordinates, contradicting \cref{rec:thm:distinct}. The greedy selection picks rows of frequency at
least $\delta_qL_T$, whose sums are at most $B_q=\ceil{C_q/\delta_q}$.
\end{proof}

With $A=(v_j)$ and $D_q=d!B_q^{d-1}$, the cofactor formula and $\abs{\det A}\ge1$ give
$\norm{A^{-1}}_{\infty\to\infty}\le D_q$. Hence for every $b\in\C^d$ at least $\delta_qL_T$ distinct
ordinates with $\norm{\nu(\gamma)}_1\le B_q$ have $\abs{\nu(\gamma)\cdot b}\ge D_q^{-1}\norm b_\infty$, and
on a squarefree line with $f|_{G^m}\ne0$ the residues there satisfy
$\abs{\Res}\ge\norm{f|_{G^m}}_\infty/(mdD_q)$. After removing any $\delta_qL_T/2$ ordinates,
$c_q'L_T\norm b_2^2\le\sum_{\norm{\nu(\gamma)}_1\le B_q}\abs{\nu(\gamma)\cdot b}^2\le C_q'L_T\norm b_2^2$.

\begin{proposition}[Both signs of a balanced combination; \stG]\label{rec:prop:bothsigns}
For $d\ge2$ there are $a_q,\tau_q>0$ and $B_q'$ such that for large $T$ and every real $b$ with
$\sum_\psi b_\psi=0$, $\norm b_2=1$, at least $a_qL_T$ distinct ordinates in $\mathcal Z_T$ have
$\norm{\nu(\gamma)}_1\le B_q'$ and $\nu(\gamma)\cdot b\ge\tau_q$, and at least $a_qL_T$ have
$\nu(\gamma)\cdot b\le-\tau_q$.
\end{proposition}

\begin{proof}
If the positive support had size $o(L_T)$ along a sequence, the argument of \cref{rec:thm:distinct}
would give $b_T\cdot\Delta_hN\le0$ on almost all windows; as $\sum b_{T,\psi}=0$ the drift $M/(2\pi)$
cancels, so $b_T\cdot Y_T\le o(1)$ with probability tending to one, against a centred nondegenerate
Gaussian limit. Rows with large sums are excluded by the count, and a uniform margin follows by
compactness of the balanced sphere, since finitely many integer types of bounded sum occur.
\end{proof}

At $H=T$ this gives $\gg T\log T$ ordinates in each direction, more than the mass
$\gg T\log\log T$ of \cref{zer:thm:short}, which on the other hand applies in intervals of
triple-logarithmic length.

\begin{theorem}[Stable recovery; \stG]\label{rec:thm:stable}
There are $\alpha_q,K_q>0$ such that for large $T$, every squarefree $m$ and every $f$, if at most
$\alpha_q\#\mathcal Z_T$ of the normalized residues $X_{m,f}(\gamma)$, $\gamma\in\mathcal Z_T$, are erased or
arbitrarily corrupted and the others have error at most $\eps$, then $f|_{G^m}$ is recovered with
supremum error at most $K_q\eps$, exactly if $\eps=0$.
\end{theorem}

\begin{proof}
Select the rows of \cref{rec:thm:frequent} and put $\alpha_q=\delta_q/(4C_q)$; then at least three
quarters of the observations of each selected type are good. Coordinatewise medians of real and
imaginary parts give $y_j$ with $\abs{y_j-v_j\cdot b}\le\sqrt2\eps$, and $Ab^{\mathrm{rec}}=y$ gives
$\norm{b^{\mathrm{rec}}-b}_\infty\le\sqrt2D_q\eps$. Output $\sum_\psi b_\psi^{\mathrm{rec}}\psi(a_u)$ for
$a_u^m=u$, with $K_q=d\sqrt2D_q$.
\end{proof}

For $q=13$ the recoverable classes are all twelve at $m=1$; $1,3,4,9,10,12$ at $m=2$; $1,5,8,12$ at
$m=3$; none at $m=4$; and $1,12$ at $m=6$. A floating-point illustration modulo $4$ with two of five
pairs corrupted and noise $10^{-6}$ recovers $f(1)=2+3i$, $f(3)=-1+i$ to about $7\cdot10^{-7}$
(\artifact{dirichlet_zero_recovery/layers/observed_example.txt}).

\paragraph{Residues of single progressions.}
For $a\in G$ and $h=\abs G$ let
$D_a(s)=\sum_{n\equiv a}\Lambda(n)n^{-s}=-\frac1h\sum_\chi\overline{\chi(a)}\frac{L'}L(s,\chi)$. At any
$\rho$ in the strip, with $m_\chi=\ord_\rho L(s,\chi)$,
\begin{equation}\label{rec:eq:progres}
 -h\Res_{s=\rho}D_a(s)=F_m(a):=\sum_\chi m_\chi\overline{\chi(a)} .
\end{equation}
For a cyclic subgroup $H\le G$ put $d_H=\abs H$, $D_H=\varphi(d_H)$, let $\recGen H$ be its generators
(the identity for $H=1$), fix $a_H\in\recGen H$, and $K_d=\Q(e^{2\pi i/d})$.

\begin{lemma}[Rational blocks; \stU]\label{rec:lem:blocks}
The map $v\mapsto(F_v(a_H))_H$ is a rational isomorphism $\Q^{\widehat G}\to\prod_HK_{d_H}$, and the
values at the generators of $H$ are the conjugates of $F_v(a_H)$. Hence for $A\subseteq G$,
$\dim_\Q\ker(v\mapsto F_v|_A)=\sum_{H:A\cap\recGen H=\varnothing}D_H$, and $A$ recovers every
nonnegative integer vector from exact residues if and only if it meets every $\recGen H$.
\end{lemma}

\begin{proof}
Applying $e^{2\pi i/d_H}\mapsto e^{2\pi ik/d_H}$ to $F_v(a_H)$ gives $F_v(a_H^k)$; if all coordinates
vanish, $F_v=0$ on $G$ and $v=0$. The generator sets partition $G$, so both sides have dimension
$\abs G$. A nonzero rational kernel vector is, after scaling, a difference of two distinct
nonnegative integer vectors.
\end{proof}

For $M\ge1$ let $\mathcal M_M$ be the nonnegative integer vectors of total at most $M$ and
$\delta_A(M)=\min_{m\ne n\in\mathcal M_M}\max_{a\in A}\abs{F_m(a)-F_n(a)}$. For $A$ meeting every generator set,
let $b_H$ count the inverse pairs $\set{a,a^{-1}}$ met by $A$ in $\recGen H$ when $d_H>2$, and
$\kappa(A)=\max(\set0\cup\set{D_H/(2b_H)-1:d_H>2})$.

\begin{theorem}[Precision exponent; \stU]\label{rec:thm:precision}
For such $A$ and every $M\ge1$, $(2M)^{-\kappa(A)}\le\delta_A(M)\le C_{G,A}M^{-\kappa(A)}$. If $A$ misses a
generator set, $\delta_A(M)=0$ for large $M$. Hence the least $\abs A$ with $\delta_A(M)\gg M^{-\theta}$
is $\#\set{H:d_H\le2}+\sum_{d_H>2}\ceil{\varphi(d_H)/(2(1+\theta))}$.
\end{theorem}

\begin{proof}
For $v=m-n\ne0$ some block $\alpha=F_v(a_H)$ is a nonzero algebraic integer with all conjugates at
most $2M$. If $d_H\le2$ it is an observed nonzero integer. Otherwise at least $2b_H$ embeddings are
observed, of modulus at most $\eps=\max_A\abs{F_v}$, so
$1\le\abs{N(\alpha)}\le\eps^{2b_H}(2M)^{D_H-2b_H}$. For the upper bound take a block attaining
$\kappa(A)>0$, integer vectors $v_1,\dots,v_D$ spanning its inverse image, and the $(L+1)^D$ vectors
$\sum n_jv_j$, $0\le n_j\le L$, whose observed values lie in a box of side $O(L)$ in $\R^{2b_H}$;
pigeonhole with $J=\ceil{(L+1)^{D/(2b_H)}}-1$ cubes per side gives two images at distance
$O(L^{1-D/(2b_H)})$, and $L=\floor{M/B}$ concludes. The count follows by choosing $b_H$ in each block.
\end{proof}

Modulo $11$, $A=\set{1,10,4,2}$ meets every generator set, so four classes suffice for exact data;
$\delta_A(M)\asymp M^{-1}$, a fifth class does not change the exponent, and adding both $5$ and $8$
gives $\delta_A(M)=1$; uniform separation needs $(\abs G+\abs{G[2]})/2$ classes. For actual zeros of
height at most $T$, $\sum_\chi m_\chi\ll_q\log(T+3)$, so absolute error below
$\tfrac12(2M)^{-\kappa(A)}$ with $M\asymp_q\log(T+3)$ identifies the vector.

The detection question asks less: a single weighted count should see every zero off the line. For
$q\ge3$ choose a generator $u_H$ of each nontrivial cyclic subgroup, put $w(u_H)=\varphi(\abs H)$ and
$w=0$ elsewhere, so $\sum_aw(a)=n-1$ with $n=\abs G$, and let
$D_w(x)=\sum_{k\le x}\Lambda(k)w(k\bmod q)-\frac{n-1}n\sum_{k\le x,(k,q)=1}\Lambda(k)$. Let
$\recav:K\to\Q$ be the normalized trace of $K=\Q(\zeta_{\exp G})$ and $C_\chi(v)=\sum_av(a)\overline{\chi(a)}$.
Since $\varphi(d)\recav(\overline{\chi(u)})=\sum_{a\in\recGen\langle u\rangle}\overline{\chi(a)}$,
\begin{equation}\label{rec:eq:trace}
 \recav C_\chi(w)=\sum_{a\ne e}\overline{\chi(a)}=\begin{cases}n-1,&\chi=\triv,\\-1,&\chi\ne\triv.\end{cases}
\end{equation}

\begin{theorem}[A sparse weighted count that detects zeros; \stU]\label{rec:thm:sparsecount}
The following are equivalent: \textup{(i)} every nonprincipal $L(s,\chi)$ modulo $q$ is zero-free in
$\Re s>1/2$; \textup{(ii)} $D_w(x)=O_q(\sqrt x\log^2(2x))$; \textup{(iii)} $D_w(x)=O_{q,\eps}(x^{1/2+\eps})$
for every $\eps>0$. Assuming RH, the same holds with $\sum_{(k,q)=1}\Lambda(k)$ replaced by $x$.
\end{theorem}

\begin{proof}
$D_w=\frac1n\sum_{\chi\ne\triv}C_\chi(w)\psi(x,\chi)$, so (i) gives (ii) by the explicit formula
\cite{ErnvallHytonenPalojarvi2022Explicit}, imprimitivity costing $O_q(\log x)$. Under (iii),
$s\int_1^\infty D_w(x)x^{-s-1}\dd x=-\frac1n\sum_{\chi\ne\triv}C_\chi(w)\frac{L'}L(s,\chi)$ is holomorphic on
$\Re s>1/2$. At a zero $\rho$ there with multiplicities $m_\chi$, \eqref{rec:eq:trace} gives
$\recav\Res_{s=\rho}=\frac1n\sum_{\chi\ne\triv}m_\chi>0$, a contradiction.
\end{proof}

Modulo $7$ this is $\psi_6+2\psi_2+2\psi_3$, and for cyclic $G$ of order $n$ the support has
$\tau(n)-1$ classes. Say that a weight $v$ \emph{detects multiplicities} if
$\sum_{\chi\ne\triv}m_\chi C_\chi(v)\ne0$ for every nonzero tuple of nonnegative integers.

\begin{theorem}[Detecting weights; \stU]\label{rec:thm:detecting}
A rational weight detects multiplicities if and only if the rational numbers $\recav C_\chi(v)$,
$\chi\ne\triv$, are all positive or all negative. A set $A\subseteq G\setminus\set e$ supports a rational
detecting weight if and only if it meets $\recGen H$ for every nontrivial cyclic $H$, and if it
misses one, no complex weight supported on $A$ detects multiplicities; the least support is the
number of nontrivial cyclic subgroups.
\end{theorem}

\begin{proof}
A strict sign gives detection after taking traces. Conversely the trace is constant on Galois
orbits of characters and the orbit sums are rational: an orbit with zero trace, or two orbits of
opposite signs with suitable positive integer multiples, gives an undetected pattern. If $A$ misses
$B=\recGen H$, put $m_\chi=\abs B-\sum_{a\in B}\chi(a)$, a nonnegative rational integer with
$m_\triv=0$; orthogonality gives $\sum_\chi m_\chi\overline{\chi(g)}=n(\abs B\mathbf 1_{g=e}-\mathbf 1_{g\in B})$,
which vanishes on $A$, so no weight on $A$ detects it. The generator sets are disjoint.
\end{proof}

The criterion for exact recovery of vectors (\cref{rec:lem:blocks}) includes the identity class,
which detection excludes; allowing $e$, the singleton weight $\mathbf 1_{\set e}$ detects.

\subsection{Quadratic characters from finitely many zeros}\label{rec:sec:quadratic}
\sources{quadratic_spectra}

Let $\chi$ be a primitive nonprincipal quadratic character of conductor $q$ and parity $\kappa$, and
$p_-(\chi)$ the least prime with $\chi(p)=-1$ (the least inert prime for the character of a
quadratic field). Under \cref{src:thm:dgrh} write the zeros as $\tfrac12+i\gamma$ with both signs. For
compact intervals $I,J$ of positive length put $v_I(\gamma)=\int_Ie^{i\gamma x}\dd x$ and
$C_\chi(I,J)=\Re\sum_\gamma\overline{v_I(\gamma)}v_J(\gamma)$, an absolutely convergent pairing since
$\abs{v_I(\gamma)}\le\min(\abs I,2/\abs\gamma)$. With $g_{I,J}(u)=\abs{I\cap(J-u)}$ and
$K_\kappa(u)=e^{-(\kappa+1/2)u}/(1-e^{-2u})$, the Guinand--Weil formula
\cite{ChirrePereiraDeLaat2021Primes} gives, for $J-I\subset(0,\infty)$,
\begin{equation}\label{rec:eq:packetweil}
 C_\chi(I,J)=-\int_0^\infty g_{I,J}(u)K_\kappa(u)\dd u-\sum_{n\ge2}\frac{\Lambda(n)\chi(n)}{\sqrt n}g_{I,J}(\log n).
\end{equation}
Indeed, for smooth $g$ supported in $(0,\infty)$ and $H(t)=\int g(u)\cos(tu)\dd u$ the conductor term
vanishes, the prime terms are as displayed, and the gamma term, from the integral representation of
$\psi$, is $-\int gK_\kappa$; the trapezoid $g_{I,J}$, whose cosine transform is
$\Re(\overline{v_I(t)}v_J(t))$, is reached by mollification.

\begin{theorem}[Sign threshold; \stG]\label{rec:thm:signthreshold}
The supremum $R_\chi$ of the $R>0$ such that $C_\chi(I,J)<0$ whenever $J-I\subset(0,R)$ equals
$\log p_-(\chi)$. For every $R>R_\chi$ there are equal-length intervals, of arbitrarily small length,
with $J-I\subset(0,R)$ and $C_\chi(I,J)>0$.
\end{theorem}

\begin{proof}
If $J-I\subset(0,\log p_-)$, every prime power in \eqref{rec:eq:packetweil} has $\chi\ge0$, and
$K_\kappa>0$, so the pairing is negative. For $r=\log p_-$, $I_h=[0,h]$ and $J=[r,r+h]$, the weight
$(h-\abs{u-r})_+$ sees for small $h$ only the prime $p_-$, contributing $h\log p_-/\sqrt{p_-}$, while
the gamma term is at most $h^2K_\kappa(r-h)$.
\end{proof}

With $n=p^k$, $0<h<\log(1+1/n)$ and $A_\kappa(r,h)=\int(h-\abs{u-r})_+K_\kappa(u)\dd u$, the same formula
gives the exact identity
$\chi(n)=-\frac{\sqrt n}{h\log p}\bigl(C_\chi(\log n,h)+A_\kappa(\log n,h)\bigr)$, where
$C_\chi(r,h)=\sum_\gamma\frac{4\sin^2(\gamma h/2)}{\gamma^2}\cos(\gamma r)$, and by the explicit zero count
of \cite{BennettMartinOBryantRechnitzer2021Counting} the part with $\abs\gamma>T$ is at most
$16(\log(qT)+1)/T$, or
$\frac{4(\log(qT/(2\pi))+1)}{\pi T}+\frac{4(2b(T)+0.149)}{T^2}$ with $b(T)=0.298\log(qT)+4.358$. For
the character modulo $3$ with $r=\log2$, $h=3/8$ and the $46$ positive ordinates up to $100$, the
estimate of $\chi(2)$ is $-0.98283822$, which rounds to $-1$; the cutoff bound is proved and the
evaluations are numerical (\artifact{quadratic_spectra/spectrum_example.txt}). A higher convolution
power recovers all small prime powers from one list.

For $Q\ge q$ and $X\ge2$ let $h=1/(8X)$, choose an even $r\ge2$ with $T=2r/h=16rX$ such that
\begin{equation}\label{rec:eq:choose}
 2^r\ge\frac{32}{\log2}\,r\sqrt X\log(eQT),
\end{equation}
and let $f$ be the $r$-fold convolution of the uniform density on $[-h/r,h/r]$, so
$\widehat f(t)=\recsinc(ht/r)^r$.

\begin{theorem}[Finite recovery with a rounding margin; \stG]\label{rec:thm:quaddecode}
Let $\Gamma$ be the multiset of ordinates $\abs\gamma\le T$ of $L(s,\chi)$, $\chi$ primitive
nonprincipal quadratic of conductor at most $Q$. For each $n=p^a\le X$ the nearest integer to
\[
 B(n;\Gamma)=-\frac{\sqrt n}{(\log p)f(0)}\sum_{\gamma\in\Gamma}\cos(\gamma\log n)\recsinc(h\gamma/r)^r
\]
is $\chi(n)$, including $0$ at ramified primes, simultaneously for all such $n$, and still with
approximate ordinates $\widetilde\gamma$ whenever
$\sum_\Gamma\abs{\widetilde\gamma-\gamma}\le\log2/(16\sqrt Xh(\log X+h))$. One can take
$T\ll X(\log(2X)+\log\log(3Q))$ with an absolute constant.
\end{theorem}

\begin{proof}
$f=v*v$ with $v$ an even density on an interval of length $h$, so $f(0)=\int v^2\ge1/h$. For $n\le X$
the neighbouring integer logarithms are more than $h$ away, so $g(u)=f(u-\log n)$ in the explicit
formula sees only $n$, and the infinite sum is $-\chi(n)\log p\,f(0)/\sqrt n-\int f(u-\log n)K_\kappa(u)\dd u$,
with $K_\kappa<2/\sqrt n$ on the support; normalized, the integral is at most $2h/\log2$. From
$N(t)\le2t\log(Qt)$ for $t\ge2$,
$\sum_{\abs\gamma>T}\abs{\widehat f(\gamma)}\le(r/h)^rr\int_T^\infty N(t)t^{-r-1}\dd t\le4T\log(eQT)2^{-r}$,
which after normalization is at most $1/4$ by \eqref{rec:eq:choose}. The test has Lipschitz constant
at most $\log X+h$ in $\gamma$, so the perturbation costs $1/16$, and
$1/4+1/(8\log2)+1/16<1/2$. Finally $r\asymp\log(2X)+\log\log(3Q)$ satisfies \eqref{rec:eq:choose}.
\end{proof}

With $X=C\log^2(2Q)$, admissible by the bound $O(\log^2(2M))$ of Lamzouri, Li and Soundararajan
\cite[Theorem~1.1]{LamzouriLiSoundararajan2015Conditional} for the least prime outside the kernel of a
nontrivial quadratic character modulo $M\le Q^2$, the theorem determines every primitive quadratic
character of conductor at most $Q$ from its zeros with $\abs\gamma\ll\log^2(2Q)\log\log(3Q)$;
\cref{rec:thm:lowzero} lowers this height to $O(\log^2(2Q)/\log\log(3Q))$ for characters of every order.

\begin{theorem}[Separation at every higher cutoff; \stG]\label{rec:thm:quaddisplace}
With $X,r,h,T$ as above and $H\ge2T$, suppose two primitive nonprincipal quadratic characters of
conductor at most $Q$ differ at a prime $p\le X$ and have the same number of zeros with
$\abs\gamma\le H$. Then every bijection between the two lists satisfies
$\sum\abs{\gamma-\gamma'}>H/(144\pi\sqrt X)$; for distinct characters, with $X=C\log^2(2Q)$, this is $\gg H/\log(2Q)$.
\end{theorem}

\begin{proof}
Use $\phi(u)=f(u)\recsinc^2(Bu)/f(0)$ with $B=H/8$: $\phi(0)=1$, $\int\phi\le\min(h,\pi/B)$, and
$\widehat\phi$ averages $\widehat f(t+V)/f(0)$ over $\abs V\le2B$, with $\abs{t+V}\ge3\abs t/4$ for
$\abs t>H$. The gamma term is again at most $2h/\log2$, the normalized tail is at most
$\frac{4h\sqrt X}{\log2}H\log(eQH)(4r/(3hH))^r\le1/4$ for $H\ge2T$ (since $r\ge4$), and the
normalized test has Lipschitz constant at most $9\pi\sqrt X/H$. A matching of smaller displacement
would force equal values at $p$ by rounding.
\end{proof}

For $q=Q=3$ and $X=5$ the rule gives $r=14$, $h=1/40$, $T=1120$; with $941$ positive ordinates the
finite sums at $n=2,3,4,5$ are $-0.990596$, $0.003337$, $1.003759$, $-0.998733$, rounding to
$-1,0,1,-1$ (\artifact{quadratic_spectra/decoder_example.txt}; the root enumeration is numerical).

\subsection{Characters, products and abelian fields from low zeros}\label{rec:sec:lowzero}
\sources{dirichlet_low_zero_identification; dirichlet_zeros_and_abelian_fields; robust_dirichlet_spectra_2026-11-26; quadratic_spectra; dirichlet-character-zero-windows; dirichlet-zero-reconstruction; dirichlet_zero_windows; dirichlet-zero-windows; grh-moving-zero-windows}

Let $\mathcal C(Q)$ be the set of primitive characters of conductor at most $Q\ge3$, the
conductor-one character included, $L_Q=\log(2Q)$, and $m_\chi(t)=\ord_{s=1/2+it}L(s,\chi)$. For
$\tau\in\R$, $H>0$ and $F\ge1$ put
\[
 I=I(\tau,H)=[\tau-H,\tau+H],\qquad J(F,\tau)=\log\log\bigl(e^eF(3+\abs\tau)\bigr),\qquad
 L_3(\tau)=\log\log\log\bigl(e^{e^e}+\abs\tau\bigr).
\]
Sums over $t\in I$ run over the distinct ordinates in the union of the spectra in question. The
height needed to identify a primitive character of conductor at most $Q$, of any order, is
$O(L_Q^2/\log\log(3Q))$: this replaces the quadratic cutoff $O(L_Q^2\log\log(3Q))$ that follows from
\cref{rec:thm:quaddecode}. In this subsection $\widehat g(t)=\int g(u)e^{itu}\dd u$. For $g$ supported in a compact
subinterval of $(0,\infty)$, with $a$ the parity and $\delta_\chi=\mathbf 1_{q=1}$, the explicit formula is
\begin{equation}\label{rec:eq:halfweil}
 \sum_\gamma m_\chi(\gamma)\widehat g(\gamma)=-\sum_{n\ge2}\frac{\Lambda(n)\chi(n)}{\sqrt n}g(\log n)
 -\int_0^\infty g(u)\frac{e^{-(a+1/2)u}}{1-e^{-2u}}\dd u+\delta_\chi\int_0^\infty g(u)\bigl(e^{u/2}+e^{-u/2}\bigr)\dd u,
\end{equation}
the positive-support case of the Dirichlet explicit formula \cite{ChirrePereiraDeLaat2021Primes}, of which
\eqref{rec:eq:packetweil} is the quadratic instance: the conductor term and the reflected prime term
vanish because $g(0)=g(-\log n)=0$, and compact $C^2$ tests follow by smoothing. Every unit interval
$[v,v+1]$ contains $O(\log(3q(2+\abs v)))$ ordinates with multiplicity
\cite{BennettMartinOBryantRechnitzer2021Counting}.

\paragraph{Neighbouring prime powers.}
For a prime $p$ and $B>0$ let
\[
 W_p(B)=\sum_{\substack{n\ne p\\ \abs{\log(n/p)}\le1/2}}\frac{\Lambda(n)}{\log p}\sqrt{\frac pn}\,\min\bigl\{1,(B\log(n/p))^{-2}\bigr\},
\]
and $\mathcal P_X$ the set of primes in $[X,2X]$, $P_X=X/\log X$. A single $W_p$ can be large; their mean is
small at bandwidth $B\gg X/\log X$.

\begin{lemma}[Mean interference from neighbours; \stU]\label{rec:lem:neighbours}
Uniformly for all sufficiently large $X$ and all $B>0$,
$\frac1{P_X}\sum_{p\in\mathcal P_X}W_p(B)\ll\frac X{B\log X}$.
\end{lemma}

\begin{proof}
On the support $Xe^{-1/2}\le n\le2Xe^{1/2}$, the prefactor is bounded, and the minimum is at most
$b(h)=\min(1,(CX/(Bh))^2)$ with $h=\abs{n-p}\ge1$. For $h$ even, sieving the classes $0$ and $\mp h$
modulo primes $\ell\le y=X^{1/16}$ (one class when $\ell\mid h$) and applying the arithmetic large sieve
\cite{Montgomery1973Large} with $R=\sqrt X$ gives
\[
 \#\set{p\in\mathcal P_X:\ p\pm h\text{ prime in }[Xe^{-1/2},2Xe^{1/2}]}\ll\frac X{\log^2X}\,s(h),\qquad
 s(h)=\prod_{\substack{\ell\mid h\\\ell>2}}\frac{\ell-1}{\ell-2};
\]
indeed, by Markov's inequality at least half of the Euler product of the sieve weights is carried by
$d\le R$, and Mertens' product formula bounds its reciprocal by $s(h)/\log^2y$. Odd $h$ give no prime
pairs. Since $\sum_{h\le z}s(h)\le z\prod_{\ell>2}(1+1/(\ell(\ell-2)))\ll z$, dyadic summation gives
$\sum_hb(h)s(h)\ll X/B$, so prime neighbours contribute $O(X^2/(B\log^2X))$ to the sum over centres.
The $O(\sqrt X)$ proper prime powers in the range contribute, each summed over all integer centres,
$O(\sum_hb(h))=O(X/B)$, which after division by $P_X$ is absorbed.
\end{proof}

\paragraph{Tests on translated windows.}
For $H\ge800$ take $B=H/4$, $m=\floor{H/200}\ge4$, $w_m(u)=(1-4u^2)_+^m$ and
\[
 g_p(u)=e^{-i\tau(u-\log p)}w_m(u-\log p)\recsinc^2\bigl(B(u-\log p)\bigr),\qquad h_p=\widehat g_p,
\]
so that $g_p(\log p)=1$. Since $\int\recsinc^2(Bu)\dd u=\pi/B$,
\begin{equation}\label{rec:eq:lowlip}
 \abs{h_p(t)}\le\pi/B,\qquad\abs{h_p'(t)}\le(\log p+\tfrac12)\pi/B\qquad(t\in\R).
\end{equation}
Integrating $m$ times by parts and using Cauchy's estimate on circles of radius one for $(1-4z^2)^m$
gives $\abs{\widehat w_m(t)}\le(10m/\abs t)^m$; the factor $\recsinc^2(Bu)$ is an average of $e^{ivu}$ over
$\abs v\le2B$, so $\abs{h_p(t)}\le(20m/\abs{t-\tau})^m$ for $\abs{t-\tau}\ge H$, and dyadic summation with the
unit-interval count gives
\begin{equation}\label{rec:eq:lowtail}
 \sum_{\abs{\gamma-\tau}>H}\abs{h_p(\gamma)}\ll He^{-c_0H}\log\bigl(2q(\abs\tau+H+3)\bigr).
\end{equation}
For primitive $\chi_1,\dots,\chi_r$ of conductors at most $Q$ and $v\in\C^r$ let
$A(n)=\sum_jv_j\chi_j(n)$, $M(t)=\sum_jv_jm_{\chi_j}(t)$ and
$Z_v(p)=-\frac{\sqrt p}{\log p}\sum_{t\in I}h_p(t)M(t)$. Then \eqref{rec:eq:halfweil} gives the identity
\begin{equation}\label{rec:eq:lowobs}
 Z_v(p)=A(p)+\sum_{n\ne p}\frac{\sqrt p}{\log p}\frac{\Lambda(n)}{\sqrt n}g_p(\log n)A(n)+E_v(p),
\end{equation}
\[
 \abs{E_v(p)}\ll\frac{\norm v_1}{B\log X}+\abs{v_{\mathrm{triv}}}\frac X{B\log X}
 +\norm v_1\frac{\sqrt X}{\log X}He^{-c_0H}\log\bigl(2Q(\abs\tau+H+3)\bigr)\qquad(p\in\mathcal P_X),
\]
with $v_{\mathrm{triv}}$ the coefficient of $\zeta$ (zero if absent): on the support of $g_p$ the gamma
kernel is $O(X^{-1/2})$ and the pole kernel $O(\sqrt X)$, and $\int\abs{g_p}\le\pi/B$. Forming $Z_v(p)$ needs no
knowledge of conductors or parities.

\begin{theorem}[Linear independence in a finite window; \stG]\label{rec:thm:family}
Let $\chi_1,\dots,\chi_r$ be distinct primitive characters of conductors at most $Q$, and
$X=C_0r^2L_Q^2$ with $C_0$ a large absolute constant. For every $v\in\C^r$ and $\tau\in\R$,
\[
 H\ge C_1\bigl\{rX/\log X+J(Q,\tau)\bigr\}\ \Longrightarrow\
 \sum_{t\in I}\Bigl\lvert\sum_{j=1}^rv_jm_{\chi_j}(t)\Bigr\rvert\ \ge\ cH\frac{\log X}{\sqrt{rX}}\norm v_2 .
\]
Consequently $(m_{\chi_j}(t))_{t\in I,\,j\le r}$ has column rank $r$, and some $r$ ordinates in $I$ give an
invertible minor.
\end{theorem}

\begin{proof}
By \cref{rec:lem:gram} for a fixed smooth $0\le w\le1$ supported in $(1,2)$, with total error
$O(\sqrt XL_Q\norm v_1^2)$ and $\abs{A(p)}\le\sqrt r\norm v_2$,
\begin{equation}\label{rec:eq:lowrms}
 \sum_{p\in\mathcal P_X}\abs{A(p)}^2\gg P_X\norm v_2^2,\qquad\frac1{P_X}\sum_{p\in\mathcal P_X}\abs{A(p)}\gg\frac{\norm v_2}{\sqrt r}.
\end{equation}
Since $\abs{g_p(\log n)}\le\min\set{1,(B\log(n/p))^{-2}}$, the off-diagonal term of \eqref{rec:eq:lowobs} has mean at
most $\norm v_1P_X^{-1}\sum_pW_p(B)\ll\sqrt r\norm v_2X/(B\log X)$ by \cref{rec:lem:neighbours}, below a quarter of
\eqref{rec:eq:lowrms} for $B\ge CrX/\log X$. The first two error terms are also small, and
$H\gg rX/\log X$ together with $H\ge CJ(Q,\tau)$ absorbs the third, as
$\log(2Q(\abs\tau+H+3))\ll\exp(J(Q,\tau))+\log(H+3)$. Hence some $p$ has $\abs{Z_v(p)}\gg\norm v_2/\sqrt r$, while
\eqref{rec:eq:lowlip} gives $\abs{Z_v(p)}\ll\frac{\sqrt X}{H\log X}\sum_{t\in I}\abs{M(t)}$.
\end{proof}

\paragraph{Translated windows at a triple-logarithmic cost.}
The term $J(Q,\tau)$ comes from the polynomial decay of a compactly supported test. For a fixed number
of characters a test whose transform decays doubly exponentially in a strip lowers the cost of
translation to $bL_3(\tau)$ for every $b>1/\pi$. Such a test is not supported in $(0,\infty)$, and the
two-sided explicit formula is used: for $g$ with $g(0)=0$ whose transform $h(t)=\int g(u)e^{itu}\dd u$ is
holomorphic and rapidly decreasing in a strip wider than $\abs{\Im t}\le\tfrac12$, dividing
\eqref{rec:eq:gwstrip} by $2\pi$ and expanding the digamma term gives, under \cref{src:thm:dgrh},
\begin{multline}\label{rec:eq:fullweil}
 \sum_\gamma m_\chi(\gamma)h(\gamma)=\delta_\chi\bigl\{h(\tfrac i2)+h(-\tfrac i2)\bigr\}\\
 -\sum_{n\ge2}\frac{\Lambda(n)}{\sqrt n}\bigl\{\chi(n)g(\log n)+\overline{\chi(n)}g(-\log n)\bigr\}\\
 -\int_0^\infty\frac{e^{-(a+1/2)u}}{1-e^{-2u}}\bigl\{g(u)+g(-u)\bigr\}\dd u,
\end{multline}
the conductor term vanishing with $g(0)$ \cite{ChirrePereiraDeLaat2021Primes}. Fix $b>1/\pi$, then
$1/b<a_*<\pi$ and $\tfrac12<d<\pi/(2a_*)$, and put
\[
 K(t)=e^{-\cosh(a_*t)},\qquad w(v)=\frac{\int_\R K(t)e^{-itv}\dd t}{\int_\R K(t)\dd t}.
\]
Then $w(0)=1$, $\abs w\le1$, and $\abs{w(v)}\ll_{a_*,d}e^{-d\abs v}$, by moving the line of integration to
$\Im t=-d\operatorname{sgn}(v)$, where $\abs{K}=\exp(-\cos(a_*d)\cosh(a_*\Re t))$ and $\cos(a_*d)>0$. For a prime $p$,
$\ell=\log p$, $B\ge1$ and $\tau\in\R$ let
\[
 g_p^\natural(u)=e^{-i\tau(u-\ell)}\sin^2\Bigl(\frac{\pi u}{2\ell}\Bigr)w(u-\ell)\recsinc^4\bigl(B(u-\ell)\bigr),\qquad
 h_p^\natural=\widehat{g_p^\natural}.
\]
Thus $g_p^\natural(\ell)=1$, and the double zero at $u=0$ removes the conductor term and controls the gamma
integral for every $\tau$.

\begin{lemma}[A doubly exponential test; \stG]\label{rec:lem:coshtest}
For $p\in\mathcal P_X$, $X$ large, $B\ge1$ and every $\chi\in\mathcal C(Q)$, uniformly in $\tau$ and with constants
depending on $a_*,d$:
\begin{enumerate}
\item $\norm{h_p^\natural}_\infty\ll B^{-1}$ and $\norm{(h_p^\natural)'}_\infty\ll(1+\log p)B^{-1}$ on $\R$;
\item with $R_0=4B+\pi/\ell$ and $H\ge R_0+1$,
$\sum_{\abs{\gamma-\tau}>H}m_\chi(\gamma)\abs{h_p^\natural(\gamma)}\ll\log\bigl(2Q(\abs\tau+H+3)\bigr)\exp\bigl(-c_0e^{a_*(H-R_0)}\bigr)$
for some $c_0>0$;
\item $\bigl\lvert-\frac{\sqrt p}{\log p}\sum_\gamma m_\chi(\gamma)h_p^\natural(\gamma)-\chi(p)\bigr\rvert\ll W_p(B)+XB^{-4}+B^{-1}+\delta_\chi\frac X{B\log X}$.
\end{enumerate}
\end{lemma}

\begin{proof}
(i) follows from $\int\recsinc^4(Bv)\dd v\ll B^{-1}$, $\int\abs v\recsinc^4(Bv)\dd v\ll B^{-2}$ and
$\abs u\le\ell+\abs{u-\ell}$; the factor in front of $\recsinc^4$ is bounded. (ii) The factor $\recsinc^4(Bv)$ is
the characteristic function of a probability measure on $[-4B,4B]$, and in the variable $v=u-\ell$,
$\sin^2(\pi u/(2\ell))=\cos^2(\pi v/(2\ell))$ is a combination of three exponentials with frequencies $0,\pm\pi/\ell$
and total absolute coefficient one. By Fourier inversion of $w$, $h_p^\natural(t)$ is $e^{it\ell}$ times an average of
translates $K(t-\tau+c)$ with $\abs c\le R_0$, so $\abs{h_p^\natural(t)}\ll\sup_{\abs c\le R_0}\exp(-\cosh(a_*(t-\tau-c)))$; the
same computation in $\abs{\Im t}\le d'$, for $\tfrac12<d'<d$, gives the decay required for \eqref{rec:eq:fullweil}.
Summing over unit intervals $H+j<\abs{t-\tau}\le H+j+1$ with the count
$O(\log(2q(\abs\tau+H+j+3)))$ gives the tail. (iii) In \eqref{rec:eq:fullweil} the prime $p$ gives $\chi(p)$ after
normalization. For $\abs{\log(n/p)}\le\tfrac12$, $\abs{g_p^\natural(\log n)}\le\min\set{1,(B\log(n/p))^{-2}}$, so these
terms cost at most $W_p(B)$. For the other $n\ge2$, $\abs{g_p^\natural(\log n)}\ll B^{-4}e^{-d\abs{\log(n/p)}}$, and
$\sum_n\Lambda(n)n^{-1/2}e^{-d\abs{\log(n/p)}}\ll\sqrt p\log p$ as $d>\tfrac12$, which costs $O(pB^{-4})$; the reflected
terms have $\abs{g_p^\natural(-\log n)}\ll B^{-4}(np)^{-d}$ and cost $O(B^{-4})$. On $0<u<\tfrac14$ the double zero
gives $\abs{g_p^\natural(\pm u)}\ll u^2p^{-d}B^{-4}$ against a gamma kernel $O(1/u)$; for $u\ge\tfrac14$ the kernel is
$O(e^{-u/2})$ and $\sqrt pe^{-u/2}e^{-d\abs{u-\ell}}\le1$, so the gamma term costs $O(B^{-1})$. Finally
$\abs{h_p^\natural(\pm i/2)}\ll\int e^{\abs u/2}e^{-d\abs{u-\ell}}\recsinc^4(B(u-\ell))\dd u\ll\sqrt p/B$.
\end{proof}

\begin{theorem}[Translated windows at a triple-logarithmic cost; \stG]\label{rec:thm:familytriple}
Fix an integer $r\ge1$ and $b>1/\pi$. There are $C,c>0$ depending only on $r,b$ such that for distinct
$\chi_1,\dots,\chi_r\in\mathcal C(Q)$, every $\tau\in\R$, every
\[
 H\ge C\frac{L_Q^2}{\log\log(3Q)}+bL_3(\tau)
\]
and every $v\in\C^r$,
\[
 \sum_{t\in I}\Bigl\lvert\sum_{j=1}^rv_jm_{\chi_j}(t)\Bigr\rvert\ \ge\ cH\frac{\log\log(3Q)}{L_Q}\norm v_2 .
\]
\end{theorem}

\begin{proof}
Let $\norm v_2=1$ and $X=A_rL_Q^2$. By \cref{rec:lem:gram}, with an entrywise Gram error at most $X/(8r)$
for $A_r$ large, the weighted mean square of $A(p)=\sum_jv_j\chi_j(p)$ is at least $\tfrac78X$ while the
total weight is at most $\tfrac54X$; as each prime carries weight at most $r\log(2X)$, at least
$b_rX/\log X$ primes of $\mathcal P_X$ have $\abs{A(p)}\ge\tfrac12$, with $b_r$ independent of $A_r$. For
$B\ge K_rX/\log X$, \cref{rec:lem:neighbours} and Markov's inequality leave one of them with
$W_p(B)\le\eta_r$. Choose $\theta>0$ with $4\theta<1$ and $\lambda=a_*(1-4\theta)$ satisfying $\lambda b>1$, and put
$B=\theta H$, which is at least $K_rX/\log X$ once $C$ is large. Taking $\eta_r$ small, then $K_r$ and $A_r$ large, makes the bound of
\cref{rec:lem:coshtest}(iii), summed against $v$ with $\norm v_1\le\sqrt r$, at most $\tfrac18$. For the tail,
$a_*(H-R_0)\ge\lambda H-O(1)$, and the logarithm of the prefactor, after multiplication by $\sqrt p/\log p$, is
$O(1+\log X+\log\log(2Q)+\log\log(H+3))+e^{L_3(\tau)}$. The hypothesis gives
$e^{\lambda H}\ge e^{\lambda CL_Q^2/\log\log(3Q)}\bigl(e^{L_3(\tau)}\bigr)^{\lambda b}$ with $\lambda b>1$, which dominates this
quantity by an arbitrarily large factor for $C$ large, uniformly in $Q,\tau,H$. Hence
\begin{equation}\label{rec:eq:tripsep}
 \Bigl\lvert\sum_{t\in I}h_p^\natural(t)\sum_jv_jm_{\chi_j}(t)\Bigr\rvert\ge\frac{\log p}{4\sqrt p},
\end{equation}
and \cref{rec:lem:coshtest}(i) with $B=\theta H$ and $\log X/\sqrt X\asymp_r\log\log(3Q)/L_Q$ concludes.
\end{proof}

Compared with \cref{rec:thm:family}, the number of characters is fixed, and the window may be
centred at any height at the additive cost $bL_3(\tau)$; \cref{rec:thm:localid} has the coefficient
$1/\pi$ itself and an explicit dependence on the number of factors, but needs the larger radius
$O(\log^2(2Q))$ and a large centre. Write $M_Q=\log(2L_Q)$.

\begin{corollary}[Stability under partial matching; \stG]\label{rec:cor:lowmatch}
Let $\chi_1\ne\chi_2$ in $\mathcal C(Q)$. Match submultisets of their zero lists in $I$ bijectively, let $E$ be the
number of unmatched occurrences, $D$ the sum of the matched displacements and
$D'=\sum\min\set{1,M_Q\abs{s-t}}$ over the matched pairs $(s,t)$.
\begin{enumerate}
\item If $H\ge C\{L_Q^2/\log\log(3Q)+L_3(\tau)\}$ with $C$ a large absolute constant, then $E+D'\gg HM_Q/L_Q$.
\item If $b>1/\pi$ and $H$ is as in \cref{rec:thm:familytriple} for $r=2$, then $E+M_QD\gg_bHM_Q/L_Q$ and
$E+D'\gg_bL_Q$.
\end{enumerate}
In particular every complete matching has $D\gg H/L_Q$. The same holds for half-open $I$.
\end{corollary}

\begin{proof}
Use \eqref{rec:eq:tripsep} for $v=(1,-1)/\sqrt2$: an unmatched occurrence contributes at most
$\norm{h_p^\natural}_\infty\ll B^{-1}$ and a matched pair at most
$\min\set{2\norm{h_p^\natural}_\infty,\norm{(h_p^\natural)'}_\infty\abs{s-t}}\ll B^{-1}\min\set{1,M_Q\abs{s-t}}$, as $\log p\asymp M_Q$. With
$B=\theta H$ this gives (i), whose window is one of \cref{rec:thm:familytriple} for $b=1$, and the first
bound of (ii). For the second bound of (ii) take instead $B=K_2X/\log X$, independent of $H$: then
$H-R_0\ge\tfrac12CL_Q^2/\log\log(3Q)+bL_3(\tau)-1$ for $C$ large, the tail estimate of the proof above is
unchanged, and $E+D'\gg B\log X/\sqrt X\asymp\sqrt X\asymp L_Q$. Unmatched endpoints of $I$ are covered by the
tail bound.
\end{proof}

The same construction reads every prime value at once on a common set of primes.

\begin{proposition}[A common set of readable primes; \stG]\label{rec:prop:read}
Let $X=AL_Q^2$ with $A$ a large absolute constant, and $H\ge B_1\{X/\log X+J(Q,\tau)\}$ with $B_1$ large.
There is a set $\mathcal P\subseteq\mathcal P_X$, independent of the character, such that with
$Z(p;\chi)=-\frac{\sqrt p}{\log p}\sum_{t\in I}h_p(t)m_\chi(t)$,
\[
 \abs{Z(p;\chi)-\chi(p)}\le\tfrac18\qquad(\chi\in\mathcal C(Q),\ p\in\mathcal P),
\]
also with either endpoint of $I$ omitted, and every pair of distinct $\chi,\psi\in\mathcal C(Q)$ has at least
$bX/\log X$ primes $p\in\mathcal P$ with $\abs{\chi(p)-\psi(p)}\ge1$, and, if $\psi\notin\set{\chi,\bar\chi}$, at least
$bX/\log X$ with $\abs{\Re\chi(p)-\Re\psi(p)}\ge\tfrac12$, for an absolute $b>0$. For every $b'>1/\pi$ the same
holds, with $\mathcal P$ and $B_1$ depending on $b'$, when $Z(p;\chi)$ is formed with $h_p^\natural$ for
$B=K_{b'}X/\log X$ and $H\ge B_1X/\log X+b'L_3(\tau)$.
\end{proposition}

\begin{proof}
On $(\Z/M\Z)^\times$, $M=\lcm(q_\chi,q_\psi)\le Q^2$, the characters remain distinct; with weights
$w_p=(\log p)W(p/X)$, $W$ a fixed smooth weight with $0\le W\le1$ supported in $(1,2)$ and $\int W>0$, \cref{rec:lem:gram} applied to the monomials of $\abs{\chi-\psi}^2$ (mean $2$ on the group)
and of $(\Re\chi-\Re\psi)^2$ (mean at least $1$ for nonconjugate pairs, the real parts being orthogonal),
each of total absolute coefficient at most $4$, shows that the prime-weighted means differ from the
group means by at most $\tfrac14$ for $A$ large, uniformly in $Q$; as the squares are at most $4$ and
$w_p\le\log(2X)$, at least $2bX/\log X$ primes witness each pair. By \eqref{rec:eq:lowobs} with $v=e_\chi$,
$\abs{Z(p;\chi)-\chi(p)}\le W_p(B)+\abs{E(p)}$, and the error is at most $1/16$ for $B_1$ large. By
\cref{rec:lem:neighbours} and Markov's inequality, all but $bX/\log X$ primes of $\mathcal P_X$ have
$W_p(H/4)\le1/16$; let $\mathcal P$ be these. The tail bound covers either endpoint. In the last
statement \cref{rec:lem:coshtest}(iii) replaces \eqref{rec:eq:lowobs}, $\mathcal P$ consists of the primes with
$W_p(K_{b'}X/\log X)$ below a fixed small threshold, which depends only on prime powers, and the tail is
controlled as in the second part of the proof of \cref{rec:cor:lowmatch}; the reading error is then
below $\tfrac18$ simultaneously for all $\chi\in\mathcal C(Q)$.
\end{proof}

\begin{theorem}[Identification from low zeros; \stG]\label{rec:thm:lowzero}
There are absolute $C,c>0$ such that, with $H_Q=L_Q^2/\log\log(3Q)$:
\begin{enumerate}
\item the multiset of ordinates $0<\gamma\le CH_Q$, with multiplicity, determines $\chi\in\mathcal C(Q)$;
\item for distinct members, even after matching occurrences in these lists at distance at most
$c/L_Q^2$, at least $cL_Q$ occurrences remain unmatched;
\item the multiset of $\abs\gamma$ over all signed ordinates with $\abs\gamma\le CH_Q$ determines the pair
$\set{\chi,\bar\chi}$, and conjugate characters have the same folded data at every height;
\item for each $b>1/\pi$ there is $C_b$ such that for every $\tau\in\R$ the multiset of ordinates in
$[\tau-H,\tau+H]$, $H\ge C_bH_Q+bL_3(\tau)$, determines $\chi\in\mathcal C(Q)$.
\end{enumerate}
\end{theorem}

\begin{proof}
Take $H\asymp H_Q$ with a large lower constant, $\tau=H$, and $I$ with its left endpoint omitted, which
is $(0,2H]$; then $H\ge C\{L_Q^2/\log\log(3Q)+L_3(\tau)\}$, as $L_3(H)=o(H)$. By \cref{rec:cor:lowmatch},
$E+M_QD\gg HM_Q/L_Q\gg L_Q$, and $D$ is absorbed: each list has $O(HL_Q)$ occurrences, so displacements
at most $c/L_Q^2$ give $D\ll cH/L_Q$, and $M_QD$ is a small multiple of $HM_Q/L_Q$ for small $c$. This proves
(i) and (ii), and (iv) is \cref{rec:thm:familytriple} for $r=2$. For (iii) take $\tau=0$: the test is then real and even
about $\log p$, so $\Re h_p(\gamma)=\cos(\gamma\log p)\widehat{g_0}(\gamma)$ is even in $\gamma$ and $\Re Z(p;\chi)$ is computed
from the folded list. For nonconjugate pairs a prime of \cref{rec:prop:read} has
$\abs{\Re\chi(p)-\Re\psi(p)}\ge\tfrac12>2\cdot\tfrac18$. Finally $L(\bar s,\bar\chi)=\overline{L(s,\chi)}$ reflects the signed list.
\end{proof}

A list of positive ordinates keeps only the zeros above the real axis, while a folded list
superimposes both half-planes; for a complex character they carry different information. The
discrepancy in (ii) may occur at an ordinate shared with different multiplicities. The character,
its conductor, parity and order are recovered, and neither conductor nor parity is needed to form the
measurements. For the quartic character $\chi$ modulo $5$ with $\chi(2)=i$, a numerical sign-change scan
gives the first positive ordinates $6.183578195451$, $8.457229174423$, $12.67494641701$ for $\chi$ and
$4.132903705213$, $9.442931129729$, $11.28289644158$ for $\bar\chi$, so the signed lists already differ in
$[-5,5]$ while the folded lists coincide. In the window $[0,72]$, a decoder of the type of
\cref{rec:lem:coshtest}, with $u/\ell$ in place of $\sin^2(\pi u/(2\ell))$ and $p=2$, $b=1$, $B=8$, applied to the $35$ ordinates found for each character
returns $\chi(2)$ and $\bar\chi(2)$ with errors below $0.00033$. These lists are not certified complete:
\begin{center}
\artifact{dirichlet-zero-reconstruction/checks/complex_zeros.txt}\\
\artifact{dirichlet-zero-reconstruction/checks/decoder_demo.txt}
\end{center}

\begin{corollary}[A dictionary of $O(\log Q)$ primes; \stG]\label{rec:cor:dict}
There are $k=O(L_Q)$ primes in $\mathcal P$, chosen independently of the unknown character, and an
absolute $d>0$, such that every distinct pair in $\mathcal C(Q)$ has $\abs{\chi(p)-\psi(p)}\ge1$ on at least $dk$ of
them, and every nonconjugate pair has $\abs{\Re\chi(p)-\Re\psi(p)}\ge\tfrac12$ on at least $dk$. A decoder
choosing a candidate within $\tfrac18$ on the most coordinates succeeds with arbitrary corruption of
fewer than $dk/2$ readings if the others have error at most $\tfrac18$; the same primes serve every
window of \cref{rec:prop:read}. For dictionaries of exact prime evaluations the order of $k$ cannot be
reduced, even for quadratic characters, nor for any scheme of $k$ complex measurements of modulus at
most $2$ that is uniquely decodable under errors $\tfrac18$ in every coordinate, nor for $k$ symbols from a
fixed finite alphabet.
\end{corollary}

\begin{proof}
A fraction $\beta>0$ of $\mathcal P$ witnesses each pair. Choose $k=\ceil{C_2L_Q}$ primes uniformly without
replacement; a hypergeometric tail bound gives probability at most $e^{-\beta k/8}$ of fewer than $\beta k/2$
witnesses for a given pair, and a union bound over at most $2Q^4$ pairs (there are at most $Q^2$
characters) gives existence. On a witnessing coordinate the two radius-$\tfrac18$ disks are disjoint.
For the lower bound, the $\gg Q/\log Q$ prime conductors in $(Q/2,Q]$ give distinct quadratic characters
and each prime evaluation takes at most three values, so $3^k\gg Q/\log Q$. A prime $q\in(Q/2,Q]$ gives
$q-2\gg Q$ primitive characters; if $[-2,2]^2$ is cut into $D_0$ cells of diameter at most $\tfrac18$, two
characters whose measurement vectors lie in the same product cell can produce the same observation,
so $D_0^k\gg Q$, and an alphabet of $s\ge2$ symbols needs $s^k\gg Q$.
\end{proof}

\begin{example}\label{rec:ex:dict101}
For $q=101$, $X=10^4$, $R=8\cdot10^4$ and the test half-width $\tfrac18$, the primes $11197$, $14107$, $11027$
(indices $60,32,39$ to the base $2$) separate all $100$ characters modulo $101$ by chords of length at
least one, with certified leakage bounds below $0.05$ obtained by enumerating every prime power in an
enlarged window. For the quartic character $\chi$ modulo $5$ with $\chi(2)=i$, interval arithmetic with
Euler--Maclaurin remainders and Rouch\'e's theorem shows unconditionally that the disk of radius
$10^{-4}$ about $\tfrac12+4.13290i$ contains exactly one simple zero of $L(s,\bar\chi)$, which lies on the
critical line by reflection, and no zero of $L(s,\chi)$. Symmetrically, the disk of radius $10^{-8}$ about
$\tfrac12+6.18357819545i$ contains exactly one simple zero of $L(s,\chi)$, while $\abs{L(s,\bar\chi)}>2$ on the
closed disk: $\abs{L(s_0,\chi)}<10^{-11}$ and $\abs{L'(s_0,\chi)}>1$ at the centre, and Cauchy's inequalities on
the circle of radius $\tfrac14$, where $\abs L<64$, bound the remaining Taylor terms. For the sieve step,
among the $5244$ primes in $(60000,120000)$ exactly $3598$ have certified near-prime-power leakage at most
$\tfrac18$ at $B=45000$, every nontrivial character difference modulo $127$ is separated by at least $1818$ of
them, and $61703$, $60013$, $60029$ separate all $125$ nonprincipal characters modulo $127$, in integer
arithmetic:
\begin{center}
\artifact{dirichlet_low_zero_identification/checks/prime_dictionary.txt}\\
\artifact{dirichlet_low_zero_identification/checks/quartic_certificate.txt}\\
\artifact{dirichlet_zero_windows/example/quartic-output.txt}\\
\artifact{dirichlet-zero-reconstruction/checks/prime_probes.txt}
\end{center}
\end{example}

\begin{corollary}[Products of $L$-functions; \stG]\label{rec:cor:products}
Fix $k\ge1$. A product $\prod_jL(s,\chi_j)^{e_j}$ with at most $k$ distinct factors $\chi_j\in\mathcal C(Q)$ and
arbitrary positive integer exponents is determined, with its factors and exponents, by the pooled
multiset of its positive zero ordinates up to $O_k(L_Q^2/\log\log(3Q))$, without labels, and also by the
pooled ordinates in any window $I(\tau,H)$ with $H\ge C_k\{L_Q^2/\log\log(3Q)+L_3(\tau)\}$. More generally, the
same windows determine the support and all nonzero complex coefficients of any combination
$\sum_jz_jm_{\chi_j}$ of at most $k$ distinct primitive zero measures.
\end{corollary}

\begin{proof}
Orders add under multiplication at every point of the open strip. Two different products, or
combinations, give a nonzero vector $v$ of coefficient differences on at most $2k$ distinct characters,
the empty product included; apply \cref{rec:thm:familytriple} with $b=1$ and $r\le2k$, taking the largest
of the finitely many constants, and for positive ordinates the window $[H,3H]$, $\tau=2H$.
\end{proof}

\paragraph{Abelian fields.}
For a finite abelian $K/\Q$ let $\mathcal X_K$ be its group of Dirichlet characters, each represented by
its primitive character. Then $\zeta_K(s)=\prod_{\chi\in\mathcal X_K}L(s,\chi)$ and
$\abs{\operatorname{disc}K}=\prod_{\chi\in\mathcal X_K}q_\chi$, and inside a common cyclotomic field $\mathcal X_K$ determines $K$, its
annihilator being the subgroup fixing $K$ \cite{Washington1997Introduction}. Write $m_K(t)$ for the
multiplicity of $\zeta_K$ at $\tfrac12+it$.

\begin{theorem}[Abelian fields of unrestricted degree; \stG]\label{rec:thm:fields}
There are absolute $C,c>0$ such that for distinct finite abelian $K,L\subset\overline\Q$ with absolute
discriminants at most $D\ge3$, $r$ the size of the symmetric difference of $\mathcal X_K$ and $\mathcal X_L$, every
$\tau\in\R$ and $H\ge C\{\log^2(2D)+J(D,\tau)\}$,
\[
 \sum_{t\in I}\abs{m_K(t)-m_L(t)}\ge cH\sqrt r\,\frac{\log\log(3D)}{\log(2D)} .
\]
In particular the zeros with $\abs t\le C'\log^2(2D)$, or the positive ones up to that height,
determine an abelian field of absolute discriminant at most $D$, independently of its degree.
\end{theorem}

\begin{proof}
After cancelling common factors, $M(t)=m_K(t)-m_L(t)=\sum_{j\le r}v_jm_{\chi_j}(t)$ with $v_j=\pm1$, the
trivial character cancelling, $S_0=\sum_j\log(2q_j)\ll\log(2D)$ and $r\ll\log(2D)$, since each $q_j\ge3$ and
$\sum\log q_j\le2\log D$. Take $X=C_2\log^2(2D)$. Summing the errors of \cref{rec:lem:gram} over all pairs
gives $O(\sqrt XrS_0)$, below half the diagonal $rX\int w$; with a second weight majorizing the
indicator of $\mathcal N_X=[e^{-1/2}X,2e^{1/2}X]$,
$\sum_{p\in\mathcal P_X}\abs{A(p)}^2\gg rP_X$ and $\sum_{p\in\mathcal N_X}\abs{A(p)}^2\ll rP_X$. The vector
$b_n=(\Lambda(n)/\log X)\sqrt{X/n}\,A(n)$ on $\mathcal N_X$ has $\norm b_2^2\ll rP_X+r^2\sqrt X/\log X\ll rP_X$, using
$\abs{A(n)}\le r$ at proper prime powers. The off-diagonal matrix
$T_{pn}=\mathbf 1_{p\ne n}\sqrt{p/X}(\log X/\log p)g_p(\log n)$ has $\abs{T_{pn}}\ll(X/B)^2\abs{n-p}^{-2}$, so its row and
column sums are $O((X/B)^2)$ and Schur's test gives $\norm{Tb}_2\ll(X/B)^2\norm b_2$, a quarter of the
central norm for $B\ge K_0X$. With $v_{\mathrm{triv}}=0$, $\norm v_1=r\ll\sqrt X$ and $H\ge C(X+J(D,\tau))$ the
errors are another quarter, so $\sum_{p\in\mathcal P_X}\abs{Z_v(p)}^2\gg rP_X$, some $\abs{Z_v(p)}\gg\sqrt r$, and
\eqref{rec:eq:lowlip} concludes. Positive ordinates suffice by the window $[H,3H]$.
\end{proof}

For bounded degree the height drops by a factor $\log\log D$.

\begin{corollary}[Abelian fields of bounded degree; \stG]\label{rec:cor:fieldsbounded}
Fix $n\ge1$. Among abelian $K\subset\overline\Q$ of degree at most $n$ and conductor at most $Q$, the zero
multiset of $\zeta_K$ in any window $I(\tau,H)$ with $H\ge C_n\{L_Q^2/\log\log(3Q)+L_3(\tau)\}$, or its positive
ordinates up to $C_n'L_Q^2/\log\log(3Q)$, determines $K$. The same holds with an upper bound $Q$ on the
absolute discriminant in place of the conductor.
\end{corollary}

\begin{proof}
The characters of $\mathcal X_K$ are $[K:\Q]\le n$ primitive characters whose conductors divide the
conductor of $K$, and each $q_\chi$ divides $\abs{\disc K}=\prod_{\chi\in\mathcal X_K}q_\chi$. By
\cref{rec:cor:products} with $k=n$ the data determine $\mathcal X_K$, which determines $K$.
\end{proof}

For example, $\Q(\zeta_5)$ and $\Q(\sqrt5,\sqrt{-3})$ share $\zeta$ and the quadratic character of
conductor $5$; the remaining constituents are the conjugate quartic characters modulo $5$, respectively
the quadratic characters of conductors $3$ and $15$, and the corollary separates the two zeta zero lists
although the zeros of the common factors occur in both.

\begin{example}[Two biquadratic fields]\label{rec:ex:biquad}
$K=\Q(\sqrt5,\sqrt{221})$ and $L=\Q(\sqrt{13},\sqrt{85})$ are totally real with discriminant $1105^2$
and share $\zeta$ and $L(s,\chi_{1105})$, while
$\chi_5(3)+\chi_{221}(3)-\chi_{13}(3)-\chi_{85}(3)=-4$. Unconditionally, $\zeta_K(\tfrac12+it)$ vanishes for some
$1.20<t<1.23$, while $\zeta_L(\tfrac12+it)\ne0$ on $[1.20,1.23]$: outward interval enclosures of the
theta integral $F_d(t)=4\int_0^\infty e^{u/2}\cos(tu)\sum_n\chi_d(n)e^{-\pi n^2e^{2u}/d}\dd u$, which differs from
$L(\tfrac12+it,\chi_d)$ by nonzero factors, give a sign change for $d=221$ and positivity for
$d=13,85,1105$, and Euler--Maclaurin gives $\Im\zeta(\tfrac12+it)<-0.41$ there. Under \cref{src:thm:dgrh},
$\sum_{0<t\le80}\abs{m_K(t)-m_L(t)}\ge7$, from the test $g(u)=\phi((u-\log3)/a)$, $a=2/7$,
$\phi(x)=(x(1-x^2)_+^7)'$, whose support meets only $\log3$, with $\abs{\widehat g(t)}\le C_6\abs t^{-6}$,
$C_6\le2629707255/16$, and the explicit zero count of
\cite[Corollary~1.2]{BennettMartinOBryantRechnitzer2021Counting} for the tail:
\begin{center}
\artifact{dirichlet_zeros_and_abelian_fields/checks/theta_example.txt}\\
\artifact{dirichlet_zeros_and_abelian_fields/checks/field_example.txt}
\end{center}
\end{example}

\paragraph{Explicit constants and a clipped matching metric.}
For finite real multisets $A,B$ let $d_{\mathrm{edit}}(A,B)$ be the least cost of a partial matching, an
unmatched occurrence costing $1$ and a matched pair $\min\set{1,\abs{a-b}}$, and let $\Gamma_\chi(H)$ be the
signed ordinates of $L(s,\chi)$ with $\abs\gamma\le H$, with multiplicity.

\begin{lemma}[The comparison metric; \stU]\label{rec:lem:edit}
$d_{\mathrm{edit}}$ is a metric on finite real multisets, an optimal matching pairs all occurrences of the
smaller one, and for $w$ with $\norm w_\infty\le S$ and Lipschitz constant at most $L$,
$\bigl\lvert\sum_{a\in A}w(a)-\sum_{b\in B}w(b)\bigr\rvert\le\max\set{2S,L}\,d_{\mathrm{edit}}(A,B)$.
\end{lemma}

\begin{proof}
Pairing two unmatched occurrences costs at most $1$ instead of $2$. For the triangle inequality
compose matchings; where an occurrence of the middle multiset is matched on one side only, delete
the surviving endpoint, whose cost $1$ is bounded by the movement plus the deletion already
charged. A matched pair contributes at most $\min\set{2S,L\abs{a-b}}$.
\end{proof}

\begin{lemma}[Explicit coarse counts and an angular prime; \stG]\label{rec:lem:angular}
For primitive $\chi$ of conductor $q>1$ and $T\ge2$, $N_\chi(T)\le T\log(qT)$, and $N_\zeta(T)\le2T\log(2T)$; so
$N_\chi(T)\le2T\log(2qT)$ for all primitive $\chi$. If $M\ge3$, $\psi$ is a nonprincipal character modulo $M$ and
$x=10^6\log^2(2M)$, some prime $p\nmid M$ with $x/e\le p\le ex$ has $\Re\psi(p)\le\tfrac12$, hence
$\abs{1-\psi(p)}\ge1$. Consequently two distinct $\chi_1,\chi_2\in\mathcal C(Q)$ satisfy $\abs{\chi_1(p)-\chi_2(p)}\ge1$ at a
prime $p<12\cdot10^6\log^2(2Q)$.
\end{lemma}

\begin{proof}
By \cref{src:thm:dgrh} and the functional equation,
$\Re(\Lambda_\chi'/\Lambda_\chi)(\tfrac12+u)=\sum_\gamma u/(u^2+\gamma^2)$ for $u>0$, and each zero with $\abs\gamma\le T$ contributes
at least $1/(2T)$ at $u=T$. At $\sigma=T+\tfrac12$, $\abs{L'/L(\sigma,\chi)}<\tfrac12$ and $\psi_\Gamma(y)\le\log y$ give
$\Re(\Lambda_\chi'/\Lambda_\chi)(T+\tfrac12)<\tfrac12\log(qT)$, using $7e<8\pi$; for $\xi$, $\zeta'/\zeta(\sigma)<0$ gives
$\xi'/\xi(T+\tfrac12)\le\log(2T)$. For the prime, apply \eqref{rec:eq:halfweil} to the triangle
$g(v)=(1-\abs{v-\log x})_+$, weighted sums $W_\chi(x)=\sum\Lambda(n)\chi(n)n^{-1/2}g(\log n)$: the gamma term is below
$2/\sqrt x$, the zero sum at most $B(q)=20\log(2q)+21$ by the count, the pole term for $\zeta$ is
$c(\sqrt x+x^{-1/2})$ with $c=(\sinh(1/4)/(1/4))^2>1$, removing ramified Euler factors costs less than
$\tfrac52\log M$, and prime powers of exponent at least two at most $\tfrac12(1+\log x)^2+4$. If
$\Re\psi(p)>\tfrac12$ throughout the support, positivity gives $\Re W_\psi\ge\tfrac12W_{\chi_{0,M}}-\tfrac32E_{\mathrm{pp}}$, hence
$c\sqrt x\le47.5\log(2M)+89+1.5(1+\log x)^2<1000\log(2M)$, a contradiction. For the last claim apply this
to $\chi_1\overline{\chi_2}$ modulo $M=\lcm(q_1,q_2)\le Q^2$.
\end{proof}

The coarse counts bound the zeros from above. In the other direction, every window, at every centre,
contains explicitly many zeros; for this statement only, $N_\chi(H;T)$ is the number of zeros with $\abs{\Im\rho-H}<T$, with
multiplicity.

\begin{proposition}[Zeros in a moving window; \stG]\label{rec:prop:movingwin}
Let $\chi$ have conductor $q>1$, $H\in\R$, $\Delta\ge1$ and $T\Delta\ge\tfrac12$. Then
\[
 N_\chi(H;T)\ge\frac\pi{16\Delta}\Bigl(1-\frac1{4\Delta^2T^2}\Bigr)\log\frac{q(2+\abs H)}\pi-\frac{32}\Delta-\frac{32e^{\pi\Delta}}{\Delta^2}.
\]
In particular, with $U=\log(q(2+\abs H)/\pi)$, $\Delta=(\log U-6)/\pi\ge2$ and $T=1/(2\Delta)+1/\Delta^2$, every window of half-width
$T=\pi/(2\log U)+O((\log U)^{-2})$ contains at least $\pi U/(32\Delta^2)-32/\Delta$ zeros, uniformly in $H$; and at $H=0$, with
$\Delta=\max(1,1/T)$, the window contains $M$ zeros as soon as $\log(2q/\pi)\ge\frac{64}{3\pi}(M\Delta+32+32e^{\pi\Delta}/\Delta)$.
\end{proposition}

\begin{proof}
Pass to the primitive inducer, and use $\widehat f(v)=\int f(u)e^{-2\pi iuv}\dd u$. Let $g$ be the inverse transform of
the probability density $k(v)=\tfrac\pi2\cos(\pi v)\mathbf 1_{[-1/2,1/2]}(v)$, that is $g(u)=\cos(\pi u)/(1-4u^2)$, so $\abs g\le1$, and
\[
 f(t)=\Bigl(1-\frac{(t-H)^2}{T^2}\Bigr)g\bigl(\Delta(t-H)\bigr)^2\le\mathbf 1_{\abs{t-H}<T}.
\]
It is entire, decays quadratically in horizontal strips, and with $K=k*k$, $b=\Delta T$,
\[
 \widehat f(v)=\frac{e^{-2\pi iHv}}\Delta\Bigl(K(v/\Delta)+\frac{K''(v/\Delta)}{4\pi^2b^2}\Bigr),
\]
supported in $[-\Delta,\Delta]$, with $K''=k'*k'$, $K(0)=\pi^2/8$, $K''(0)=-\pi^4/8$, so $\widehat f(0)=\frac{\pi^2}{8\Delta}(1-\frac1{4b^2})$; near the edge,
with $s=1-v/\Delta$, $\abs K\le\pi^4s^3/24$ and $\abs{K''}\le\pi^4s/4$, hence $\abs{\widehat f(v)}\le7(\Delta-v)/\Delta^2$ for $0\le v\le\Delta$. By
\cref{src:thm:dgrh} the Guinand--Weil formula \cite{CarneiroChandeeMilinovich2015Note} gives $\sum_\gamma f(\gamma)$ as
$(2\pi)^{-1}\widehat f(0)\log(q/\pi)$ plus the archimedean term $(2\pi)^{-1}\int f(t)\Re\psi(\sigma+it/2)\dd t$, $\sigma\in\set{\tfrac14,\tfrac34}$, minus
prime terms of absolute value at most $\pi^{-1}\sum_{n\le e^{2\pi\Delta}}\Lambda(n)n^{-1/2}\abs{\widehat f(\log n/(2\pi))}$. From $\sum_{n\le x}\Lambda(n)\le4x$ and
a split at $e^{2\pi\Delta-j}$, $\sum_{n\le X}\Lambda(n)n^{-1/2}\log(X/n)\le4\sqrt e(1-e^{-1/2})^{-2}\sqrt X$, and the prime terms are at most
$32e^{\pi\Delta}/\Delta^2$. For the archimedean term, $\abs{\psi(z)-\log z}\le1/\Re z$ for $\Re z>0$, $\tfrac1{12}(2+\abs t)\le\abs{\sigma+it/2}\le2+\abs t$,
and $2+\abs t$, $2+\abs H$ are within a factor $1+\abs u/2$, $u=\Delta(t-H)$, so
$\abs{\Re\psi(\sigma+it/2)-\log(2+\abs H)}\le7+\log(1+\abs u/2)$; with $\abs{(1-u^2/b^2)g(u)^2}\le5$ for $\abs u\le1$ and $\le5/(9u^2)$ beyond,
the term is $(2\pi)^{-1}\widehat f(0)\log(2+\abs H)$ up to $32/\Delta$. Since $\sum_\gamma f(\gamma)\le N_\chi(H;T)$ the bound follows. For the moving
window, $\Delta\ge2$ gives $1-1/(4\Delta^2T^2)\ge1/\Delta$ and $32e^{\pi\Delta}/\Delta^2=32e^{-6}U/\Delta^2<\pi U/(32\Delta^2)$; at $H=0$ the factor in
brackets is at least $\tfrac34$.
\end{proof}

The centred shrinking-window phenomenon, with a sharper second-order constant, is due to Zhao
\cite[Theorems~1 and~8]{Zhao2025Arxiv}; the proposition adds explicit constants, every centre and a free
bandwidth. For $T\le1$ the bandwidth $\Delta=(1+T)/(2T)$ lowers the exponential scale of the conductor threshold to
$e^{\pi/(2T)}$ times factors polynomial in $1/T$. Bennett, Martin, O'Bryant and Rechnitzer note that their
qualitative Conjecture~1.3 follows from Selberg's work under Dirichlet GRH
\cite{BennettMartinOBryantRechnitzer2021Counting}; their Conjecture~1.4 is not addressed here.

\begin{theorem}[Comparison from a separating prime; \stG]\label{rec:thm:interface}
Let $X\ge2$, $A,B\ge1$, suppose $N_\chi(t)\le A(t+1)\log(B(t+2))$ for $t\ge2$ for the characters considered,
and put $h=1/(64X)$, $D=1+\log(B(256X+2))$, $Z=80A\sqrt XD$, with $m\ge8$ even and $m\ge\frac2{\log3}\log Z$,
$H_0=256mX$. If $\abs{\chi_1(p)-\chi_2(p)}\ge1$ at a prime $p\le X$, then for every $H\ge H_0$,
\[
 d_{\mathrm{edit}}\bigl(\Gamma_{\chi_1}(H),\Gamma_{\chi_2}(H)\bigr)\ge\frac{5\log2}{128\pi}\frac H{\sqrt X}.
\]
\end{theorem}

\begin{proof}
Let $f$ be the density of a sum of $m$ independent uniforms on $[-h/m,h/m]$, so
$\widehat f(t)=\recsinc(ht/m)^m$ and $f(0)\ge1/h$ (write $f=g*g$), and put
$\phi_H(u)=\frac{f(u)}{f(0)}\recsinc(Hu/8)^2$, $K_H=\widehat{\phi_H}$, $w_{p,H}(t)=\frac{\sqrt p}{\log p}e^{it\log p}K_H(t)$. Then
$\phi_H(0)=1$, $\int\phi_H\le\min(h,8\pi/H)$, and the support of the frequency test meets the logarithm of
one positive integer only, $p$. By \eqref{rec:eq:halfweil},
$\chi(p)=\delta_\chi P_H(p)-G_{a,H}(p)-\sum_\gamma w_{p,H}(\gamma)$, with $0\le G_{a,H}\le4h/\log2\le\tfrac1{16}$ and
$0\le P_H(p)\le32\pi X/H<\tfrac1{16}$. As $\recsinc(Hu/8)^2$ is the characteristic function of a triangular
density on $[-H/4,H/4]$, $\abs{K_H(t)}\le h(4m/(3h\abs t))^m$ for $\abs t\ge H$, and Stieltjes summation against
the count gives the tail $E(H)$, decreasing in $H$, with
$E(H_0)\le10A\sqrt XDm^23^{-m}\le\tfrac18$, since $m^2\le3^{m/2}$ and $3^{m/2}\ge Z$. So the two zero sums over
$\Gamma_{\chi_j}(H)$ differ by at least $\tfrac58$, while $\max\set{2\norm{w_{p,H}}_\infty,\operatorname{Lip}w_{p,H}}\le16\pi\sqrt X/(H\log2)$;
\cref{rec:lem:edit} concludes.
\end{proof}

\begin{theorem}[Robust identification with explicit constants; \stG]\label{rec:thm:robust}
Put $X_Q=\ceil{12\cdot10^6L_Q^2}$, $D_Q=1+\log(2Q(256X_Q+2))$,
$m_Q=2\ceil{\tfrac12\max\set{8,\frac2{\log3}\log(160\sqrt{X_Q}D_Q)}}$ and $H_Q'=256m_QX_Q=O(L_Q^2\log\log(3Q))$. For
$H\ge H_Q'$ and distinct $\chi_1,\chi_2\in\mathcal C(Q)$,
\[
 d_{\mathrm{edit}}\bigl(\Gamma_{\chi_1}(H),\Gamma_{\chi_2}(H)\bigr)\ge\Delta_Q(H)=\frac{5\log2}{128\pi}\frac H{\sqrt{X_Q}}>\frac H{10^6\log(2Q)} .
\]
Hence a finite multiset within $d_{\mathrm{edit}}$-distance $\Delta_Q(H)/2$ of some spectrum, for example after
$r$ erasures, $s$ insertions and movements $e_j$ with $r+s+\sum_j\min\set{1,\abs{e_j}}<\Delta_Q(H)/2$, identifies
the character, and $\abs{\Gamma_{\chi_1}(H)}+\abs{\Gamma_{\chi_2}(H)}-2\abs{\Gamma_{\chi_1}(H)\cap\Gamma_{\chi_2}(H)}\ge\Delta_Q(H)$.
\end{theorem}

\begin{proof}
\Cref{rec:lem:angular} gives the separating prime $p\le X_Q$ and the count with $A=2$, $B=2Q$; apply
\cref{rec:thm:interface} and $5\log2/(128\pi\sqrt{13\cdot10^6})>10^{-6}$. The rest is the triangle inequality.
\end{proof}

The readings $a_p(\widetilde\Gamma)=-\sum_{t\in\widetilde\Gamma}w_{p,H}(t)$, $p\le X_Q$, of a list $\widetilde\Gamma$ with corruption cost
at most $\eta$ satisfy $\abs{a_p-\chi(p)}\le\tfrac3{16}+16\pi\sqrt{X_Q}\eta/(H\log2)$, so retaining the candidates
compatible at every $p\le X_Q$ is a finite decoder needing neither conductor nor parity. The oriented
phase is necessary, since $\Gamma_{\bar\chi}(H)=-\Gamma_\chi(H)$. For the quartic character modulo $5$ with
$\chi(2)=i$, its conjugate and its square, numerical lists to height $40$ and a damaged list (two
erasures, three insertions, movements below $0.012$) are nearest to the true character with gap
$16.91$ (\artifact{robust_dirichlet_spectra_2026-11-26/demo/demo_results.json}; numerical, not certified
complete). Since $\min\set{1,M_Qx}\le M_Q\min\set{1,x}$, \cref{rec:cor:lowmatch}(i) gives
$d_{\mathrm{edit}}\gg H/L_Q$ already for $H\ge C\{L_Q^2/\log\log(3Q)+L_3(\tau)\}$, with an unspecified constant;
\cref{rec:thm:robust} is the version with explicit constants, at a larger height. For quadratic
characters \cref{rec:thm:quaddisplace} gives the unclipped total displacement with the smaller
constant $1/(144\pi)$ from the somewhat lower height $2T$.

\subsection{Twist multisets from balanced zero statistics}\label{rec:sec:homometry}
\sources{dirichlet_zero_packet_homometry; dirichlet_packet_reconstruction; dirichlet-zero-reconstruction}

A \emph{packet} is a finite nonempty multiset of characters; pooling the zeros of its $L$-functions,
let $Z_A(T)$ be the multiset of ordinates in $(0,T]$. An observer who sees only differences of
heights can measure, for $d\ge2$ and real $t_j\ne0$ with $\sum t_j=0$,
\[
 S_A(t;T)=\frac1T\sum_{\gamma\in Z_A(T)}e^{it\gamma},\qquad \mathcal B_A(\mathbf t;T)=\prod_{j=1}^dS_A(t_j;T),
\]
since a common translation of all heights cancels. Fix a modulus $Q$, $G=(\Z/Q\Z)^\times$, a character
$\psi$ modulo $Q$ and a subgroup $H\le\widehat G$ of exponent $E$; a function $A:H\to\Z_{\ge0}$ denotes the
packet with $A(h)$ copies of $\psi h$, all factors being primitive of conductor $Q$. With
$\zeta=e^{2\pi i/E}$ and $m_A=\sum_hA(h)$, put
$f_{A,a}(u)=\sum_{h:h(a)=\zeta^u}A(h)$ and $F_A(a,k)=\sum_uf_{A,a}(u)\zeta^{ku}$ for $a\in G$. By the
Landau formula of Ford and Zaharescu \cite[Lemma~2.1]{FordZaharescu2015Unnormalized} and
\cref{src:thm:dgrh}, $S_A(t;T)=L_A(t)+O_{A,t}(\log(2T)/T)$, where
\begin{equation}\label{rec:eq:landau}
 L_A(k\log p)=-\frac{\log p}{2\pi p^{\abs k/2}}\psi(p)^kF_A(p,k)\qquad(p\nmid Q,\ k\in\Z\setminus\set0)
\end{equation}
and $L_A(t)=0$ at all other nonzero frequencies.

\begin{theorem}[Local classification; \stG]\label{rec:thm:localclass}
The following are equivalent: \textup{(i)} the limits of all balanced measurements agree for $A,B$;
\textup{(ii)} they agree for the finitely many tuples at a single prime $p$, of orders
$2\le d\le\max(2,E)$, one prime in each reduced class; \textup{(iii)} for each $a\in G$, $f_{B,a}$ is a
cyclic rotation of $f_{A,a}$.
\end{theorem}

\begin{proof}
Given (iii), a tuple with a slot outside the support of \eqref{rec:eq:landau} has zero limit for both;
otherwise $t_j=k_j\log p_j$ and unique factorization gives $\sum_{p_j=p}k_j=0$ at each prime, so the
rotation factors cancel. Conversely, $(E\log p,-E\log p)$ gives $m_A$; at the chosen prime the tests
give every product $\prod_jF(b_j)$ with $\sum b_j\equiv0\pmod E$, $d\le E$. For real $y$ on $\Z/E\Z$ the
power sums $\sum_sz_s^d$ of $z_s=\sum_uy_uf(u-s)$ are combinations of these products, so Newton's
identities give the multiset $\set{z_s}$; with $y_u=(m_A+1)^u$ each $z_s$ encodes one rotation in base
$m_A+1$.
\end{proof}

\begin{corollary}[Coprime packet size; \stG]\label{rec:cor:coprime}
If $\gcd(m_A,E)=1$, all balanced data agree if and only if $B$ is obtained from $A$ by multiplying
every member by one character of $H$.
\end{corollary}

\begin{proof}
With $D_A=\prod_hh^{A(h)}$ the local rotations $v_a\in\mu_E$ satisfy $v_a^{m_A}=(D_B/D_A)(a)$; if
$m_Au\equiv1\pmod E$, then $v_a=(D_B/D_A)^u(a)$ for all $a$, and orthogonality on $G$ gives the
multiplicities.
\end{proof}

For quadratic twists ($E=2$) let $c_A(h)=\sum_uA(u)A(hu)$ count pairwise quotients. Then
$\widehat c_A=\abs{\widehat A}^2$ on $\widehat H$, every element of $\widehat H$ is evaluation at some $a\in G$,
and a two-bin histogram of known mass is determined up to a swap by the difference of its bins;
hence all balanced limits agree if and only if all pair limits agree, if and only if $c_A=c_B$
(homometry). For packets satisfying (iii) and real $w\in C_c^\infty(\R\setminus\set0)$ with $\int w=0$,
the pair statistic $\sum_{\gamma,\gamma'\in Z_A(T)}w(\gamma-\gamma')$ of the two packets differs by
$O(T/(\log T)^{1/3})$, by the cross-character formula of
\cite[Theorem~2]{FordZaharescu2015Unnormalized}, whose main term depends on
$\Lambda(n)^2\abs{\sum_{\chi\in A}\chi(n)}^2$, unchanged by local rotation.

\begin{theorem}[Twenty-eight indistinguishable twist classes; \stG]\label{rec:thm:six}
Let $Q=5\cdot13\cdot17\cdot29=32045$, $\psi_i$ a character of order four modulo $q_i$,
$\psi=\prod\psi_i$ and $h_x=\prod_i\psi_i^{2x_i}$ for $x\in\F_2^4$, so every $\psi h_x$ is primitive of
conductor $Q$. The six-factor packets supported where $x_1x_2+x_3x_4=1$ and where
$x_1x_2+x_3x_4+x_1x_3=1$ have the same balanced limits at every order. Exactly $448$ six-element
subsets of $\F_2^4$ share their data, forming $28$ classes under a common twist, and
\[
 \sum_{0<\gamma\le T}\abs{m_A(\gamma)-m_B(\gamma)}\ge\frac{\sqrt2\log2}\pi T-O(\log(2T)),
\]
the same bound holding with $H$ in place of $T$ on every window $(T,T+H]$, up to $O(\log(2T+2H))$.
\end{theorem}

\begin{proof}
The Walsh transforms of $(-1)^{q_A}$ and $(-1)^{q_B}$ are $4(-1)^{v_1v_2+v_3v_4}$ and the same times
$(-1)^{v_2v_4}$, so both sets have size six and all nontrivial indicator coefficients of modulus two,
and homometry applies; the bilinear parts differ, and twists act by translation. Difference counts
$6$ at zero and $2$ elsewhere mean that the $15$ pair sums are distinct and nonzero; such a six-set
sums to zero with every five points affinely independent, and conversely, so there are
$16\cdot15\cdot14\cdot12\cdot8/6!=448$ of them, in translation orbits of size $16$. At $p\nmid Q$ the
coefficients differ exactly when $(p/q_2)=(p/q_4)=-1$, for one quarter of the classes and for $p=2$,
with difference four; apply \eqref{rec:eq:landau} at $\log2$.
\end{proof}

\begin{theorem}[Uniformity in the frequencies; \stG]\label{rec:thm:homuniform}
If the local histograms of $A$ and $B$ are rotations of each other at every class, then for fixed
$d\ge2$, $a>0$, $0<\delta<2$ and $\eps>0$,
\[
 \sup_{\sum t_j=0,\ a\le\abs{t_j}\le(2-\delta)\log T}\abs{\mathcal B_A(\mathbf t;T)-\mathcal B_B(\mathbf t;T)}
 \ll T^{-\delta/2+\eps}.
\]
\end{theorem}

\begin{proof}
With $X=T^{2-\delta}$, $x=e^{\abs t}$ and $n$ a nearest integer to $x$, the uniform Landau formula
approximates $S_A(t;T)$ by $0$ unless $n=p^k$ is an unramified prime power, and then by
$-\frac{\log p}{2\pi\sqrt x}\psi(p)^{sk}F_A(p,sk)K_T(t-s\log n)$ with $K_T(u)=T^{-1}\int_0^Te^{ivu}\dd v$,
$s=\operatorname{sgn}t$, with error $O(X^{1/2+\eta}\log(2XT)/T)$. If the products $U,V$ of the $n_j$
over positive and negative $t_j$ agree, the leading products agree; otherwise some detuning is at
least $\abs{\log(U/V)}/d\ge1/(d(1+\min(U,V)))$, and $\abs{K_T(u)}\le\min(1,2/(T\abs u))$ bounds both
products by $O((\log2X)^d/T)$.
\end{proof}

\begin{theorem}[One cubic measurement at odd prime exponent; \stG]\label{rec:thm:cubic}
Let $E=\ell$ be an odd prime and choose one prime $p\equiv a\pmod Q$ for each $a\in G$. The limit at
$(\log p,\log p,-2\log p)$ determines $f_{A,a}$ up to rotation; if $\ell\nmid m_A$, these measurements
determine the packet up to a common twist.
\end{theorem}

\begin{proof}
With $\alpha_f=\sum_uf(u)\zeta^u$ and $\sigma(\zeta)=\zeta^2$, the normalized cubic is
$C(f)=\alpha_f^2\overline{\sigma(\alpha_f)}$. If $C(f)=0$ all entries are equal. Otherwise, for
$C(g)=C(f)$ and $b=\alpha_g/\alpha_f$, $\sigma(b)=\bar b^{-2}$, so the radii along the
$\sigma$-orbit satisfy $r_{j+1}=r_j^{-2}$ and are all $1$; then $\sigma(b)=b^2$, $b$ is a root of
unity of odd order in $\Q(\zeta)$, hence $b=\zeta^s$, and $g$ minus the rotated $f$ is constant, hence
zero. The last claim is \cref{rec:cor:coprime}.
\end{proof}

A nonzero difference of two cubics is an algebraic integer with conjugates at most $2m^3$, so it
has modulus at least $(2m^3)^{-(\ell-3)/2}$. The quadratic data cannot replace the cubic: one copy
at $0$ and two at $1$, and its reflection, have equal pair data but are not rotations of each other.
More collisions come from the Maiorana--McFarland bent functions \cite{Dillon1974Elementary}: with
$2r$ primes $q_i\equiv1\pmod4$ and quadratic twists identified with $\F_2^r\times\F_2^r$, the supports
of $x\cdot y+g(y)$, $g(0)=0$, all have the same balanced data, and they fall into exactly
$2^{2^r-1-r}$ common-twist classes; for $r\ge3$ the supports for $g=0$ and $g=y_1y_2y_3$ are not even
affinely equivalent. The finite enumerations are
\artifact{dirichlet_zero_packet_homometry/packet_checks.txt} and
\artifact{dirichlet_zero_packet_homometry/cubic_checks.txt}.

\paragraph{The least number of local measurements.}
At one unramified prime the balanced limits give products $\prod_iF(k_i)$ with $\sum_ik_i\equiv0\pmod E$ of the
Fourier coefficients $F(k)=\sum_xf(x)\zeta_E^{kx}$ of the local histogram $f:C_E=\Z/E\Z\to\Z_{\ge0}$, of known mass
$m=\sum_xf(x)$; rotation multiplies $F(k)$ by $\zeta_E^{kt}$, and the question is how many individual products,
of what orders $s$, determine $f$ up to rotation. Let $\mathcal D=\set{d\mid E:d>1}$, and $\mathcal I$ the set of unordered
pairs $\set{d,e}\subset\mathcal D$ with $\gcd(d,e)>1$.

\begin{theorem}[Exact measurement counts; \stU]\label{rec:thm:probecount}
The least size of a fixed separating list of products, and of a deterministic adaptive procedure,
in the worst case, are
\[
 N_{\mathrm{fixed}}=\abs{\mathcal D}+\abs{\mathcal I}=\tfrac12\bigl(\tau(E)^2+\tau(E)-1-\tau(E^2)\bigr),\qquad N_{\mathrm{adaptive}}=\abs{\mathcal D}=\tau(E)-1 .
\]
A fixed list of minimum size exists with every product of order at most six, and the lower bounds hold
with unrestricted orders, even at the single mass $m=2E^2$. When the mass is not supplied, both minima
increase by exactly one.
\end{theorem}

For $E=p^r$ the counts are $r(r+1)/2$ and $r$; for $E=12$ they are $13$ and $5$; for $E=30$ they are $22$
and $7$.

\begin{lemma}[Products on one Galois orbit; \stU]\label{rec:lem:probefactor}
For $d\in\mathcal D$ put $k_d=E/d$, $\alpha_d=F(k_d)\in K_d=\Q(\zeta_d)$, so $F(uk_d)=\sigma_u(\alpha_d)$ for
$u\in(\Z/d\Z)^\times$. If $a>b\ge0$ and $P(\alpha)=\alpha^a\prod_{j\le b}\sigma_{u_j}(\alpha)$, then $P(\alpha)=0$ only for $\alpha=0$, and
$P(\alpha)=P(\beta)\ne0$ forces $\beta/\alpha$ to be a root of unity of $K_d$, lying in $\mu_d$ if $d$ and $a+b$ are odd.
\end{lemma}

\begin{proof}
Applying all $\sigma_v$ to $P(z)=1$, $z=\beta/\alpha$, gives multiplicative relations with exponent matrix
$aI+\sum_jP_j$, $P_j$ permutation matrices, which is nonsingular by the sup norm; the adjugate gives
$z^{\det}=1$. The roots of unity of $K_d$ are $\mu_d$ for even $d$ and $\mu_{2d}$ for odd $d$, and for odd $d$ an
extra root $-\zeta_d^t$ has image in $-\mu_d$ when $a+b$ is odd.
\end{proof}

\begin{proof}[Proof of \cref{rec:thm:probecount}]
Every residue modulo odd $d$, and every even residue modulo even $d$, is a sum of two units, by
the Chinese remainder theorem. \emph{Pure measurements.} For $d=2$ use the units $(1,-1)$, for odd $d$
$(1,1,-2)$, for even $d>2$ with $3\nmid d$ $(1,1,1,-3)$, and for $6\mid d$ $(1,1,1,1,u,v)$ with $u+v=-4$. By
\cref{rec:lem:probefactor} equal values give equal activity of each order and a unique relative phase
$t_d\bmod d$ with $\beta_d=\zeta_d^{t_d}\alpha_d$. \emph{Bridges.} For $\set{d,e}\in\mathcal I$, $g=\gcd(d,e)$, take units of order $d$
summing to $d/g$ and of order $e$ summing to $-e/g$, with $\ell(n,g)$ units on the side of order $n$:
$\ell=1$ if $n=g$, $3$ if $n>g$ is even with $n/g$ odd, and $2$ otherwise. The product enforces
$t_d\equiv t_e\pmod g$ when both orders are active, and the generalized Chinese remainder theorem gives
a common rotation; Fourier inversion recovers $f$. \emph{Lower bounds.} With the Ramanujan sums
$C_d(x)=\sum_{u\in(\Z/d\Z)^\times}\zeta_d^{ux}$ and $M=2E$: distinguishing $M$ from $M+C_d$ needs a product all of
whose indices have order $d$, also along any adaptive path; and distinguishing $M+C_d(x)+C_e(x)$ from
$M+C_d(x)+C_e(x-1)$, for $\gcd(d,e)>1$, needs a product whose set of orders is exactly $\set{d,e}$. All
these histograms have mass $2E^2$. \emph{Adaptive upper bound.} Process the divisors in a fixed order;
with $L=\lcm$ of the processed active orders and $g=\gcd(d,L)$, $h=d/g$, ask for $F(k_d)^hB$, with $B$ a
nonzero product of known coefficients of index sum $-E/g$. This determines $\alpha_d$ up to $\mu_h$, exactly
the ambiguity left by rotations by multiples of $L$. \emph{Unknown mass.} The product $F(E)F(-E)=m^2$ supplies
the mass, and the constant histograms $M$ and $M+1$ agree in every product containing a nonzero index
modulo $E$, so every list, fixed or adaptive along the transcript of $M$, needs a product all of whose
indices are $\equiv0\pmod E$ in addition to the products above.
\end{proof}

\begin{theorem}[Least maximum order; \stU]\label{rec:thm:probedegree}
For $E>1$ the least possible maximum order over all separating lists is
$\max(\set{c_d:d\in\mathcal D}\cup\set{\ell(d,g)+\ell(e,g):\set{d,e}\in\mathcal I})$, with $c_d=2,3,4,6$ the lengths of the pure
rows; the fixed list above attains both minima at once. The optimal adaptive count is attained
with orders at most $6+3\omega(E)$, with order three for odd prime powers and four for powers of two.
\end{theorem}

\begin{proof}
Inputs with arbitrary nonzero entries on one order are realized by the histograms
$f_\alpha(x)=M+\Tr_{K_d/\Q}(A\alpha\zeta_d^{-x})$, for large integers $A$ and $M$. For $d>2$ compare $\alpha=1$ with $b\in K_d$ of
modulus one at every embedding and not a root of unity (for example $(n+\zeta_d)/(n+\zeta_d^{-1})$): all
balanced quadratic products agree, and for even $d$ balanced pure products have even length. For
$6\mid d$ take $b\in K_6$, on which a unit acts as $b^{\pm1}$ according to its class modulo $6$, so that a
balanced row of length below six cannot see $b$. For the pairs shift the $e$-part by $1$ or $g/2$; the
partial index sum of the $d$-part must be a nonzero multiple of $E/g$, which needs $\ell(n,g)$ units on
each side. For the adaptive orders replace $F(k_d)^h$ by $\alpha_d^a\sigma_u(\alpha_d)\sigma_v(\alpha_d)$ with $a\in\set{3,4}$ and
$h-a=u+v$, and decompose $-E/g$ into primary components, each a sum of at most three indices of an
active order.
\end{proof}

\begin{proposition}[When the mass removes the ambiguity; \stU]\label{rec:prop:probemass}
If the convolution matrix $A_f$ has $\gcd(\det A_f,E)=1$, the pure measurements determine $f$ up to rotation
among integral histograms of the same mass; this holds for $E=p^r$ whenever $p\nmid m$. For $E=p^r$, $r\ge2$,
the single pure measurement at the indices of order $E$ separates all rotation classes of mass $m$ exactly
when $m\le p$.
\end{proposition}

\begin{proof}
If $b$ has the same pure measurements, the ratio $G/F$ takes values in $\mu_E$; its inverse transform $v$
satisfies $A_fv=b$, so $(\det A_f)v$ and $Ev$ are integral, hence $v\in\Z^E$, and Parseval with $\sum v=1$
makes $v$ a point mass. For $E=p^r$, $\mathbb F_p[X]/((X-1)^E)$ has $f$ as a unit when $p\nmid m$. For the
last claim the aligned difference is divisible by $\Phi_E(X)=\sum_{i<p}X^{iE/p}$, hence constant on fibres
modulo $E/p$, and a nonzero difference costs at least $p$ units of mass on each side; for $m>p$ the
histograms $(m-p)\mathbf 1_{\set0}+\mathbf 1_U$ and $(m-p)\mathbf 1_{\set0}+\mathbf 1_V$ for distinct fibres $U,V$ collide.
\end{proof}

\begin{corollary}[Logarithmic windows of zeros; \stG]\label{rec:cor:probewindow}
In the setting of \cref{rec:thm:localclass}, fix $Q,p\nmid Q,E$ and the mass $m$. There is $C$ such that for
all large $T$ the fixed list of \cref{rec:thm:probecount}, evaluated on the relative ordinates of the pooled
positive zeros in every window $(T,T+C\log(2QT))$, determines the local histogram at $p$ up to rotation.
\end{corollary}

\begin{proof}
With $S(t)=Y^{-1}\sum_{T<\gamma<T+Y}e^{it\gamma}$, the Landau formula of Banks \cite{Banks2023Comparing} under
\cref{src:thm:dgrh} gives $S(k\log p)=-\frac{\log p}{2\pi p^{\abs k/2}}\psi(p)^kF(k)+O(\log(2QT)/Y)$ for fixed $k\ne0$. Lift each
row to nonzero integers with sum exactly zero; the product of the $S(k_j\log p)$ depends only on
differences of ordinates and the factors $\psi(p)^{k_j}$ cancel. Distinct values of a product of length
$s\le6$ differ by at least $(2m^s)^{-(\varphi(E)-2)/2}$ for $E>2$, as a nonzero algebraic integer of $K_E$ with
conjugates at most $2m^s$, and by at least one for $E\le2$; choose $C$ to make the errors smaller than
half this and round.
\end{proof}

All lower-bound examples occur arithmetically: for a prime $Q\equiv1\pmod{2E}$, a generator $\psi$ of the
characters modulo $Q$, $H$ of order $E$ and $p$ a primitive root modulo $Q$, every histogram is realized by
primitive twists $\psi h$. For $E=30$, $Q=61$, $p=2$, the histograms $60+C_6(x)+C_{10}(x)$ and
$60+C_6(x)+C_{10}(x-1)$ agree in every product of order at most five, while the sextic with indices
$(5,5,5,3,3,-21)$ gives $\pm729{,}000{,}000$:
\begin{center}
\artifact{dirichlet_packet_reconstruction/zero-packets/check_probes.txt}\\
\artifact{dirichlet_packet_reconstruction/zero-packets/rigidity_output.txt}
\end{center}
